% 本文件由 LaTeX 排版 Agent 根据有效 translation.json 自动生成。
% author=Athanasius; work_id=2806; title=书信

\chapter{书信}

% structure: p0001 source=h1 target=omitted reason=page-title-already-rendered-by-parent

\Emph{书信 1}\newline{}\Emph{书信 2}\newline{}\Emph{书信 3}\newline{}\Emph{书信 4}\newline{}\Emph{书信 5}\newline{}\Emph{书信 10}\newline{}\Emph{书信 11}\newline{}\Emph{书信 13}\newline{}\Emph{书信 14}\newline{}\Emph{书信 17}\newline{}\Emph{书信 18}\newline{}\Emph{书信 19}\newline{}\Emph{书信 20}\newline{}\Emph{书信 22}\newline{}\Emph{书信 24}\newline{}\Emph{书信 27}\newline{}\Emph{书信 28}\newline{}\Emph{书信 29}\newline{}\Emph{书信 39}\newline{}\Emph{书信 40}\newline{}\Emph{书信 42}\newline{}\Emph{书信 43}\newline{}\Emph{书信 44}\newline{}\Emph{书信 45}\newline{}\Emph{书信 46}\newline{}\Emph{书信 47}\newline{}\Emph{书信 48}\newline{}\Emph{书信 49}\newline{}\Emph{书信 50}\newline{}\Emph{书信 51}\newline{}\Emph{书信 52}\newline{}\Emph{书信 53}\newline{}\Emph{书信 54}\newline{}\Emph{书信 55}\newline{}\Emph{书信 56}\newline{}\Emph{书信 57}\newline{}\Emph{书信 58}\newline{}\Emph{书信 59}\newline{}\Emph{书信 60}\newline{}\Emph{书信 61}\newline{}\Emph{书信 62}\newline{}\Emph{书信 63}\newline{}\Emph{书信 64}\par

\SourceRecord{Translated by Archibald Robertson.；Nicene and Post-Nicene Fathers, Second Series；Vol. 4.；Edited by Philip Schaff and Henry Wace.；Buffalo, NY: Christian Literature Publishing Co.,；1892.；<http://www.newadvent.org/fathers/2806.htm>.}

% structure: p0001 source=h1 target=section reason=page-title

\section{书信一}

\Emph{为329年。复活节，法尔穆提月十一日，四月伊德前第八日；戴克里先纪元45年；\Person{君士坦丁}奥古斯都第八次、\Person{君士坦丁}凯撒第四次任执政官；总督\Person{塞普蒂米乌斯·泽尼乌斯}；第二指示周期。}\par

\Emph{论禁食、号角与庆节}\par

来啊，我心爱的，节期召唤我们去守这庆节。再一次，\BibleQuote{正义的太阳}\BibleRef{拉 4:2}，使祂的神圣光芒升起照耀我们，预先宣告了这庆节的时间；我们应顺从祂，在这时间内庆祝，免得时间逝去，欢乐也离我们而去。分辨时间是操练德行最急迫的职责之一；因此，圣\Person{保禄}在教导他的门徒\Person{弟茂德}时，嘱咐他观察时机，说：\Quote{务要宣讲真道，不论顺境逆境，总要坚持不变}——好让他既知道何时顺、何时逆，便能做合乎时宜的事，避免不合时宜的过失。正如智慧之王\Person{撒罗满}所做的那样，万有的天主将一切按时节分配，为使人得救恩的佳音，在适当的时候传遍各处。于是，\BibleQuote{天主的智慧}\BibleRef{格前 1:24}，我们的主及救主耶稣基督，并非不合时宜，而是适时地\BibleQuote{寄居在圣善的灵魂中，造就天主的朋友和先知}\BibleRef{智 7:27}；虽然许多人曾祈求祂来临，说：\Quote{但愿天主的救恩由\Place{熙雍}发出！}——而且，如\WorkTitle{雅歌}所载，那新娘也在祈求和说：\BibleQuote{巴不得你像我的兄弟，像吮吸过我母亲乳汁的兄弟！}\BibleRef{歌 8:1}——巴不得你像人子一样，为我们的缘故承担了人的情感！——然而，万有的天主，时间和季节的创造者，祂比我们更了解我们的事务，如同一位良医，在适当的时机劝人顺从——唯有在这时机我们才能获得医治——祂也在适当的时候，而非不适时地差遣了祂，说：\BibleQuote{在悦纳的时候，我应允了你；在救恩的日子，我帮助了你。}\BibleRef{依 49:8}\par

2. 为此，圣\Person{保禄}敦促我们留意这时机，写道：\BibleQuote{看，如今正是悦纳的时候；看，如今正是救恩的时日。}\BibleRef{格后 6:2} 天主也按时节，藉着\Person{梅瑟}召唤以色列子民参加肋未人的庆节，说：\BibleQuote{你每年三次应为我举行庆节}\BibleRef{出 23:14}（我心爱的，其中之一就是眼下这个庆节），并藉着司祭们的号角声，催促他们遵守；正如圣咏作者所命令的：\Quote{要在月朔吹响号角，在月圆之日，庆祝我们的节期。} 这节经文既吩咐我们在月朔吹号，也在月圆之日吹号，祂使得月圆满之日成为一个神圣的日子；月圆之日当时是一个预像，正如号角也是预像。号角声有时，如前面所说的，是为召人赴庆节；有时是为禁食，有时是为战争。这并非不加庄严，也不是偶然，号角声的设立，是为让所有人都来听从宣告。这教训不仅应从我这得知，更应来自圣经。当天主显现给\Person{梅瑟}时，正如\BibleBook{户籍}{纪}所载，祂说：\BibleQuote{你制造两个银号筒，用来召集会众}\BibleRef{户 10:1}；这对那爱祂的人非常合适。我们由此可知道，这一切都是为\Person{梅瑟}的时代所预备的——是的，在预像存留的时候，都应遵守，这一切的安排都是为那\BibleQuote{直到改革的时期}\BibleRef{希 9:10}。因为（祂说）\BibleQuote{若你们在本土要出去作战，攻击那压迫你们的仇敌，（因为这类事只涉及本土地，而非更远，）你们就吹响号角，在上主面前必蒙记念，也必从敌人手中被救出来。}\BibleRef{户 10:9} 他们不仅在战争中吹号，法律也规定了庆节的号角。再听他说下去：\Quote{在你们的庆日、节日、月朔，你们都要吹响号角。} 谁若听到法律命令关乎号角，切勿以为这是轻微可轻视的事；这是一件奇妙又可畏的事。因为号角的声音超越任何其他声音或乐器，它令人警醒，令人畏惧；以色列人因此接受教导，因为他们当时还是个孩子。但为使这宣告不被视为纯属人为，而是超越人力，它的声音就如同他们面对山颤栗时所发出的声音那样\BibleRef{出 19:16}；因此他们想起了当时所颁的法律，并遵守了它。\par

3. 因为法律曾是奇妙的，影子曾是卓越的；否则，它不会在听的人心中产生敬畏，引发崇敬；尤其是在那些当时不仅听见，而且看见这些事的人心中。这些原是预像，是作为影子而行的。但我们应转向其意义，从此远离影象，接近真理，并仰望我们救主的司祭号角，它们发出呼声，召唤我们，有时是要我们作战，正如可赞颂的\Person{保禄}所说：\BibleQuote{我们不是与血肉搏斗，而是与率领者、掌权者、这黑暗世界的霸主、以及天上的邪恶神灵搏斗}\BibleRef{弗 6:12}。另一次，召唤乃是为着童贞、克己以及婚姻和谐，说道：对童贞者，讲论有关童贞的事；对喜爱节德之路的人，讲论有关节德的事；对已婚的人，讲论有关崇高婚姻的事；如此将各自的德行和相应的光荣报偿分配给各人。有时，召唤乃是为着守斋，有时则是为着庆节。再听同一位\EditorialNote{宗徒}吹响号角并宣告道：\BibleQuote{我们的逾越节羔羊基督已被祭献；因此我们过节，不可用旧酵母，也不可用邪恶与恶毒的酵母}。如果你愿听一只远较这一切更响亮的号角，就聆听我们的救主说：\BibleQuote{在庆节最盛大也是最后一天，耶稣站着大声喊道：谁若渴，到我这里来喝吧}\BibleRef{若 7:37}。因为救主不仅要召唤我们赴一节，而是要赴\Quote{大庆节}，只要我们预备好去听，并遵行每一号角的宣告。\par

4. 如上文所述，既有种种不同宣告，你就当如在预像中一般，听先知吹响号角；继而转向真理之后，要预备听这号角的宣告，因为他说道：\BibleQuote{在熙雍吹响号角，祝圣一个斋期}\BibleRef{岳 2:15}。这是一只警告的号角，并以极大的恳切命令，当我们守斋时，应当圣化这斋戒。因为并非所有呼求天主的人，都圣化天主，因为有人玷污祂；然而不是玷污祂本身——那是不可能的——而是玷污他们内心对祂的观念；因为祂是圣的，并喜悦于圣人。因此，可赞颂的\Person{保禄}指责那些侮辱天主的人：\BibleQuote{违犯法律的人侮辱天主}\BibleRef{罗 2:23}。那么，为与那些玷污斋戒者分别开来，他在这里说：\Quote{祝圣斋期}。因为许多人涌向斋戒，却以心中的思念玷污了自己，有时是对兄弟行恶，有时是胆敢欺诈。此外，更不必提，许多人自高举于邻人之上，因而造成极大的损害。因为守斋的矜夸未曾给\Person{法利塞人}带来任何益处，虽然他每周守斋两次\BibleRef{路 18:12}，只因他自高举于税吏之上。同样，\OriginalTerm{圣言}{Word}也曾以同样的方式，借着先知\Person{依撒意亚}劝勉以色列子民，责备他们那类的斋戒，说：\BibleQuote{这不是我所选定的斋戒和日子，并非要人克苦己身；你们即便屈颈如钩，将苦衣和灰尘铺在身下，你们也不能称这是可蒙悦纳的斋戒}\BibleRef{依 58:5}。为要能显示我们守斋时当如何为人，以及斋戒应有的性质，再听天主命令\Person{梅瑟}，正如在\WorkTitle{肋未纪}中所记载的：\BibleQuote{上主训示梅瑟说：七月十日，应是一个赎罪之日，你们应召集圣会，克苦己身，并向上主呈献全燔祭。}而后，为在这事上界定法律，祂继续说道：\BibleQuote{凡不克苦己身的灵魂，应从百姓中铲除。}\par

5. 看哪，我的弟兄们，斋戒能成就何等大事，法律又是以何种方式命令我们守斋。不仅应以肉身守斋，还应以灵魂守斋。当灵魂不随从邪恶的见解，而以相称的德行为食时，它就受到贬抑。因为德行与恶行是灵魂的食粮，它能随意吃这两种食粮中的一种，并按照自己的意愿偏向其中之一。如果它倾向于德行，就会以德行、以正义、以节德、以温良、以刚毅为滋养，正如\Person{保禄}所说：\BibleQuote{以真理之言得到滋养}\BibleRef{弟前 4:6}。我们的主正是如此，祂曾说：\BibleQuote{我的食物就是承行我在天之父的旨意}\BibleRef{若 4:34}。但如果灵魂不是这样，而是向下堕落，那么它所赖以滋养的，便只有罪恶。因为\OriginalTerm{圣神}{Holy Ghost}如此描述罪人及其食物，论及\OriginalTerm{魔鬼}{devil}说：\Quote{我已将他交给\Place{埃塞俄比亚}人民作食物。}因为这就是罪人的食物。正如我们的主和救主\Person{耶稣基督}，作为来自天上的食粮，是圣人们的精神食粮，正如经上说：\BibleQuote{你们若不吃人子的肉，不喝他的血，在你们内便没有生命}\BibleRef{若 6:53}；同样，魔鬼是不洁之人的食物，是那些不作任何光明之事，只行黑暗之事者的食物。因此，为使他们远离并转离恶行，祂命令他们以德行的食粮作为滋养；那就是心灵谦卑、甘受侮辱的卑微、承认天主。因为这样的斋戒不仅为灵魂赢得宽恕，而且若善守圣洁，还能准备圣人，使之超脱尘世之上。\par

6. 我接下来要说的确实奇妙，甚至是那些极为神奇的事；然而离真理不远，你们可以从\OriginalTerm{神圣经典}{sacred writings}中得知。那位伟大的人物\Person{梅瑟}，当他守斋时，曾与天主交谈，并领受了法律。伟大而圣善的\Person{厄里亚}，当他守斋时，被认定为配得\OriginalTerm{神视}{divine visions}，最后如同那位升天者一样被提去。\Person{达尼尔}，当他守斋时，虽然年纪很轻，却被托付了\OriginalTerm{奥秘}{mystery}，唯独他明白了君王的秘密，并被认为配得\OriginalTerm{神视}{divine visions}。但这些人的守斋时间之长令人惊奇，日子延长，所以谁也不要轻易陷入不信；反而要相信并知道，\OriginalTerm{默观天主}{contemplation of God}，以及从祂而来的话语，足以滋养那些聆听的人，并代替一切食物。因为天使也是靠着时常瞻仰在天上的圣父和\OriginalTerm{救主}{Saviour}的面容而得以维持，并没有别的方法。同样，\Person{梅瑟}在与天主交谈期间，身体虽在守斋，却由天主的话得到了滋养。当他下到人群中，天主离开他上升之后，他便像其他人一样感到饥饿了。因为并没有记载他守斋超过四十天——那些就是他与天主交谈的日子。通常，每一位圣人也都曾被认为配得类似这样\OriginalTerm{超性的滋养}{transcendent nourishment}。\par

7. 因此，我亲爱的各位，我们的灵魂既由\OriginalTerm{神圣食粮}{divine food}，即\OriginalTerm{圣言}{Word}所滋养，并按天主的旨意，在身体的外在事物上守斋，就让我们以相称的方式庆祝这伟大而带来救恩的庆节。即使是无知的犹太人，也通过\OriginalTerm{预像}{type}领受了这神圣食粮，就是在逾越节吃羔羊。但由于不明白预像，他们直到今天仍在吃羔羊，错在离开了那城和真理。只要\Place{犹太}和那城存在，就有预像、羔羊和\OriginalTerm{影子}{shadow}，因为法律这样命令说：\Quote{这些事不可在另一座城举行；而应在犹太地区，不得在犹太地区以外的地方。} 此外，法律还命令他们奉献全燔祭和祭献，除了\Place{耶路撒冷}的祭台外，不可有别的祭台。正因为如此，只有在那一座城里有祭台和圣殿，不准在任何其他城举行这些礼仪，这样，当那座城终结时，那些\OriginalTerm{象征的}{figurative}事物也就随之废除了。\par

8. 现在你们观察：自我们的\OriginalTerm{救主}{Saviour}来临以后，那座城已经终结，犹太人的整个地区都被摧毁了；因此，由这些事的见证（我们亲眼证实，无需其他证据），\OriginalTerm{影子}{shadow}必然要终结。这些事不应从我这里学习，而是先知圣声早已预言，呼喊说：\BibleQuote{看，传报喜讯、宣布和平者的脚已站在山上！}\BibleRef{鸿 1:15}；他所宣布的信息无非是继续对他们说：\BibleQuote{犹大啊，遵守你的庆节，向上主偿还你的誓愿！因为他们不再回到旧事；已经完成，已被除去：那曾向脸上吹气，并拯救你脱离苦难者已上升。} 那么，那位上升的是谁？可以这样对犹太人说，以使影子的夸耀也被废除；听这句话不是无益的：\Quote{已经完成，那吹气者已上升。} 因为在那位吹气者上升以前，什么也没有完成。但他一上升，就完成了。那么，犹太人啊，正如我前面所说，他是谁？如果是\Person{梅瑟}，这断言就是假的，因为那时人民还没有来到唯一受命举行这些礼仪的地方。如果是\Person{撒慕尔}，或任何其他先知，那也是曲解真理；因为直到那时，这些事仍在犹太举行，城也还在。因为城在，这些事就必须举行。所以，我亲爱的，上升的绝非这些人物。但若你们愿意听真实的事，并远离犹太人的无稽之谈，请看我们的\OriginalTerm{救主}{Saviour}上升了，并且\BibleQuote{向脸上吹气，并对门徒们说：\InnerQuote{领受圣神吧！}}\BibleRef{若 20:22}。因为这事一完成，一切都终结了，因为祭台破碎了，\OriginalTerm{圣殿的帐幔}{the veil of the temple}裂开了；虽然城还未被毁灭，但那\OriginalTerm{可憎之物}{abomination}已准备坐在圣殿中心，城和那些古老的规条也要领受它们的最终\OriginalTerm{终结}{consummation}。\par

9. 既然我们已经越过了那\OriginalTerm{影子}{shadow}的时期，不再在其下举行礼仪，而是转向了主，正如经上说：\BibleQuote{主就是那神，主的神在哪里，那里就有自由。}\BibleRef{格后 3:17}；我们听从圣号，不再宰杀物质的羔羊，而是那被宰杀的真正羔羊，我们的主耶稣基督，\BibleQuote{祂如同羊被牵去宰杀，又像羔羊在剪毛者前缄默。}\BibleRef{依 53:7}；我们被祂的\OriginalTerm{宝血}{precious blood}所洁净，这宝血说话比\Person{亚伯尔}的血更美，我们的脚穿上了福音的准备，手持主的\OriginalTerm{棍杖}{rod and staff}，那圣人曾因此得到安慰，说：\BibleQuote{你的棍杖，它们安慰了我；} 总之，在各方面都准备好了，一无所挂虑，因为正如圣\Person{保禄}所说的：\BibleQuote{主临近了。}\BibleRef{斐 4:5}；又如我们的救主所说：\Quote{在你们想不到的时辰，主就来了——所以我们过节，不可用\OriginalTerm{旧酵母}{old leaven}，也不可用邪恶和凶恶的酵母，而要用真诚和真理的\OriginalTerm{无酵饼}{unleavened bread}。脱去\OriginalTerm{旧人}{old man}和他的行为，穿上\OriginalTerm{新人}{new man}，就是照天主在真理中所造的。} 以心谦和纯洁的良心，昼夜默思法律；并除去一切伪善和欺诈，远离骄傲和欺骗，让我们承担起对天主和对近人的爱德，这样，我们成为新的受造物，领受新酒，就是圣神，我们才能适当地庆祝这个庆节，即这新果实的月份。\par

10. 我们于帕尔穆特月第五日（3月31日）开始圣斋期，并再增加六天，以对应那六个神圣而伟大的日子——这六日是此世创造的象征；然后于同月帕尔穆特月第十日（4月5日），即当周的圣安息日，停止斋戒休息。当圣周的第一日黎明升起、光照我们时，即在同月第十一日（4月6日），从这一天起我们再一周一周地计算整整七周，以便在神圣的五旬节守节——这一日对犹太人曾是预像的七七节，他们在那一天赦免并清还债务；实际上，这一天在各方面都是一个解救之日。让我们在\OriginalTerm{大周}{Great Week}的第一日守节，作为来世的象征，在此我们预先领受保证，日后我们将获得永生。然后，当我们离开此世，我们将与基督共守完美的庆节，那时我们要像圣人们一样欢呼说：\BibleQuote{我要进入奇妙帐幕的居所，进入天主的殿宇；在欢乐和称谢声中，在欢腾的歌唱声中；}在那里痛苦、悲伤和叹息都已远遁，欢欣与喜乐将临到我们头上！愿我们被判定堪当参与这些事。\par

11. 我们要记念穷人，不忘善待异乡人；首要的是，我们应以全部灵魂、全部力量、全部能力爱天主，并爱近人如自己\BibleRef{玛 22:37-39}。如此，我们就能领受那眼所未见、耳所未闻、人心所未想到的，就是天主为爱祂的人所预备的这一切\BibleRef{格前 2:9}，藉着祂的唯一圣子，我们的主和救主耶稣基督；愿光荣和权能，藉着祂，并在圣神内，归于唯一的圣父，至于无穷世之世。阿们。\par

你们要以圣吻彼此问候。与我同在的众弟兄都问候你们。\par

\SourceRecord{Translated by R. Payne-Smith.；Nicene and Post-Nicene Fathers, Second Series；Vol. 4.；Edited by Philip Schaff and Henry Wace.；Buffalo, NY: Christian Literature Publishing Co.,；1892.；<http://www.newadvent.org/fathers/2806001.htm>.}

% structure: p0001 source=h1 target=section reason=page-title

\section{第二封信}

\Emph{为330年。复活节在\OriginalTerm{法尔穆提月}{Pharmuthi}24日；\OriginalTerm{5月朔日前13日}{xiii Kal. Mai.}；\OriginalTerm{戴克里先纪元}{Æra Dioclet.}46年；\OriginalTerm{执政官}{Coss.} \Person{加里奇阿努斯}与\Person{瓦莱里乌斯·叙马库斯}；\OriginalTerm{总督}{Præfect} \Person{马格尼尼亚努斯}；\OriginalTerm{第三小指示周期}{Indict. iii}。}\par

我的弟兄们，复活节与喜乐再次来临；主再次带领我们进入这个时节；使我们能依照惯例，以祂的言语滋养后，合宜地守这节期。那么，让我们与那些曾经宣告同样节期、并为我们树立在基督内生活榜样的圣人们，一同庆祝这属天的喜乐吧。因为他们不仅受委讬宣讲福音，而且，若我们查考，我们就会看见，正如经上所载，福音的能力在他们身上彰显出来。\BibleQuote{所以你们该效法我}\BibleRef{格前 4:16}，他写信给\Place{格林多}人说。如今，这宗徒的规劝也训示我们众人，因为他所颁给个别人的命令，同时也在各处命给每一个人，因为他原是\Quote{在信仰和真理上，做了外邦人的教师}。而且，大致而言，所有圣人的命令也以同样的方式催促我们，正如\Person{撒罗满}引用\BibleBook{}{箴言}说：\BibleQuote{我儿，要听你父亲的训诲，专心领受明白，因为我给你们一个好的礼物，不要离弃我的言语；因为我在我父亲面前是顺命的儿子，在我母亲眼中是娇儿}\BibleRef{箴 4:1}。因为一位正义的父亲会良好地教养子女，当他勤于按照自己正直的品行教导他人时，这样，即使遭遇反对，他也能在听到那话时不致羞愧：\BibleQuote{你教导别人，却不教导你自己吗？}\BibleRef{罗 2:21}但更要像那善仆，既能救自己，也能得着别人；这样，当他所受讬的恩宠加倍后，他或许能听到：\BibleQuote{好，善良忠信的仆人，你在少许事上忠信，我要派你管理许多事：进入你主人的福乐吧！}\BibleRef{玛 25:21}。\par

2. 那么，让我们如合宜的，就如任何时候，尤其在这节期的日子，不但要听，而且还要实行我们救主的诫命；这样，效法了圣人们的行为之后，我们方可共同进入那在天上而非短暂、却真实常存的主的福乐；作恶的人既自绝于此福乐，留给他们的就只有他们行径的果实：忧愁、痛苦及带折磨的呻吟。让人看看这些人变成了何等模样，他们既不带有圣人们的品行，也不带有那正确悟性——人起初即以此成为理性的并带有天主的肖像。但他们竟蒙羞地被比作无知的畜牲，在非法的享乐中与它们相似，就被说成是\BibleQuote{纵欲的马}\BibleRef{耶 5:8}；同样，因为他们的狡诈、错谬和致死的罪，他们被称作\Quote{毒蛇的种类}，正如\Person{若翰}所说。既然已经如此堕落，像蛇一般爬伏于尘土中，他们的心思只放在可见的事物上，便以为这些是好的，乐在其中，服事自己的私欲而不服事天主。\par

3. 然而即使处在这种状态，那正是为此而来，要寻找并拯救迷失者的热爱世人的\OriginalTerm{圣言}{Word}，也力图制止他们如此愚昧，祂呼喊说：\BibleQuote{不要像那无知的骡马，须用嚼环辔头勒住它的嘴}。因为他们粗心大意、仿效恶人，先知在心神内祈祷说：\BibleQuote{你于我，有如\Place{腓尼基}的商人一样}。而施行惩罚的圣神也用这些话申斥他们：\BibleQuote{主啊，在祢的城里，祢必轻视他们的形象}。就这样，他们变成了愚人的模样，悟性如此低落，以致由于过度的思量，竟将天主的智慧比作他们自己，以为与他们自己的伎俩相似。因此，\Quote{他们自负为智者，反而成了愚妄，将不朽的天主的光荣，变为可朽的人、飞禽、走兽和爬虫形状的偶像。因此，天主就任凭他们陷入堕落的心神，去行那些不当的事。}因为他们不听从那谴责他们的先知之声说：\BibleQuote{你们将主比拟谁，以什么像与他相比呢？}\BibleRef{依 40:18}也不听\Person{达味}为这样的人祈祷而唱的：\BibleQuote{制造偶像的人，都与偶像相似；凡信赖偶像的人，也与偶像一样}。他们对真理视而不见，竟把石头当作天主，因而像无知的受造物一样，行在黑暗中，正如先知呼喊的：\BibleQuote{他们听是听，但不了解；看是看，却不明白；因为这百姓的心迟钝了，耳朵沉重了}\BibleRef{依 6:9}。\par

4. 那些不守\OriginalTerm{庆节}{feast}的人，甚至至今仍保持这些恶习，佯作并捏造庆节之名，却引入哀悼之日而非欢乐之日；\BibleQuote{主说：\InnerQuote{恶人绝无平安！}\BibleRef{依 48:22}}正如智慧书所说：\Quote{欢笑与喜乐从他们口中被夺去。}这便是恶人的庆节。然而，主明智的仆人，真正穿上了那依照\OriginalTerm{天主}{God}所造的新人\BibleRef{弗 4:24}，领受了\OriginalTerm{福音}{Gospel}的话，并将给弟茂德的训令视为一条普遍诫命，这训令说：\BibleQuote{你要在言语、品行、爱德、信德和洁德上，作信徒的典范。\BibleRef{弟前 4:12}}他们如此善守庆节，甚至不信者看见他们的端正品行，也会说：\BibleQuote{天主确实与你们同在。\BibleRef{格前 14:25}}因为凡接待\OriginalTerm{宗徒}{apostle}的，就是接待那派遣他来的\BibleRef{玛 10:40}；同样，凡效法\OriginalTerm{圣人}{saints}的人，便在各方面以主为终向与目标，正如保禄效法主，进而说：\BibleQuote{你们该效法我，如我效法了基督一样。\BibleRef{格前 11:1}}因为首先有\OriginalTerm{救主}{Saviour}自己的话，祂从\OriginalTerm{天主性}{divinity}的高处，与门徒交谈时，说：\BibleQuote{跟我学吧！因为我是良善心谦的，这样你们必要找得你们灵魂的安息。\BibleRef{玛 11:29}}随后，祂把水倒在盆里，取手巾束上腰，给门徒洗脚，并对他们说：\BibleQuote{你们明白我给你们做的吗？你们称我为师、为主，说得正对，我原本就是。那么我这为主、为师的，给你们洗脚，你们也该彼此洗脚；因我给你们立了榜样，叫你们照我给你们做的去做。\BibleRef{若 13:12}}\par

5. 啊！弟兄们，我们当如何赞叹救主的慈爱？人当以何等的能力，何等响亮的号角，来宣扬祂的这些恩惠！我们不仅应带有祂的肖像，更应从祂领受天上生活的典范与榜样；祂既已开始，我们也应继续下去，受苦难时不恐吓，受辱骂时不还口，却应祝福诅咒我们的人，并在一切事上将自己托付于按正义审判的天主。因为那些如此立志，并按照福音塑造自己的人，将成为有分于基督的人，成为效法宗徒品行的人；正因此，他们堪受保禄的称赞，他曾这样称赞格林多人，说：\BibleQuote{我称赞你们，因为你们在一切事上都记念我。\BibleRef{格前 11:2}}后来，因有些人使用他的话，却随从私欲选择听取，并胆敢曲解，就如依默纳约和亚历山大之徒，以及在他们之前的撒杜塞人，他们如他所说，\Quote{在信德上覆舟}，并嘲弄复活的\OriginalTerm{奥迹}{mystery}，所以他紧接着说：\Quote{我怎样传授给你们，你们就怎样持守。}这的确意味着，我们不该有不同于导师所传授的思想。\par

6. 因为那些恶人不仅在外部伪装，主所说的\Quote{披上羊皮}，显像粉白的坟墓；他们还口诵天主的圣言，内心却怀藏恶念。最初如此伪装的是蛇，即自始就发明邪恶的魔鬼，它伪装与厄娃交谈，立即欺骗了她。但在它之后并与它为伍的，是一切发明不法\OriginalTerm{异端}{heresy}的人；他们虽提及圣经，却不持守圣人所传授的意见，反将它们视为人的传授，因而陷入错误，因为他们不明了圣经，也不明了其能力。\BibleRef{玛 22:29}因此，保禄理所当然地称赞格林多人\BibleRef{格前 11:2}，因为他们的意见与他的传授相符。而主极其正当的责斥犹太人，说：\BibleQuote{你们为什么因着你们的传授而违犯天主的诫命呢？\BibleRef{玛 15:3}}因为他们随从自己的见识，改变了从天主领受的诫命，宁愿恪守人的传授。关于这些，不久之后，真福保禄又对处于这危险中的迦拉达人发出训令，写道：\BibleQuote{如果有人给你们宣讲不同于你们所领受的，他该受诅咒。\BibleRef{迦 1:9}}\par

7. 因为圣人的言语与人的虚构幻想之间毫无相通之处；圣人是真理的执事，宣讲天国，而那些走向相反方向的人，只知吃喝，以为他们的结局就是归于虚无，他们说：\BibleQuote{我们吃喝吧！因为明天就要死了。\BibleRef{依 22:13}}因此，真福路加指摘人的虚构，而传授圣人的记述，他在福音的开端说：\BibleQuote{德敖斐罗阁下，已有许多人着手将其间所发生的事，照那些从起初就亲眼看见并为\OriginalTerm{圣言}{Logos}服役的人，所传授给我们的，据实记述。我也在从头仔细观察一切之后，觉得也应该按着次序，写给您，使您知道所接受教导的真确无误。\BibleRef{路 1:1}}因为每位圣人怎样领受，就怎样完整不变地传授，以坚固这奥迹的道理。圣言要我们做这些人的门徒，这些人按理也应是我们的导师，只需听从他们，因为只有他们才宣认\Quote{这话是确实的，值得完全接纳}\BibleRef{弟前 1:15}；他们并非因耳闻而成为门徒，乃是亲眼看见并为圣言服役，将他们从祂所听见的传授下来。\par

如今，有些人记述了我们的\OriginalTerm{救主}{Saviour}所行的奇事，并宣扬了祂永恒的\OriginalTerm{天主性}{Godhead}。另一些人则记载了祂由\OriginalTerm{童贞女}{Virgin}降生成人，并宣告了\OriginalTerm{神圣的逾越节}{holy passover}的庆典，说：\BibleQuote{我们的逾越节羔羊基督，已被祭杀作了牺牲。}\BibleRef{格前 5:7}；好使我们，无论是个人还是团体，以及世间所有教会都能记念，正如经上所载：\BibleQuote{达味的后裔耶稣基督从死者中复活了，这是我的福音}\BibleRef{弟后 2:8}。我们也不可忘记保禄向格林多人所宣告的；我指的是祂的复活，借此\BibleQuote{毁灭那握有死亡权势者——魔鬼}\BibleRef{希 2:14}；并使我们与祂一同复活，解开了死亡的束缚，赐予祝福代替诅咒，喜乐代替忧伤，盛宴代替哀悼。在这\OriginalTerm{复活节}{Easter}的神圣喜乐中，这喜乐常存于我们心中，我们便永远欢喜，正如保禄所命令的：\BibleQuote{不断祈祷，事事感谢}\BibleRef{得前 5:17}。因此，我们并未疏于通知节期，一如我们从\OriginalTerm{教父们}{Fathers}所领受的。我们再次书写，并再度遵守\OriginalTerm{宗徒的传统}{apostolic traditions}，在聚集祈祷时彼此提醒；共同举行庆典，同心合意地真诚感谢主。如此感谢祂，并效法圣人，我们正如\OriginalTerm{圣咏作者}{Psalmist}所说：\BibleQuote{我必要时时赞美上主}。因此，当我们正当地庆祝这节期，我们便配得那在天堂里的喜乐。\par

8. 我们于\OriginalTerm{法慕诺斯月}{Phamenoth}\EditorialNote{谷 9}十三日开始\OriginalTerm{四十日的斋期}{fast of forty days}。在连续禁食之后，让我们于\OriginalTerm{法尔穆提月}{Pharmuthi}十八日（四月十三日）开始\OriginalTerm{圣周}{holy Paschal week}。然后于同月二十三日（四月十八日）安息，随后在二十四日（四月十九日），即一周的第一日，庆祝节期；此后加上\OriginalTerm{五旬节}{Pentecost}那七个星期，在我们主基督耶稣内全然喜乐欢跃，愿光荣和权能借着祂，在\OriginalTerm{圣神}{Holy Ghost}内，归于圣父，至于无穷世之世。阿们。\par

与我同在的弟兄们问候你们。你们要以\OriginalTerm{圣吻}{holy kiss}彼此问候。\par

圣\Person{亚大纳削}，\Place{亚历山大}的主教，第二封《\WorkTitle{节期书信}》终。\par

\SourceRecord{Translated by R. Payne-Smith.；Nicene and Post-Nicene Fathers, Second Series；Vol. 4.；Edited by Philip Schaff and Henry Wace.；Buffalo, NY: Christian Literature Publishing Co.,；1892.；<http://www.newadvent.org/fathers/2806002.htm>.}

% structure: p0001 source=h1 target=section reason=page-title

\section{书信三}

\Emph{主历331年。复活节在法尔穆提月十六日；四月伊德日前三天；戴克里先纪元47年；执政官为安尼乌斯·巴苏斯与阿布拉比乌斯；总督弗洛伦提乌斯；第四征税期。}\par

我亲爱的弟兄们，庆节的日子再次临近我们，这日子超越一切，应专务祈祷；法律命令我们遵守它，若我们缄默度过，则是不虔敬的事。因为虽然我们被那些磨难我们的人所压制，以致因他们的缘故，我们无法向你们宣布这时期；然而感谢\Quote{安慰忧苦者的天主}，我们没有被控告者的邪恶所胜过而缄默；反而顺从真理的声音，在庆节之日与你们一同高声呼喊。因为万有的天主曾命令说：\Quote{你吩咐以色列子民遵守逾越节。}圣神在圣咏中劝勉说：\Quote{在月朔之日，在你们的庆节之日，要吹响号角。}先知也呼喊说：\Quote{犹大啊，你要遵守你的庆节。}我向你们发送消息，并非视你们为无知；而是向知道的人公布，好使你们明白：虽然人使我们分离，天主却使我们成为同伴，我们趋赴同一庆节，并不断朝拜同一的主。我们守这庆节并非像那些谨守时日的人，因为我们知道宗徒曾责备这样的人，他如此说：\BibleQuote{你们谨守某日、某月、某时、某年！}\BibleRef{迦 4:10}相反，我们因这庆节而视此日为庄严；好使我们所有在各处事奉天主的人，能在祈祷中一同蒙天主悦纳。因为真福保禄在宣布类似这样的喜乐临近时，并非宣布日子，而是宣布主，我们是为了祂的缘故而守这庆节，他说：\BibleQuote{我们的逾越节羔羊基督，已被祭杀作了牺牲；}\BibleRef{格前 5:7}好使我们众人，在默观圣言的永恒时，前来事奉祂。\par

二、庆节是什么呢？岂不是灵魂的事奉吗？那事奉又是什么呢？岂不是向天主的长久祈祷和无尽的感恩吗？忘恩负义的人远离这些，理所当然地被剥夺由此而来的喜乐；因为\Quote{欢欣与喜乐已从他们口中消失。}因此，[天主的]话不容他们享有平安；\BibleQuote{恶人决无平安，主说。}\BibleRef{依 48:22}他们痛苦忧伤地劳碌。所以，甚至那欠一万塔冷通的人，在主的面前，福音也不赐予宽恕。\BibleRef{玛 18:24}因为他虽蒙赦免了大事，却在小事上忘恩，以致他也为先前的事受罚。这确实公义，因为他自己既蒙受了仁慈，就应当怜悯与他同为仆人的。那领了一个塔冷通的，把它包在手巾里藏在地下，同样因忘恩负义而被逐出，听到这样的话：\BibleQuote{可恶懒惰的仆人！你既知道我在没有下种的地方收割，在没有散布的地方聚敛；那么，你就该把我的银子交给钱庄里的人，待我回来时，把我的本钱连本带利取回。所以，你们把这个塔冷通从他手中夺过来，给那有十个塔冷通的。}\BibleRef{玛 25:26}当然，当他被要求将属于主人的东西交还时，他本应承认施与者的仁慈和所赐之物的价值。因为施与者并非严苛之人；若是，他起初就不会施与；所赐之物也并非无益空虚，否则他就不会受责备。但施与者是美善的，所赐之物也是能结果实的。因此，如同箴言所说：\BibleQuote{在播种时囤粮不粜的，必受诅咒。}\BibleRef{箴 11:26}同样，忽视恩宠，不加耕耘而将其隐藏的人，理当被当作邪恶忘恩的人被逐出。为此，他称赞那些增加了[塔冷通]的人，说：\BibleQuote{好！善良忠信的仆人，你在小事上忠信，我要委派你管理许多大事；进入你主人的福乐吧！}\BibleRef{玛 25:23}\par

三、这是正确且合理的；因为，如圣经所宣称，他们赚得的与所领受的一样多。我亲爱的诸位，我们的意志应当与天主的恩宠同步，不要落后；免得我们的意志闲散时，所赐的恩宠开始离去，仇敌见我们空虚赤裸，就进入[我们内]，正如福音中所说那个被魔鬼离去之人的情况；\BibleQuote{因为他在干旱之地走过，带了七个比自己更恶的邪魔进入那里；回去后，见房屋空着，便住下来，那人后来的处境比先前更坏了。}因为离弃德行便给不洁的邪魔的进入留了余地。此外，还有宗徒的训令，即赐予我们的恩宠不应变为无益；因为他特别写给门徒的那些话，通过他，也对我们强调说：\BibleQuote{不要疏忽你心内的神恩。因为耕种自己田地的，必得饱食；但懒散人的道路布满荆棘。}因此圣神预先警告人不要陷入其中，说：\BibleQuote{你们要开垦荒地，不要撒种在荆棘中。}因为当一个人轻视所赐的恩宠，立刻陷入世俗的挂虑，便将自己交给私欲；这样，在迫害的时候他就跌倒，变得全无果实。先知指出了这种疏忽的结局，说：\BibleQuote{凡懈慢执行主工作的人，是可咒骂的。}\BibleRef{耶 48:10}因为主的仆人应当勤勉谨慎，更当炽热如火焰，好使他藉着热切的心灵，摧毁一切肉欲之罪后，得以接近天主；按照圣人们的说法，天主被称为\BibleQuote{吞噬的烈火}。\par

4. 因此，万有的天主，\BibleQuote{祂使自己的天使为风[spirits]}，是灵，\BibleQuote{祂的仆役为火焰}。因此，在离开埃及时，因为群众不属于这品性，祂禁止他们挨近那座天主正为他们颁布法律的山；但祂却召了圣梅瑟上前，因为他心灵炽热，拥有不灭的恩宠，说：\Quote{\BibleQuote{让梅瑟独自前来}\BibleRef{出 24:2}}。他也进入了云中，当山冒烟时，他并未受伤；反而，借着\BibleQuote{主的言语，是纯净的言语，如同在泥炉中炼过七次的白银}，他经净化后下来。因此，圣保禄切望我们不要让我们所得的圣神的恩宠变冷，便劝勉说：\Quote{\BibleQuote{不要消灭圣神}\BibleRef{得前 5:19}}。因为这样，如果我们坚守起初赐给我们的圣神到底，我们就仍能分享基督。因为他说的\Quote{不要消灭}，不是因为圣神被置于人的权势下，会遭受什么；而是因为恶劣和不知感恩的人显然是想要消灭圣神，他们像不洁的人一样，用不圣洁的行为迫害圣神。\Quote{\BibleQuote{因为智慧不进入存心不良的灵魂里，也不住在一个屈服于罪恶的身体内；她必远离欺诈，远离无知的思念；邪恶，一被识别，她就退避。}\BibleRef{智 1:5}}如今，这些不明了、欺诈、贪爱罪恶的人，仍在黑暗中行走，没有那\Quote{\BibleQuote{普照每人的真光，正进入这世界}\BibleRef{若 1:9}}。这样的火曾抓住耶肋米亚先知，那时主的话在他内如火一般，他说：\Quote{\BibleQuote{我从各处逃离，再不能忍受。}}我们的主耶稣基督，因祂的良善和爱人之心，来到世上为把这火投在地上，并说：\Quote{\BibleQuote{我是多么切望它已经燃烧起来！}\BibleRef{路 12:49}}就如祂在\BibleBook{厄则克耳}{先知书}中作证，祂渴望人的悔改，胜过他的丧亡；这样，罪恶能在众人心中被彻底烧尽，好使灵魂净化后能结出果实；因为祂所撒的种子会结实：\Quote{\BibleQuote{有的三十倍，有的六十倍，有的一百倍。}\BibleRef{谷 4:20}}就如，那些与\Person{克罗帕}同路的人\EditorialNote{\BibleRef{路 24}}，虽然起初由于缺乏知识而软弱，但后来被救主的话点燃，结出了认识祂的果实。圣保禄亦然，当他被这火抓住时，没有与血肉商量，而是体验了这恩宠，成为圣言的宣讲者。但那九个被洁净了癞病，却对救治他们的主不知感恩的人，却不是这样；犹达斯也是如此，他得到了宗徒的位分，被称为主的门徒，最后却\Quote{\BibleQuote{同救主一起吃饭时，竟举脚踢祂，成了叛徒。}}但这样的人必得到他们愚妄的应有报应，因为他们的期望将因忘恩负义而落空；因为对忘恩负义者毫无希望，最后的永火，就是为魔鬼和他的使者所预备的，等候着那些忽视天主之光的人。这就是忘恩负义者的结局。\par

5. 然而，主的忠信而真实的仆人，知道主喜爱感恩的人，便不住讚颂祂，时时向主谢恩。无论是在安逸或磨难之时，他们以感恩之心向天主献上讚颂，不计较这些短暂的事物，而敬拜主，时间的天主。古时的\Person{约伯}，拥有超越众人的坚忍，在顺境中思念这些事；在逆境中，他忍耐承受，受苦时仍感谢。谦卑的\Person{达味}也在磨难中咏唱讚颂说：\BibleQuote{我就要时时讚颂上主。} 真福\Person{保禄}也在他所有的书信中，可以说，从未停止感谢天主。在安逸时他不懈怠，在磨难中他夸耀，深知：\BibleQuote{磨难生忍耐，忍耐生老练，老练生望德，而望德不使人蒙羞。}\BibleRef{罗 5:3} 我们身为这些人的追随者，应时时感恩，绝不懈怠；尤其在此刻，当异端者煽动对我们进行的磨难之时，我们要讚颂主，并说出圣人们的话：\BibleQuote{这一切都临到了我们身上，我们却没有忘记祢。} 因为正如犹太人那时，虽受到厄东人帐棚的攻击，被\Place{耶路撒冷}的敌人压迫，却并未放弃，反而更加歌颂天主；同样，我亲爱的弟兄们，虽然我们宣讲主的话受阻，我们更要宣扬它，受磨难时更要咏唱圣咏，因我们配为真理而受轻蔑、苦心劳碌。不仅如此，就是备受困扰，我们也要感谢。因为那位时时感谢的真福宗徒，同样劝我们亲近天主说：\BibleQuote{你们要以感恩之心，将你们的请求呈于天主。}\BibleRef{斐 4:6} 且愿我们常持此心，他又说：\BibleQuote{时时感谢，不断祈祷。}\BibleRef{得前 5:17} 因他知道信徒在感恩时变得坚强，喜乐中越过了敌人的城墙，就如那些圣人所说：\BibleQuote{赖祢，我们冲破了敌人的重围；仗赖我的天主，我跳过了城墙。} 我们要常常坚定不移，尤其在此时，尽管许多磨难临到我们，许多异端者向我们发怒。因此，我亲爱的弟兄们，我们要以感恩之情庆祝即将来临的圣节，\BibleQuote{束上我们思想的腰}\BibleRef{伯前 1:13}，效法我们的救主\Person{耶稣基督}，经上论及祂说：\BibleQuote{正义是他腰间束带，忠实是他胁下佩带。}\BibleRef{依 11:5} 我们各人手持出自\Person{叶瑟}根的杖，脚上穿着福音的准备，按\Person{保禄}所说守这节：\BibleQuote{不可用旧酵母，而要用纯洁和真诚的无酵饼。}\BibleRef{格前 5:8} 虔诚相信我们藉着基督已得和好，不偏离对祂的信德，也不与异端者和异于真理的人同流合污，他们的言谈和意志使他们堕落。反而在磨难中喜乐，冲破铁与黑暗的炉，平安度过那可怕的\Place{红海}。同样，当看到异端者混乱时，我们要与\Person{梅瑟}一起高唱伟大的讚美歌，说：\BibleQuote{我要歌颂上主，因他大显荣耀。}\BibleRef{出 15:1} 这样，歌颂讚美，眼见我们内的罪被投入海中，我们便越过进入旷野。先以四十日斋戒、祈祷、守斋、操练和善工净化自己，我们就能在\Place{耶路撒冷}享用神圣的逾越节。\par

6. 四十日斋期的开始是在\OriginalTerm{法默诺特月}{Phamenoth}第五日\EditorialNote{见：谷 1}；当我们如我所言，先以那些日子净化和准备自己之后，我们在\OriginalTerm{法尔穆提月}{Pharmuthi}第十日（4月5日）开始伟大复活节的圣周。在这圣周中，我亲爱的弟兄们，我们应当使用更长的祈祷、守斋和醒寤，好使我们能以宝贵的血涂抹我们的门楣，逃避毁灭者。\BibleRef{出 12:7} 然后，我们在法尔穆提月第十五日（4月10日）安息，因为在那星期六的傍晚，我们听到天使的讯息：\BibleQuote{你们为什么在死人中找活人呢？他已经复活了。}\BibleRef{路 24:5} 立即，那个伟大的主日迎接我们，我指的是同月即法尔穆提月的第十六日（4月11日），在这一日我们的主复活了，赐给我们与近人的和平。当我们按照祂的旨意守了这节以后，让我们从圣周的第一日起，添加上\OriginalTerm{五旬节}{Pentecost}的七周；当我们那时领受了圣神的恩宠，要时时感谢主；藉着祂，光荣和权能归于圣父，在圣神内，至于无穷之世。阿们。\par

你们要以圣吻彼此问候。与我在一起的弟兄们都问候你们。我祈祷，亲爱的、我切望的弟兄们，愿你们健康，愿你们在主内记念我们。\par

\Person{圣亚大纳修}的第三封节日书信至此结束。\par

\SourceRecord{Translated by R. Payne-Smith.；Nicene and Post-Nicene Fathers, Second Series；Vol. 4.；Edited by Philip Schaff and Henry Wace.；Buffalo, NY: Christian Literature Publishing Co.,；1892.；<http://www.newadvent.org/fathers/2806003.htm>.}

% structure: p0001 source=h1 target=section reason=page-title

\section{书信四}

\Emph{主历332年。复活节在Pharmuthi月七日，即四月诺尼日前第四日；戴克里先纪年48年；执政官：\Person{法比乌斯·帕卡提亚努斯}、\Person{麦启利乌斯·希拉里亚努斯}；长官：\Person{希吉努斯}；第五次指示周期。}\par

\Emph{他从皇帝宫廷由一名士兵送出此信。}\par

我亲爱的，我把这封信送达你们，虽较惯常时间为迟；但我深信你们会因我漫长的旅程，并因我身患疾病而予以宽恕。因这两重阻碍，且遭遇异常猛烈的暴风雨，我延迟了致书于你们。然而，尽管长途跋涉及沉疴缠身，我仍未忘记向你们发布庆节通告，为尽我职分，此刻向你们宣布此庆节。因为，虽然此信的日期较惯常的宣布为晚，仍应视为合时，既然我们的仇敌因无故迫害我们，已蒙教会羞辱谴责，我们如今可高唱节日的颂赞之歌，奏响对抗\Person{法郎}的凯歌：\BibleQuote{我们要向主歌唱，因他大获全胜，将马和骑士投于海中。}\BibleRef{出 15:1}\par

2. 我亲爱的，从一庆节迈向另一庆节，实为美善；节期的集会、圣善的守夜，再度唤醒我们的心灵，并催迫我们的理智保持警醒，以静观美善之事。我们不应如同哀悼者一般度过这些日子，而是借着饱享精神的食粮，努力抑制我们肉性的贪欲。因为借着这些，我们才有力量克胜仇敌，如同有福的\Person{友弟德}\BibleRef{友 13:8}一样，她先借着斋戒和祈祷锻炼自己，然后克胜敌人，斩杀了\Person{敖罗斐乃}。又有福的\Person{艾斯德尔}\BibleRef{艾 4:16}，当毁灭即将临到全族，以色列民族危在旦夕之际，她别无他法，唯以斋戒和向天主祈祷，克胜了暴君的狂怒，将民族的危亡化为安宁。正如这些日子在以色列被尊为庆节，同样，在古代，当敌人被杀，或危害民众的阴谋被粉碎，以色列获救时，也曾制定庆节。为此，古时的真福\Person{梅瑟}规定了逾越节的大庆节，我们也庆祝它，正是因为\Person{法郎}被杀，民众脱离奴役。因为在那些时候，尤其是当压迫民众的暴君被杀之时，\Place{犹地亚}才遵守现世的庆节和节日。\par

3. 然而，我亲爱的，如今那敌对整个世界的暴君——魔鬼，既已被杀，我们便不再是赴现世的庆节，而是永远的、天上的庆节。我们不是在预像中展示它，而是在真理中亲临其境。因为旧约的子民，饱食了那不能言语的羔羊之肉，完成了庆节，并以血涂抹门框，祈求保护对抗那毁灭者。但现在我们，饱飨\OriginalTerm{圣父的圣言}{Word of the Father}，且以\OriginalTerm{新约的血}{blood of the New Testament}\BibleRef{玛 26:28}封印了我们心门的门楣，认明了救主所赐的恩宠，他曾说：\BibleQuote{看，我已授予你们权柄，使你们践踏在蛇蝎之上，并能制伏仇敌的一切势力。} 因为死亡不再为王；从此死亡已逝，生命取而代之，因我主曾言：\BibleQuote{我是生命}\BibleRef{若 14:6}。因此万有充满喜乐与欢欣；如经上记载：\BibleQuote{主为王，愿大地踊跃。} 当死亡为王时，我们\BibleQuote{坐在\Place{巴比伦}河畔，一想起熙雍即泪流满面}，悲哀哭泣，因为感受到被掳之苦；但如今死亡与魔鬼的国度既已废除，万有便全然充满喜乐与欢欣。天主不再仅仅在\Place{犹地亚}被人认识，而是在全地，\BibleQuote{他们的声音传遍普世，他们的言语达于地极。} 我亲爱的，接下来已不言而喻：我们赴此庆节，不应身披污秽之衣，而应以洁衣佩饰我们的心灵。为此，我们需要穿上我们的主耶稣\BibleRef{罗 13:14}，好能同他一起庆祝此庆节。当我们热爱德行，嫉恶如仇，当我们操练节制，克治邪淫，当我们喜好正义先于邪恶，当我们尊崇知足，心神强健，当我们不忘穷人，向所有人敞开大门，当我们扶助谦虚，憎恶骄傲时，我们便穿上了他。\par

4. 古时的以色列，借着这些事，首先在预像中奋力取胜，前来赴宴，因为那时这些事是预兆和预像。然而，我亲爱的，如今预像已得应验，预像已经成全，我们便不应再将庆节视为象征，也不该按照犹太人过时不合宜的遵守方式，上到属地的\Place{耶路撒冷}去宰杀逾越节羔羊，免得时令流逝，我们反被视为行事不合时宜；反而，我们应遵照宗徒的训令，超越预像，高唱新的赞歌。因为他们察觉这事，并与\OriginalTerm{真理的本身}{the Truth}聚集一处，便前来对我们的救主说：\BibleQuote{你愿意我们在哪里为你预备逾越节晚餐？}\BibleRef{玛 26:17}因为不再需要执行属于地上\Place{耶路撒冷}的事；庆节也不应单在那里庆祝，而是随天主的圣意在任何地方。他愿意在每一处，使\BibleQuote{在各处都有乳香和洁净的祭品献于我的名}\BibleRef{拉 1:11}。因为，尽管按历史记载，逾越节只能留在\Place{耶路撒冷}举行，然而当属于那时候的事得以完成，属于预像的事已经过去，福音的宣讲即将遍及各地时；尤其当门徒们在各处开展庆节时，他们询问救主：\BibleQuote{你要我们在哪里预备？}救主那时正将预像转变为属灵的，他应许他们不再吃羔羊的肉，而是吃他自己的，说：\BibleQuote{你们拿去，吃，喝；这是我的身体，这是我的血。}当我们如此借这些滋养，我亲爱的，我们也必真正举行逾越庆节。\par

5. 我们从\OriginalTerm{法尔穆提月}{Pharmuthi}一日（3月27日）开始，在同月的第六日（4月1日）休息，即第七天的晚上；那神圣的一周的第一日，也就是同法尔穆提月七日（4月2日），已为我们升起，我们也庆祝接着而来的神圣的\OriginalTerm{五旬节}{Pentecost}节日，通过这些日子预显将来的世界，这样我们从此就能永远与基督在一起，赞美在万有之上的天主，在基督耶稣内，并藉着祂，同众圣人一起，向主说：阿们。你们要以圣吻彼此致意。所有与我在一起的弟兄都问候你们。我们从\Place{宫中}寄出这封信，经由一位侍从官员之手，是那真正敬畏天主的\OriginalTerm{禁卫军长官}{Præfect of the Prætorium}\Person{阿布拉维乌斯}交给他的。因为我正在\Place{宫中}，是受\Person{君士坦丁}皇帝召见。但那里的\OriginalTerm{梅勒提乌斯派}{Meletians}心怀嫉妒，在皇帝面前图谋败坏我们。然而，他们被许多事驳斥，蒙羞受辱，被当作诽谤者从那里赶了出去。被赶走的人是\Person{卡利尼库斯}、\Person{伊西翁}、\Person{尤德蒙}，以及\Person{格洛乌斯·希拉坎蒙}，他因自己名字的可耻，自称\Person{欧洛吉乌斯}。\par

圣\Person{亚大纳削}的\WorkTitle{第四封节日书信}到此结束。\par

\SourceRecord{Translated by R. Payne-Smith.；Nicene and Post-Nicene Fathers, Second Series；Vol. 4.；Edited by Philip Schaff and Henry Wace.；Buffalo, NY: Christian Literature Publishing Co.,；1892.；<http://www.newadvent.org/fathers/2806004.htm>.}

% structure: p0001 source=h1 target=section reason=page-title

\section{第五封信}

\Emph{时为333年复活节，执政官达尔马提乌斯与泽诺菲卢斯；总督帕特努斯；第六征税期；五月朔日前十七日，即法穆提月二十日；月龄十五；神祇七日；戴克里先纪元49年。}\par

我的弟兄们，我们按时从节期走向节期，按时从祈祷走向祈祷，我们从斋戒走向斋戒，将圣日连接圣日。那带给我们一个新开始的时刻又到了，即宣告那至福的逾越节，在那时主曾被祭献。我们就像吃下生命之粮，且因不断渴慕而时时畅饮祂的宝血，如同从泉源一般，使灵魂欢欣。因为我们持续而热切地渴望；祂常为渴者预备；而对渴者，有我们救主的话，就是祂在节期那天，以祂的慈爱说出的：\BibleQuote{人若渴了，到我这里来喝吧！} 并不是只在有人就近祂时，祂才解除他的渴；而是无论何时有人寻求，都可自由地来到救主面前。因为节期的恩宠不限于一时，其辉煌的光辉也不消减；它常近在眼前，光照那些切望者的心灵。因为其中蕴藏着恒久的力量，为那些心灵受到光照，并昼夜默思\WorkTitle{圣经}的人，就像那位领受祝福的人，正如在\WorkTitle{圣咏集}中所记载的：\BibleQuote{凡不随从恶人的计谋，不插足于罪人的道路，不参与讥讽者的席位，而专心爱好上主法律的，昼夜默思祂的法律的，才真是有福的！} 因为光照他的不是太阳、月亮，或众多星辰，他乃是闪耀着超越万有的天主至高光辉。\par

2. 因为我亲爱的，正是天主，那位最初为我们设立此节期的主，赐予我们年复一年地庆祝它。祂既使祂的圣子被牺牲以带来救恩，又赐给我们这神圣节期的缘由，每年此季节宣布节期时，都为这事作见证。这节期也引领我们从十字架经过这个世界，走向那摆在我们前面的；天主现今也从其中生出光荣救恩的喜乐，把我们带到同一的聚会中，并在每一处使我们在精神上合而为一；又为我们指定共同的祈祷，以及由此节期涌出的共同恩宠。这就是祂慈爱的奇妙之处：祂使远隔各处的人聚集在同一地方；并使那些身体上显得遥远的人，在精神的合一中相近相亲。\par

3. 所以，我亲爱的，我们为何不照节期所宜的感谢这恩宠呢？我们为何不向我们的恩主回报呢？对天主做出相称的回报确实不可能；然而，我们领受了恩惠的礼物，却不承认它，实在邪恶。本性本身就显明我们的无能为力；但我们的意志却谴责我们的忘恩负义。因此，圣保禄赞叹天主恩赐的伟大时，曾说：\BibleQuote{谁堪当此任呢？} \BibleRef{格后 2:17} 因为祂借着救主的宝血使世界获得自由；然后，祂又借着救主的死亡，使坟墓被踏平，为那些上升的人开辟了一条通达天门的无阻之路。所以，圣人们中的一位承认这恩宠，却无力回报，说道：\BibleQuote{我应何以为主报谢他所赐我的一切厚恩？} 因为他以生命代替了死亡，以自由代替了奴役，以天国代替了坟墓。从前，\BibleQuote{死亡自亚当掌权直到梅瑟}；但如今，天主的声音却说：\BibleQuote{今天你就要与我同在乐园。} 圣人们体认到这事，便说道：\BibleQuote{若非上主扶助我，我的灵魂几乎已居阴间..} 除此以外，虽然他无力回报，却仍承认这恩赐，最后写道：\BibleQuote{我要举起救恩的杯，呼求上主的名；在祂眼中，圣人们的逝世极为珍贵！}\par

关于这杯，主曾说：\BibleQuote{你们能饮我将要饮的杯吗？} 门徒们答应了，主就说：\BibleQuote{我的杯你们固然要饮；但坐在我右边或左边，不是我可以给的，而是我父给谁预备了，就给谁。} \BibleRef{玛 20:22} 因此，我亲爱的，我们既无力回报，就该领受这恩赐。让我们按自己的力量，把握时机。因为尽管本性不能以不配主圣言的事物来回报如此鸿恩，但我们在坚持虔敬时，仍应感谢祂。我们如何能更持久地保持虔敬，胜于认明天主，因祂对人类的爱，赐予我们如此恩惠？（这样我们便顺从遵守法律，并谨守其诫命。而且，我们也不会像忘恩负义之人，被算为法律的违犯者，或行那些可憎之事，因为主喜爱感恩的人）；当我们将自己献给主，如同圣人们那样，完全把自己献出，从此不再为自己而活，而是为那为我们而死的主而活，如同圣保禄所做的那样，他曾说：\BibleQuote{我已与基督同钉十字架，所以我生活，已不是我生活，而是基督在我内生活。} \BibleRef{迦 2:20}\par

4. 弟兄们，我们生命的真谛是在于舍弃一切肉身之事，只坚定持守我们救主的事。因此，当前的时期要求我们，不仅应说出这样的话，还应效法圣人的行为。而我们效法他们，就是当我们承认那位死去的，不再为自己而活，而是让基督由此在我们内生活；就是当我们尽力向主献上回报，尽管我们回报时献上的并非我们自己的所有，而是我们先前从祂所领受的，这尤其出于祂的恩宠，因为祂竟要求我们交还祂自己的恩赐，仿佛是从我们献上的。祂亲自为此作证说：\Quote{我的祭献就是我自己的恩赐。}就是说，你们献给我的，原是你们的，因为是从我领受的，但这些都是天主的恩赐。且让我们将一切德行和那在祂内的真圣德献给主，并以虔诚之心，用祂为我们祝圣的事物，向祂守节。让我们就这样投身于圣斋期，因为这斋期是祂所规定的，且借着它我们找到通往天主的道路。但不要让我们像外教人、无知的犹太人，或当今的异端者和裂教者一样。因为外教人以为节日的成全在于食物的丰盛；犹太人迷失于预像和影子，仍以为如此；裂教者在分离的地方，怀着虚幻的想象守节。但是，弟兄们，让我们胜过外教人，即以心灵的真挚和身体的纯洁守节；胜过犹太人，即不再接受预像和影子，而是如同被真理之光荣耀地照耀，并瞻仰那正义的太阳\BibleRef{拉 4:2}；胜过裂教者，即不撕裂基督的长袍，而是在同一屋宇——即在\OriginalTerm{天主教会}{Catholic Church}内，吃主的逾越节。是祂制定了神圣的法律，引领我们走向德行，并劝勉我们在这节期中斋戒。因为逾越节确实是戒除邪恶以修德行的节期，是由死亡进入生命。这一点甚至可从古时的预像得知。因为那时他们奋力从\Place{埃及}前往\Place{耶路撒冷}，如今我们却是由死亡进入生命；他们那时是从\Person{法郎}转向\Person{梅瑟}，如今我们却是从魔鬼转向\Person{救主}。正如那时，解救的预像每年作证，如今我们纪念我们的救恩。我们斋戒，默想死亡，为能获得生命；我们警醒，不像哀恸者，却像那些等待主从婚宴归来的人，好使我们在凯旋中互相争先，急于宣告战胜死亡的标记。\par

5. 因此，我亲爱的，但愿我们如\OriginalTerm{圣言}{Word}所要求的，在此时此地就完全地约束自己，并如此生活，总不忘却天主的高尚作为，也不偏离德行的实践！正如宗徒之声也劝勉说：\BibleQuote{忆念耶稣基督，祂从死者中复活了}\BibleRef{弟后 2:8}。这并非规定了某一限制的纪念时期，因为祂应常在我们思想中。但因许多人的懈怠，我们日复一日地拖延。那么就让我们在这些日子里开始吧。为此，准许一段纪念的时期，为能向圣人彰显他们蒙召的赏报，并在责备疏忽者的同时劝勉他们。因此，在余下的所有日子里，让我们恒心行善，按本分悔改我们所忽略的一切，无论是什么；因为没有人能免于污秽，即使他在世上的时日不过一小时，正如那位毅力超群的\Person{约伯}所证实的。然而，\Quote{伸向前方之事}，我们应祈求，使我们不致不相称地吃这逾越节，免得陷入危险。因为以纯洁守这节期的人，逾越节便是天粮；而亵渎轻慢守这节期的人，它却是危险和耻辱。因为经上记载：\BibleQuote{凡不配地吃喝的人，就是干犯主的身体和血}\BibleRef{格前 11:27}。因此，我们不要仅仅去履行节期的礼仪，而要准备好亲近\OriginalTerm{天主羔羊}{divine Lamb}，并触及\OriginalTerm{天粮}{heavenly food}。让我们洁净双手，净化身体。让我们保持整个心灵远离欺诈；不纵情于放荡和私欲，而是全身心专注于我们的主，和神圣的教义；这样，成为完全纯洁，我们才能领受\OriginalTerm{圣言}{Word}。\par

6. 我们在法尔穆提月十四日（4月9日），即一周的[第一]晚开始圣斋期；并于同月法尔穆提月十九日（4月14日）结束，圣周的第一日于同月法尔穆提月二十日（4月15日）向我们破晓，我们加上五旬节的七个星期；伴随着祈祷、与邻人的共融、彼此相爱，以及那超越一切的和平意愿。因为这样，我们才能成为天国的继承人，靠赖我们的主耶稣基督，愿光荣和权能借着祂归于圣父，至于无穷之世。阿们。所有与我在一起的弟兄都问候你们。你们要以圣吻彼此问候。\par

\Person{圣亚大纳削}的\WorkTitle{第五封节期书信}于此结束。\par

\SourceRecord{Translated by R. Payne-Smith.；Nicene and Post-Nicene Fathers, Second Series；Vol. 4.；Edited by Philip Schaff and Henry Wace.；Buffalo, NY: Christian Literature Publishing Co.,；1892.；<http://www.newadvent.org/fathers/2806005.htm>.}

% structure: p0001 source=h1 target=section reason=page-title

\section{书信十}

\Emph{时为338年，执政官\Person{乌尔苏斯}与\Person{波勒米乌斯}；总督仍为\Place{赫利奥波利斯}人\Person{狄奥多鲁斯}，亦为公教徒。} \Emph{其后，\Person{菲拉格里乌斯}续任第二年；\OriginalTerm{小纪}{Indiction}十一；复活节在\OriginalTerm{四月朔日}{Kal. Ap.}前第七日，即\OriginalTerm{法梅诺斯月}{Phamenoth}三十日；月龄十八又半；\OriginalTerm{戴克里先纪元}{Æra Dioclet.}五十四年。}\par

尽管我已远离你们，我的弟兄们，我仍未忘记在你们中间通行的、由教父们传给我们的惯例，即不可缄默不语，而不将年度圣节的时间及其庆祝之日通知你们。因为，虽然那些苦难——你们无疑已有所闻——阻碍了我，严峻的考验加在我身上，遥远的距离又将我们隔绝；同时真理的仇敌追踪我们的足迹，设下陷阱以查获我们的书信，好借他们的控告加深我们创伤的痛楚；然而，主在我们的磨难中坚固并安慰我们，因此我们虽身处这等阴谋与圈套之中，也毫不畏惧，仍要从地极向你们指示并宣报这使人得救的复活节。再者，当我写信给\Place{亚历山大里亚}的司铎时，我力劝他们借由这些书信转达给你们，尽管我知悉反对者加给他们的恐惧。但我仍勉励他们记念宗徒的直言之勇，并说：\BibleQuote{没有什么能使我们与基督的爱相隔绝：是困苦吗？是窘迫吗？是迫害吗？是饥饿吗？是赤贫吗？是危险吗？是刀剑吗？} \BibleRef{罗 8:35}。同样，我既自己庆祝这节筵，也切愿你们，我所亲爱的，同样庆祝；我自知应做这样的通报，便不敢迟延履行此责，以免因宗徒的劝谕受审判：\Quote{凡人应得的，你们要付清。}\par

2. 我既将一切事务交托于天主，便切愿与你们一同庆祝这节筵，不顾及我们之间的距离。因为尽管空间将我们分隔，但那位节筵的赐予者，且祂自己就是我们的节筵，祂也是圣神的赋予者，却借着心灵的契合、和谐与和平的连系，将我们聚集在一起。因为我们若同心同思，并为彼此献上同样的祈祷，就没有空间能隔绝我们，反而是主将我们召集并联合为一。因为祂曾应许：\BibleQuote{那里有两三个人，因我的名聚在一起，我就在他们中间} \BibleRef{玛 18:20}，那么显然，祂既在每一处聚在一起的人中间，就联合他们，并收纳他们众人的祈祷，如他们就在身边，又垂听他们同声高呼的\Quote{阿们}。我的弟兄们，当我写信给你们时，已承受了与此相似的苦难，并我所提及的一切考验。\par

3. 为免你们全然忧愁，我现今只简短提醒你们这些事；因为人若稍得安乐，便忘记在磨难中所受的苦楚，那是不可取的；免得他如同忘恩负义者，被排除于神圣集会之外。因为人颂扬天主，莫过于他历经苦难之后；同样，感谢天主，也莫过于他脱离劳苦与试探、获得安息之时。就如\Person{赫则克雅}，当\Place{亚述}人灭亡后，赞美主，称谢说：\BibleQuote{主是我的救援；我一生岁月，在主的殿内，不断弹琴颂扬。} \BibleRef{依 38:20}。那在\Place{巴比伦}受试炼的、英勇而蒙福的三人——\Person{哈纳尼雅}、\Person{米沙耳}、\Person{阿匝黎雅}——当他们得保安全，火焰为他们化为甘露时，便称谢、赞美并\Quote{向天主说出颂祝的言辞}。我的弟兄们，我怀着这些心思，也像他们一样给你们写信；因为即使在我们这个时代，天主，人所不能的，祂使之成为可能。那人所不能成全的事，主却显示为易成，祂以此带领我们来到你们中间。因为祂并不把我们交于那欲吞噬我们的人作掠物。其实他们所图谋以邪恶湮没的，并不是我们，而是教会、信仰与虔敬。\par

4. 然而，美善的天主向我们丰厚地赐下祂的\OriginalTerm{慈爱}{loving-kindness (hesed)}，不仅是在祂藉着自己的\OriginalTerm{圣言}{Word (Logos)}赐予我们众人共同的\OriginalTerm{救恩}{salvation}之时，更是在现在，当仇敌迫害我们、企图捉拿我们的时候。正如真福\Person{保禄}在某处描述基督那深不可测的丰饶时所说的：\BibleQuote{然而富于仁慈的天主，因着祂爱我们的大爱，竟在我们因过犯和罪恶死了的时候，使我们同基督一起活过来 \BibleRef{弗 2:4}}。因为人的力量和一切受造物的力量，是软弱而贫乏的；但那超越人、非受造的\OriginalTerm{德能}{Power (Dynamis)}，却是丰裕而深不可测的，没有起始，永恒常在。因此，祂不仅拥有一种医治的方法，反而因着祂的丰裕，藉着自己的圣言，以各种方式为我们的救恩而工作；祂对待我们时，不受任何限制或阻碍；正是因为祂丰裕而多彩，祂根据每个灵魂的不同容量而改变自己。因为祂是天主的圣言、德能和\OriginalTerm{智慧}{Wisdom (Sophia)}，正如\Person{撒罗满}论智慧时作证说：\BibleQuote{智慧虽是唯一的，却无所不能；她虽永存不变，却使万象更新；她进入圣善的灵魂，使他们成为天主的朋友和先知。}因此，对那些尚未达到完美道路的人，祂就像一只哺乳的母羊，这是由保禄所施行的：\BibleQuote{我给你们喝的奶，不是饭 \BibleRef{格前 3:2}}。对那些已超越了孩童的完全身量，但在成全上仍软弱的人，祂按他们的能力成为他们的食物，这同样是保禄所施行的：\BibleQuote{软弱的人可以吃蔬菜}。但人一旦开始行走在完美的道路上，就不再以先前的事物为食，而是以圣言为饼，以肉身为食，因为经上记载着：\BibleQuote{硬食是成年人的食物，他们的官能因练习而熟练，能分辨善恶 \BibleRef{希 5:14}}。再者，当圣言被播种时，它在人类生命中并不产生单一果实，而是多种而丰盛的；因为它有的结一百倍，有的六十倍，有的三十倍，正如\Person{救主}——那位\OriginalTerm{恩宠}{grace}的播种者，\OriginalTerm{圣神}{Spirit}的赋予者——所教导的。这并非可疑之事，也不是不可证实的；我们确实有能力看见祂所播种的田地；因为在\OriginalTerm{教会}{Church}中，圣言是多样的，果实是丰盛的。这样的田地，不仅以\OriginalTerm{贞女}{virgins}为装饰，也不仅以\OriginalTerm{隐修士}{monks}为装饰，更以尊贵的婚姻和各人的贞洁为装饰。因为播种时，祂并未强迫意志超越能力。仁慈也不仅限于成全的人，也降临于居中的和第三等级的人之中，好使祂普遍地拯救所有人进入救恩。为此，祂在父那里预备了许多住处 \BibleRef{若 14:2}，这样，尽管住所因道德成就的进步而各有不同，但我们众人都在同一围墙之内，都进入同一栅栏之中，那仇敌已被赶出，其全军也从那里被驱逐。因为离开光明便是黑暗，离开祝福便是诅咒，魔鬼也与圣人隔离，罪恶远离德行。因此，福音斥责\Person{撒殚}说：\BibleQuote{撒殚，退到我身后去！\BibleRef{玛 4:10}} 但它召唤我们进入，说：\BibleQuote{你们要从窄门进去}。又说：\BibleQuote{我父所祝福的，你们来吧，承受给你们预备的国度吧}。同样，圣神从前也在\BibleBook{}{圣咏集}中呼喊说：\BibleQuote{你们要以咏颂进入祂的门庭}。因为人藉着德行进入天主，如\Person{梅瑟}进入天主所在的浓云中。但人因恶习而离开上主的面；如\Person{加音} \BibleRef{创 4:16} 在杀死兄弟后，就他内心的意向而言，已离开了天主的面；而圣咏作者却进入说：\BibleQuote{我要走向天主的祭台，走向让我喜乐的天主}。关于魔鬼，圣经作证说，魔鬼离开天主面前后，打击了\Person{约伯}，使他生出恶疮。这就是那些离开天主面前之人的特征——击打和伤害天主的人。这也是那些背弃信德之人的特征——伤害和迫害忠信者。相反，圣人接纳他们，视他们为朋友；正如真福\Person{达味}也曾坦然地说：\BibleQuote{我的眼目注视地上的忠信者，使他们与我同住}。但对于信德软弱的人，保禄劝我们尤其要接纳他们。因为德行是爱人的，正如在品性相反的人身上，罪恶是憎恨人的。因此，身为罪人的\Person{撒乌耳}迫害\Person{达味}，而\Person{达味}虽有好机会，却不杀\Person{撒乌耳}。\Person{厄撒乌}也迫害\Person{雅各伯}，而\Person{雅各伯}却以温和胜过了他的邪恶。那十一个弟兄卖了\Person{若瑟}，\Person{若瑟}却以慈爱怜悯了他们。\par

但我们何须多言？我们的主和救主，当祂受法利塞人迫害时，却为他们的丧亡哭泣。祂受了凌辱，却不威胁；受苦难时不威胁，甚至被杀害时也不威胁。反倒为那些胆敢如此行的人哀恸。祂这位救主，为人类受苦，他们却轻视并抛弃了生命、光明与恩宠。这些本都藉着那位替我们受苦的救主而属于他们。的确，祂为他们的黑暗与盲目哭泣。因为他们若明白圣咏中所载的，就不会如此狂妄地反对救主，圣神曾说：\BibleQuote{万邦为什么嚣张，众民为什么妄想？}\BibleRef{咏~2:1}。他们若思想过\Person{梅瑟}的预言，就不会把那位本是他们的生命者钉死。他们若用心查考所记载的，就不会故意应验那反对他们自己的预言，致使他们的城如今一片荒凉，恩宠被夺去，他们自己失去了法律，不再被称为子女，而是外方人。因为圣咏中早已如此宣示说，\BibleQuote{外邦的子民竟向我行诈。} 又藉\Person{依撒意亚}先知说：\BibleQuote{我养育了子女，抚养了他们，他们竟背叛了我。}\BibleRef{依~1:2} 他们不再被称为天主的子民、圣洁的邦国，反而是\Place{索多玛}的统治者、\Place{哈摩辣}的百姓；甚至在这一点上，他们超过了索多玛人的罪恶，正如先知也说道：\BibleQuote{索多玛比你更有正义。}\BibleRef{则~16:52} 因为索多玛人是对天使发狂，这些人却是对抗万有的主、天主和君王，并且竟敢杀害天使的主，却不知道那被他们杀害的基督仍然活着。但那些合谋反对主的犹太人却死了，他们在这些短暂的事物中稍享欢乐，却丧失了永恒的事物。他们不明白这一点——不朽的应许并不顾及暂时的享乐，而在于对那些永恒事物的希望。因为圣人经过许多患难、劳苦和忧闷，才进入天国；但当他到达那里，痛苦、忧伤、叹息都将逃遁时，他从此就享受安息；正如\Person{约伯}，在此世受试炼之后，后来成为主的密友。但那贪爱享乐的人，一时欢愉，日后却度悲痛的生活；就如\Person{厄撒乌}，曾享有现世的食物，后来却因此被定罪。\par

我们可以以色列子民与\Place{埃及}人出离\Place{埃及}的事，作为这种区别的预象。因为埃及人因对以色列的不义而暂享欢乐，当他们出去之时，全都淹死在深海之中；但天主的子民，虽被监工们暂时打击和伤害，当他们出了\Place{埃及}，却安然渡过大海，并在旷野中行走如同居住之境。因为那地方虽无人迹，一片荒凉，但因着法律的恩赐，以及天使的交往，便不再是旷野，反而远胜于人烟之地。正如\Person{厄里叟}，当他以为自己在旷野中孑然一身时，却有众天使与他同在；同样在这情形下，虽然百姓起初受苦难、在旷野中，但那些保持忠信的人最终进入了预许之地。同样，那些在此世遭受暂时苦难的人，最终忍耐到底，获得安慰，而在此世迫害他人者，却被践踏，没有好结局。就如福音所证实的，那位富家人\BibleRef{路~16:19}，在此世暂享欢乐之后，在那里忍饥受饿，在此世开怀畅饮之后，在那里极度口渴。但\Person{拉匝禄}，在世俗事物中受苦之后，在天堂得到安息；他饥饿想吃麦粒磨成的饼，在那里却饱饫了那比玛纳更好的，即主，祂曾降下来说：\BibleQuote{我是从天上降下来的食粮，赐给人类生命。}\BibleRef{若~6:51}。\par

哦！我极亲爱的诸位，如果我们能从苦难中获得安慰，从劳苦中获得安息，从疾病中获得健康，从死亡中获得不朽，那么就不该为临于人类的暂时灾祸而愁苦。我们不该因遭遇试探而慌乱。如果那与基督为敌的团伙合谋反对虔敬，我们也不该惧怕；反而我们应通过这些事更加悦乐天主，并将这些事视为对善德生活的考验和操练。因为，若没有先前的劳苦与忧闷，又怎能期望忍耐呢？若没有仇敌的攻击，刚毅又怎能被试探呢？若没有侮辱与不义，心胸宽广又怎能彰显呢？若没有反基督者的谤毁在先，坚忍又怎能被证明呢？最后，若没有极为邪恶者的不义先行显露，一个人又如何能以肉眼看见美德呢？因此，我们的主和救主\Person{耶稣基督}站在我们面前，当祂要指示人如何忍受苦难时，祂被击打却耐心忍受，受人辱骂却不还骂，受苦时也不威胁，却将自己的背交给击打者，将脸颊交给掌掴者，也不掩面躲避唾污\BibleRef{伯前~2:23}；最后，甘愿被领往死亡，好使我们在祂身上看见一切美德与不朽的肖像，并使我们效法这些榜样，真正踏在蛇蝎和仇敌的一切势力之上。\par

8. 同样，\Person{保禄}在效法主的榜样时，也劝勉我们说：\BibleQuote{你们要效法我，如同我效法了基督一样。}\BibleRef{格前 11:1} 他这样得胜了魔鬼的一切分裂，写道：\BibleQuote{因为我深信：无论是死亡，是生活，是天使，是掌权者，是现存的，或是将来的事物，是有权能者，是崇高或深远的势力，或其他任何受造之物，都不能使我们与天主的爱相隔绝，即是与我们的主耶稣基督之内的爱相隔绝。}\BibleRef{罗 8:38} 因为仇敌借着患难、试探和劳苦临近我们，千方百计要毁灭我们。然而在基督内的人，与那些反对之事争斗，以坚忍对抗忿怒，以温柔对抗侮辱，以德性对抗邪恶，便获得胜利，并宣告说：\BibleQuote{我赖加强我力量的那位，能应付一切。}又说：\BibleQuote{靠着那爱我们的主，我们在这一切事上，大获全胜。}这就是主的恩宠，也是主为人类儿女恢复的方法。因为他受苦，为那些在他内受苦的人预备免于苦难的自由；他降下，为把我们举扬起来；他承担了诞生之难，为使我们能爱那\OriginalTerm{非受生}{unbegotten}者；他进入腐朽，为使腐朽穿上不死不灭；他为我们而成为软弱的，为使我们在能力中复活；他下降至死亡，为赐予我们不死不灭，并将生命赋予死者。最后，他成为人，好使我们这些如人般死亡的人能再活过来，并使死亡不再统治我们；因为宗徒的言论宣称：\BibleQuote{死亡不再统治我们了。}\par

9. 但如今，因为他们不如此考虑这些事，那些\OriginalTerm{亚略派狂热者}{Ario-maniacs}，作为基督的敌对者和异端者，用自己的舌头打击那位援助他们的主，亵渎那位曾解救他们的主，并对救主持有种种不同的意见。因他为了世人而降下，他们便否认了他本有的天主性；见到他生于童贞女，便怀疑他是否真是天主子；认为他在时间内降生成人，便否认了他的永恒性；看待他为我们受苦，便不相信他是由不可朽坏的圣父所生的不可朽坏的圣子。最后，因他为我们忍受了一切，他们便否认了与祂本有永恒有关的事；这些忘恩之人行事，轻看救主，不但不承认他的恩宠，反而侮辱他。对他们，可以正当地说这样的话：哦！你这忘恩的反对基督者，全然邪恶，杀死自己主人的凶手，心灵盲目，思想如犹太人，如果你曾明白圣经，听从众圣者所说：\BibleQuote{愿你的慈颜光照我们，使我们得救。}或又说：\BibleQuote{求你发出你的光明和你的真理。}那么你便会知道，主并非为自己，而是为我们而降下；因此，你倒更会敬佩他的慈爱。假若你曾思虑过圣父是什么，圣子是什么，你就不会亵渎圣子，好像祂有可变的性体。假若你曾理解祂对我们的慈爱工程，你就不会使子远离父，也不会视那位使我们与父和好的主为外人。我知道这些话不但令那些与基督争辩的人感到难堪，也令那些裂教者难堪；因为他们同气相求，如同志趣相投的人。因为他们学会了撕裂天主的无缝长袍：他们认为分裂不可分割的父与子并非怪事。\par

10. 我确实知道，当这些事情被宣讲时，他们会对我们咬牙切齿，连同那煽动他们的魔鬼，因为他们因宣讲关于救赎主的真正荣耀而烦恼不安。但那位时常嘲笑魔鬼的主，如今也这样做，他说：\BibleQuote{我在父内，父在我内。}\BibleRef{若 14:11} 这就是那位在父内显现，而父也在他内显现的主；他，真正是圣父的儿子，最终为了我们而降生成人，好能替我们把自己献给父，并借着他的奉献和祭献救赎我们。这就是那曾引领古时人民离开埃及的那位；但后来他把我们众人，或更好说全人类，从死亡中救赎出来，并将他们从坟墓中领出。这就是那位在古时曾作为羔羊被祭献，借羔羊所预表的那位；但后来他为我们被宰杀，因为\BibleQuote{我们的逾越节羔羊基督，已被祭杀作了牺牲。}\BibleRef{格前 5:7} 这就是那位救我们脱离猎人的网罗，我说离开那些反对基督者，并离开裂教者，再次拯救他的教会的那位。因为我们当时是欺骗的受害者，如今他借着自己亲自拯救了我们。\par

11. 那么，我的弟兄们，为了这些事，我们的本分岂不是要赞美并感谢万有的君王天主吗？首先让我们用\BibleBook{}{圣咏}的话高呼：\BibleQuote{愿主受赞美，因祂没有将我们交出，成为他们牙齿的猎物。} \BibleRef{咏 124:6} 让我们按祂为我们的得救而祝圣的方式守这节日——即神圣的复活节——以便我们能同天使一起庆祝那在天上的节日。古时，犹太人的百姓从苦难进入安逸时，便守这节日，为他们的胜利高唱赞美歌。同样，在艾斯德尔的时代，百姓因脱免了死刑的命令，也向主守了节日，视之为庆节，向主谢恩，并赞美祂，因为祂改变了他们的处境。因此，让我们向主履行誓言，承认我们的罪过，在言谈、品行和生活方式上向主守节；赞美我们的主，祂虽稍微惩戒了我们，却没有全然不顾或离弃我们，也没有向我们完全沉默。因为祂既带领我们脱离反基督者的骗人的、著名的埃及，使我们经过许多考验和磨难，如同在旷野中，进入祂的圣教会，以致从此以后，我们可以照常给你们写信，并收到你们的信；为此，我特别亲自感谢天主，并劝你们同我一起、也为我的缘故感谢祂，这是宗徒的习惯，而那些反基督者和裂教者却想终止并断绝这习惯。主不允许这事，反而恢复并保存了祂藉宗徒所规定的事，使我们得以一同守节，一同度圣日，按照教父们的传统和命令。\par

12. 我们在梅基尔月十九日（二月十三日）开始为期四十天的斋戒；在法梅诺特月二十四日（三月二十日）开始神圣的复活节斋戒。我们在法梅诺特月二十九日（三月二十五日），即第七日深夜，停止斋戒。这样，我们在法梅诺特月三十日（三月二十六日）黎明的一周的第一天守这节日；从这天起，直到五旬节，我们度过圣日，连续七周，一周接一周。因为当我们先恰当地默想了这些事，我们将能堪当配得那些永恒的事，赖我们的主耶稣基督，藉着祂，愿光荣和权能归于圣父，直到无穷世之世。阿们。你们要以圣吻彼此问候，在你们的圣祷中纪念我们。所有与我在一起的弟兄都问候你们，时刻纪念你们。我祈求你们在主内得享健康，我亲爱的弟兄们，我们爱你们胜过一切。\par

\Person{圣亚大纳削} 的第十封信至此结束。\par

\SourceRecord{Translated by R. Payne-Smith.；Nicene and Post-Nicene Fathers, Second Series；Vol. 4.；Edited by Philip Schaff and Henry Wace.；Buffalo, NY: Christian Literature Publishing Co.,；1892.；<http://www.newadvent.org/fathers/2806010.htm>.}

% structure: p0001 source=h1 target=section reason=page-title

\section{第十一封信}

\Emph{为 339 年。\OriginalTerm{执政官}{Coss.}：\Person{康士坦提乌斯·奥古斯都二世}、\Person{康士坦斯一世}；\OriginalTerm{总督}{Præfect}，\Person{斐拉格里乌斯}（卡帕多西亚人），第二次任职；\OriginalTerm{第 12 小纪}{Indict. xii}；\OriginalTerm{复活节}{Easter-day} \OriginalTerm{五月前第十七日}{xvii Kal. Mai}，\OriginalTerm{法穆蒂月二十日}{xx Pharmuthi}；\OriginalTerm{戴克里先纪元 55 年}{Æra Dioclet. 55}.}\par

有福的保禄，身被一切德行，被称为\OriginalTerm{主的忠信者}{faithful of the Lord}——因为他自觉在他内只有美德和称赞，或与爱德和虔敬相谐的事——他越来越依赖这些事，甚至被提升到天上，被携往\Place{乐园}\BibleRef{格后 12:4}；以至于，他既超越了人的交往，就被高举于众人之上。当他下来时，他向每个人宣讲：\BibleQuote{我们如今所知道的只是局部，我们作先知所讲的也只是局部；如今我只知道局部；但那时我就要完全知道，如同我完全被知道一样}\BibleRef{格前 13:9}。因为事实上，他被那些在天上的圣人所认识，视他为同城的公民。就一切未来和圆满的事而言，他在这里所知道的只是局部；但就那些由主交托并委任于他的事而言，他是完全的；正如他所说：\BibleQuote{我们凡是完全的人，都该怀有这种心意}\BibleRef{斐 3:15}。因为正如基督的福音是那由\Place{以色列}法律所提供的职事的实现和完成，同样，未来的事将是现今存在之事的完成，那时福音得以完全，信者们将领受那些如今未看见，却仍盼望的事物，正如保禄所说：\BibleQuote{人若看见了，还希望什么呢？但如果我们希望那未看见的，就必须耐心等待}\BibleRef{罗 8:24}。既然那有福之人具有这样的品性，且宗徒的恩宠托付给了他，他就写下了，切望\BibleQuote{所有的人都如同他一样}\BibleRef{格前 7:7}。因为德行是博爱的，而天国的团体是庞大的，因为在那里有千千万万的万军事奉主。人虽然经过一条狭窄而艰难的道路进入其中，然而一进入，便看见无边无际的空间，一个比任何其他地方都大的地方，如同那些亲眼目睹并继承这些事的见证人所宣告的。\Quote{你曾使我们遭受苦难。}但后来，在述说了他们的苦难之后，他们又说：\Quote{你使我们进入了宽广之地；}又说：\Quote{在苦难中，你宽慰了我们。}的确，我的弟兄们，圣人们在此世的路程是狭窄的；因为他们或因渴望未来的事而劳苦，正如那说：\Quote{我寄居的时日拖延，我有祸了！}的人；或是为别人的得救而忧伤耗尽，如同保禄写信给格林多人说：\BibleQuote{免得我再来的时候，天主叫我在你们面前羞愧，并且我为那许多从前犯了罪，却不悔改他们所行的污秽、奸淫和放荡的人而哀恸}\BibleRef{格后 12:21}。正如\Person{撒慕尔}为\Person{撒乌耳}的毁灭而哀恸，\Person{耶肋米亚}为百姓的被掳而哭泣。但在经历了这苦难、忧伤和叹息之后，当他们离开此世时，某种属神的喜乐、愉悦和欢欣便会迎接他们，而患难、忧伤和叹息都要逃遁远离。\par

我的弟兄们，我们既然处于这样的境况，就绝不要在德行的道路上懈怠；因为为此他劝勉我们说：\BibleQuote{你们该效法我，如同我效法了基督一样}\BibleRef{格前 11:1}。因为他给出这劝告不仅是为格林多人，既然他不只是他们的宗徒，而是\BibleQuote{在信仰和真理上做了外邦人的教师}\BibleRef{弟前 2:7}，他便藉着他们劝诫我们众人；总之，他写给每个人的，是对众人共同适用的诫命。因此，他在写给不同的人时，对有些人是劝勉，例如在\BibleBook{罗马}{书}、\BibleBook{厄弗所}{书}和\BibleBook{费肋孟}{书}中；对有些人是谴责和愤慨，如在格林多人和迦拉达人的情况中；对有些人给予劝戒，如对哥罗森人和得撒洛尼人；对斐理伯人他则赞许喜悦；对希伯来人他教导他们法律只是影子。但对于他所拣选的子女\Person{弟茂德}和\Person{弟铎}，当他们在近处时，他便给予指导；在远处时，便提醒他们。因为他对一切人成为一切；他自己既是完全的人，便使他的教导适应每个人的需要，为能尽一切方法拯救其中一些人。因此他的话语并非不结果实；而是在各处都被栽种，并直到今日仍结实累累。\par

我亲爱的诸位，这是为什么？因为我们理当探究宗徒的心意。不仅在书信的开端，也在其结尾和中间，他都运用了劝告和训诫。因此我希望，藉着你们的祈祷，我绝不会错误表述那位圣者的计划。他既精通这些神圣之事，且知晓神圣教导的能力，就认为必须首先传扬关于基督的圣言，以及关于祂的奥秘；然后再指出品行的改正，好让他们在认识主之后，能真诚地切望遵行祂所命令的事。因为当法律的向导不为人知时，人就不会轻易地去遵行它们。忠信的天主之仆\Person{梅瑟}就采用了这个方法；当他颁布神圣法律安排的话语时，他首先宣告了有关认识天主的事：\BibleQuote{以色列，你要听！上主我们的天主是唯一的上主。}\BibleRef{申 6:4}之后，他将天主向百姓隐约显示，并教导他们应当信靠谁，且使他们明白谁是真正的天主，然后他才颁布那使人能蒙祂喜悦的法律，说：\BibleQuote{不可奸淫；不可偷盗；}连同其他诫命。因为按照宗徒的教导：\BibleQuote{凡接近天主的人，必须信祂存在，且信祂赏报寻求祂的人}\BibleRef{希 11:6}。而如今，祂是藉着德行的行为而被寻求的，正如先知所说：\BibleQuote{趁上主可找到的时候，你们应寻找祂；趁祂在近处的时候，你们该呼求祂。罪人应离开自己的行径，恶人该抛弃自己的思念}\BibleRef{依 55:6}。\par

4. 若一个人不因\WorkTitle{牧者}的见证而感到反感，也是好的。该书在开头说：\Quote{在万有之前，你当相信有一位天主，他创造并建立了这一切，从无中召唤它们存在。}此外，真福圣史们——他们记载了主的话——在福音的开头，写了关于我们救主的事；这样，他们先使人认识主，造物主，好让人在叙述所发生的事时能够相信他们。因为他们若没有教导他是造物主，写道：\BibleQuote{在起初已有圣言？}\BibleRef{若 1:1}，当他们记载有关那位生来瞎眼的人，以及其他复明的瞎子、由死者中复活的人、水变成酒、癞病人得洁净等事时，人怎能相信他们呢？或者按玛窦所载，那位由达味的后裔所生者，是\Person{厄玛奴耳}，是永生天主之子，犹太人和亚略派转面不顾他，我们却承认并朝拜他。因此，宗徒按所宜的，写信给不同的人，但他特别提醒他自己的儿子，\Quote{不可轻看从他所得的教训}，并叮嘱他：\BibleQuote{要记念耶稣基督，他从死者中复活，是达味的后裔，就照我所传的福音。}在谈到这些事已交托给他，要常存记念时，他立即写信给他说：\BibleQuote{你要专心于此等事，全神贯注于其中}\BibleRef{弟前 4:15}。因为恒常的默想和铭记神圣的言语，能坚定对天主的虔敬，产生一种与他不分离、不流于形式的爱；正如他怀着这种心态，论及自己和有同样心态的人时，大胆地说：\BibleQuote{谁能使我们与天主之爱隔绝？}\BibleRef{罗 8:35}因为这样的人，在主内坚固，对他有不可动摇的志向，且同一精神（因为\BibleQuote{但那与圣神相结合的，便是与他成为一神}\BibleRef{格前 6:17}），确实像\Quote{如\Place{熙雍}山}。纵然有千般磨难冲击他们，却建立在磐石即基督之上。疏忽的人在基督内毫无喜乐；他们既没有常存的善意，便被一时的打击所玷污，重视现世之事过于一切，在信德上动摇不定，应受谴责。因为\BibleQuote{世俗的焦虑和财富的迷惑，把话窒息了}\BibleRef{玛 13:22}；或者，正如耶稣在那针对他们的比喻中所说的，他们没有坚固所宣讲给他们的信仰，只是暂时存留，一旦遇到迫害，或因圣言而遭遇磨难，立刻就跌倒了。至于那些图谋邪恶的人，我们说他们所想的不是真理，而是虚妄；不是正义，而是不义，因为他们的舌头学会了说谎。他们行了恶，却不知悔改。他们沉迷于恶行，一往直前，不知回头，甚至践踏了有关爱近人的诫命，不但不爱他们，反而图谋陷害他们，正如圣人作证说：\BibleQuote{那寻索我命的，开口讲论虚妄，终日策划诡计。}但这种思念的原因无非是缺乏训诲，正如神圣的箴言早已宣告：\Quote{那离弃父亲训诲的儿子，心里默思邪恶的话。}但这种思念既是邪恶的，圣神用这样的话谴责，又用别的话责备说：\BibleQuote{你们的手沾满了血，你们的手指染满了罪；你们的口唇说谎言，你们的舌头言不义：没有人讲公义，也没有人凭诚实判断}\BibleRef{依 59:3}。但他立即指出这种邪恶思念的结局说：\BibleQuote{他们依赖虚无，说虚话；他们怀孕毒害，产生邪恶。他们孵的是蝮蛇的蛋，结的是蜘蛛的网；谁吃了这蛋，必定要死；打碎了蛋，必发现毒蛇在内。}他还预先宣布这种人的希望何在：\BibleQuote{因为正义追不上他们，他们等待光明，却只见黑暗；期待光亮，却只在幽暗中行走。他们像瞎子摸索墙壁，像无眼之人摸索；他们中午跌倒，如在深夜；死了便悲叹。他们一同吼叫如熊，或如鸽子哀鸣。}\par

这就是邪恶的果实，这些报酬给予其党羽，因为邪恶不能拯救自己的人。其实，它反而攻击他们，首先撕裂他们，尤其毁灭他们。那些遭受这些的人有祸了；因为\BibleQuote{它比双刃剑更锐利}\BibleRef{希 4:12}，预先迅速斩杀那些将要抓住它的人。因为他们的舌头，正如圣咏作者的证言，是\Quote{利剑，他们的牙齿是枪箭}。但奇妙的是，那被人图谋陷害的人往往毫发无伤，而他们却被自己的枪刺透：因为在触及他人之前，他们自身内已怀着愤怒、怒气、恶意、诡诈、仇恨、刻薄。即使他们未能将这些施于他人，它们也会立即返回并攻击他们自己，正如他所祈求的：\Quote{愿他们的刀剑刺入他们自己的心}。还有这样一句箴言：\Quote{恶人被自己的罪恶之链紧系}。\par

5. 犹太人在他们的幻想中，以及他们同意对主行不义时，忘记了他们是在为自己招致愤怒。因此，\OriginalTerm{圣言}{Word}为他们哀叹，说：\BibleQuote{万邦为什么嚣张，众民为什么妄想？}因为犹太人的幻想确实是虚妄的，图谋害死生命，并构思荒谬之事反对\OriginalTerm{圣父}{Father}的圣言。谁看到他们的流散和城市的荒凉，谁不恰当地说：\BibleQuote{祸哉他们，因为他们心怀恶念，自言自语说，让我们捆绑义人，因为他不讨好我们。}我的弟兄们，这完全正确；因为当他们在\OriginalTerm{圣经}{Scriptures}上犯错时，他们不知道\BibleQuote{挖掘陷阱的必自陷其中，拆毁墙壁的必被蛇咬伤。}\BibleRef{训 10:8}如果他们没有转面背离主，他们就会畏惧神圣的\WorkTitle{圣咏}中先前所写的：\BibleQuote{外邦人陷入自己挖掘的陷阱，自己被自己所暗藏的罗网缠住了脚。上主在施行审判时彰显了自己：恶人被自己双手的作为捉住。}让他们观察此事，以及\BibleQuote{他们所不知的罗网必临到他们，他们暗藏的网必捉住他们。}但他们不明白这些事，如果明白，\BibleQuote{他们决不至于将光荣的主钉在十字架上}。\BibleRef{格前 2:8}\par

6. 因此，主的那义德而忠信的仆人，他们\BibleQuote{成为天国的\OriginalTerm{门徒}{disciples}，能从其中拿出新的和旧的东西；}并且他们\BibleQuote{默想上主的话，无论坐在家里，或躺卧或起身，或行在路上；}——因为他们由于\OriginalTerm{圣神}{Spirit}的应许而满怀希望，这圣神曾说，\BibleQuote{凡不随从恶人的计谋，不插足于罪人的道路，不参与讥讽者的席位，而专心致志于上主法律，并昼夜默想上主法律的人，才是福的；}——他们扎根于\OriginalTerm{信德}{faith}，喜乐于\OriginalTerm{望德}{hope}，热切于心神，便有胆量说，\BibleQuote{我的口要说智慧，我心的默想是明达的。}又说，\BibleQuote{我默念你的一切作为，我默想你手的一切化工。}以及，\BibleQuote{我若在床上记念了你，便在清晨默想你。}然后，在胆量上成长，他们说，\BibleQuote{我心的默想常在你面前。}这样的人结局如何？他立刻引用：\BibleQuote{上主是我的助佑和我的救赎主。}因为对于那些如此省察自己，并使自己的心合乎上主的人，不会有灾祸临到；因为实在，他们的心因信赖上主而坚固，正如经上所写：\BibleQuote{依赖上主的人有如熙雍山，住在耶路撒冷的永不动摇。}因为无论何时，那狡猾者若胆敢攻击他们，主要是为破坏\OriginalTerm{圣人}{saints}们的秩序，在弟兄中制造分裂；即使在此事上，主也与他们同在，不但作他们的复仇者，而且当他们已受打击时，作他们的拯救者。因为这是天主的应许：\BibleQuote{上主必替你们作战。}\BibleRef{出 14:14}从此，虽然外在的患难和试炼临到他们，但他们照着\OriginalTerm{宗徒的}{apostolic}话，并\BibleQuote{在患难中要坚忍，恒心祈祷}\BibleRef{罗 12:12}及默想法律，他们抵抗那些临到他们的事，蒙天主悦纳，并说出经上所写的话：\BibleQuote{我虽遭遇忧患和痛苦，但你的诫命仍是我的默想。}\par

7. 而人不仅在行动上，也在心思念虑中，趋向德行。因此他接着说：\BibleQuote{我的眼睛先于晨光，以默想你的言语。}因为身体完善的属灵默想，应当先于肉身的行动。我们的\OriginalTerm{救主}{Saviour}在教导这同一事时，不正也是从心思念虑开始吗？祂说：\BibleQuote{凡注视妇女，有意贪恋她的，已在心里犯了奸淫；}又：\BibleQuote{凡向自己弟兄发怒的，就该负杀人的罪责。}因为没有愤怒，就阻止了杀人；而先除去了贪恋，奸淫的指控就无从说起。因此，默想法律是必要的，我亲爱的，并要不断地守德，\BibleQuote{好使\OriginalTerm{圣人}{saint}完备，适于行各种善工。}\BibleRef{弟后 3:17}因为藉着这些事，有\OriginalTerm{永生}{eternal life}的应许，正如\Person{保禄}写给\Person{弟茂德}的，称恒常的默想为操练，并说：\BibleQuote{要在虔敬上操练自己；因为身体的操练益处不多，惟独虔敬在一切事上都有益处，因为有今生与永生的应许。}\BibleRef{弟前 4:7}\par

8. 我的弟兄们，那人的德行真值得钦佩！因为他藉着\Person{弟茂德}嘱咐众人，应当对虔敬一事最为关注，而在一切之上，将对天主的\OriginalTerm{信德}{faith}置于首位。因为不义的人有什么\OriginalTerm{恩宠}{grace}呢，即使他假装遵守诫命？不，不义的人连法律的一部分也不能遵守，因为他的心思怎样，他的行为也必怎样；正如\OriginalTerm{圣神}{Spirit}谴责这样的人，说：\BibleQuote{愚妄的人心中说：\InnerQuote{没有天主。}}之后，\OriginalTerm{圣言}{Word}表明行为与思想相应，说：\BibleQuote{他们都败坏了，在计谋中沦为卑污。}所以，不义的人在各方面败坏他的身体；偷盗、犯奸淫、诅咒、醉酒，并做类似的事。正如\Person{耶肋米亚}，先知，谴责以色列，呼喊说：\BibleQuote{但愿我在旷野有个旅客的寓所！好叫我离开我的人民，离他们远去：因为他们都是些淫乱之徒，是一群反叛，他们弯起舌头像弓一样；在地上得势的是谎言而非真理，他们恶上加恶；但他们却不认识我。}\BibleRef{耶 9:2}因此，为了邪恶和虚谎，并为了他们恶上加恶的行为，他谴责了他们的行径；但因他们不认识主，且无信德，他就指控他们为不义。\par

9. 因为信德与虔敬彼此相联，如同姊妹；凡信祂的人必有虔敬，而虔敬的人也越发相信。所以，身处邪恶之中的人，无疑也偏离了信德；从虔敬上跌倒的，也从真信德上跌倒了。例如，\Person{保禄}为证实这点，劝勉他的门徒说：\BibleQuote{你要躲避亵慢的言谈，因为这样的人必趋更不虔敬，他们的言论就像毒癌一样蔓延；其中有\Person{依默纳约}和\Person{非肋托}。}\BibleRef{弟后 2:16} 他说明他们的邪恶在于何处，说：\Quote{他们离开了信德，说复活已经是过去的事了。} 此外，为表明信德与虔敬密不可分，宗徒又说：\BibleQuote{凡愿意在基督耶稣内虔敬度日的人，都要遭受迫害。} 随后，为使人在迫害中不致放弃虔敬，他劝他们持守信德，补充说：\BibleQuote{所以，你要坚持你所学和所确信的事。} 正如兄弟相助，彼此成为对方的墙垣；同样，信德与虔敬一起增长，相互依存，练就了其中一项，必然因另一项而坚固。因此，宗徒愿门徒在虔敬上始终操练，并为信德而奋斗，便劝他们说：\BibleQuote{你要奋力打信德的好仗，争取永生。}\BibleRef{弟前 4:7} 因为人若先弃绝偶像的邪恶，并正确宣认那位真是天主的，接着他便凭着信德与那些攻击祂的人作战。\par

10. 我们所说这两样——信德与虔敬，所望相同，即永生；因为他说：\BibleQuote{你要奋力打信德的好仗，争取永生。} 又说：\BibleQuote{要在虔敬上操练自己，因为虔敬在各方面都有益处，有今生和来生的应许。}\BibleRef{弟前 4:7} 为此，那些如今已脱离教会的\OriginalTerm{亚略狂徒}{Ario-maniacs}，既是基督的敌对者，他们挖了一个不信的坑，自己却陷了进去；他们既在虔敬上败坏，就\BibleQuote{倾覆了心地纯朴者的信德，}\BibleRef{罗 16:18} 亵渎天主子，说祂是受造物，且是从虚无中而有。然而，正如当初宗徒已警告过人们要提防\Person{非肋托}和\Person{依默纳约}的追随者，如今也同样预先警告众人要提防他们那样的不虔，说：\BibleQuote{天主坚固的基础已立定，上面盖有印记：主认识属于祂的人；又说：凡呼求主名的人，应远离邪恶。}\BibleRef{弟后 2:19} 因为人若想正当地庆祝庆节，就应该远离邪恶与不义的行径；因为被恶人的玷污所沾染的人，不能向上主我们的天主祭献逾越节。因此，当时在埃及的百姓说：\BibleQuote{我们不能在埃及向上主我们的天主祭献逾越节。}\BibleRef{出 8:26} 因为至高天主愿意他们远远离开\Person{法郎}的臣仆和铁炉，好使他们脱离邪恶，并谨慎地摒弃一切异端思想，得以领受对天主的认识及善行。因为祂说：\BibleQuote{你们应远离他们，从他们中间出来，不要触摸不洁之物。}\BibleRef{格后 6:17} 人除非借着省察自己的行为，就无法离弃罪恶而获得善行；当人在虔敬上受过操练之后，他便会持守信德的宣认，保禄在打完了那场好仗之后，也获得了这宣认，即那为他存留的义德冠冕；公义的审判者必将赐给这冠冕，不独赐给他，也赐给所有像他一样的人。\par

11. 这样的默想与虔敬的操练，既是圣人们一贯的品行，如今当我们若怀有这种心态时，天主圣言也切愿我们与他们一同守这庆节。庆节是什么，岂不是对天主不断的敬拜，对虔敬的认识，以及发自全心的同心合意、恒常不懈的祈祷吗？因此，\Person{保禄}切愿我们常处于这种心态，便吩咐说：\BibleQuote{应常欢乐，不断祈祷，事事感谢。} 所以，我们不可各自分开，而要团结一致，共同守这庆节，正如先知所劝勉的：\BibleQuote{请大家前来，向上主欢呼，向拯救我们的天主歌舞。} 那么，有谁如此疏忽，或如此不听从这神圣的声音，竟不放下一切，奔赴这普遍且共同的庆节集会呢？这庆节并不仅限于一处，因为不仅一处地方单独守节；而是\BibleQuote{它们的声音传遍普世，它们的言语达于地极。} 祭献也不仅在一处举行，而是\BibleQuote{在万民中，到处有人向天主焚烧馨香，并奉献洁净的祭品。} 因此，当各地所有的人，同样将赞颂与祈祷奉献给仁慈良善的圣父时，当在各处的整个天主教会，欢欣喜乐地一同举行对天主的同一敬礼时，当众人都齐声高唱赞美之歌，并同声说\Quote{阿们}时，我的弟兄们，那将是何等蒙福！那时，谁不会参与其中，正当地祈祷呢？因为一切反对势力，尤其是\Place{耶里哥}的城墙，都要倒塌；圣神的恩赐那时将丰厚地倾注在所有人身上；每一个人感觉到圣神的降临，都要说：\BibleQuote{求祢使我们清晨就饱享祢的慈爱，使我们终生欢跃喜乐。}\par

12. 既然如此，让我们与圣人们一同向主欢呼，谁也不可在这些事上疏忽自己的本分；不要把\Person{尤西比乌斯}一党出于嫉妒而加于我们的磨难和试探放在心上，尤其是在此时。他们如今还想伤害我们，借控告置我们于死地，为的是我们所怀的虔敬，而我们将主作为这虔敬的助佑。但是，作为天主忠信的仆人，我们深知祂在我们遭难时是我们的救恩：因我们的主曾预先许诺说：\BibleQuote{几时人为了我而辱骂迫害你们，捏造一切坏话毁谤你们，你们是有福的。你们欢喜踊跃罢！因为你们在天上的赏报是丰厚的}\BibleRef{玛 5:11}。再者，救赎主亲口说过，磨难并非加于世上每一个人，而是只加给那些怀着圣洁敬畏之情对待祂的人：因此，敌人愈是围困我们，我们愈要自由；他们虽辱骂我们，我们仍要聚集；他们愈想使我们背离虔敬，我们便愈要大胆宣讲，说：\BibleQuote{这一切都临到了我们身上，我们却没有忘记祢}，也没有与\OriginalTerm{亚略疯狂者}{Ario-maniacs}同流合污，他们说祢是从虚无中存在的。那自永远与父同在的\OriginalTerm{圣言}{Word}，也出自父。\par

13. 所以我亲爱的弟兄们，我们要守这节，绝不可把它当作痛苦哀伤的时机来庆祝，也不可因虔敬而招致的现世磨难而涉足于异端者。但凡有助于喜乐欢欣之事，我们便当留心；免得我们心中忧愁如同\Person{加音}那样；而是像忠信良善的主之仆人一样，听到那句话：\BibleQuote{进入你主人的福乐罢！}\BibleRef{玛 25:21}。因为我们并没有设立哀伤忧愁之日，如有些人所认为的复活节那样；我们乃是满怀着喜乐欢欣地守这节。所以我们守这节，既不按照犹太人骗人的错谬，也不按照\OriginalTerm{亚略派}{Arians}的教训；他们将圣子从神性中除去，把祂列于受造之中；我们乃是仰望从主而来的正确道理。因为犹太人的诡诈和亚略派无边的邪恶，只引起忧伤的反思：前者起初杀害了主；后者却夺走了祂战胜死亡——犹太人把祂推入死亡——的地位，因为他们说祂不是创造者，而是受造物。因为祂若是受造物，便会被死亡束缚；但若祂没有被死亡束缚，依经上所载，祂就不是受造物，而是受造物的主，也是这不朽节日的主题。\par

14. 因为死亡的主，愿意废除死亡；祂既是主，所愿者即告完成；因为我们众人都已由死亡进入了生命。但犹太人和与他们同类者的幻想，却是虚空，因为结果并非他们所预想，反倒与他们为敌；\BibleQuote{在天上坐于高座者，必对他们嗤笑；上主必对他们加以耻笑。}因此，当我们的救主被带去受死时，祂阻止了跟随祂哭泣的妇女，说：\BibleQuote{不要为我哭泣}\BibleRef{路 23:28}；意思是要表明，主的死亡不是悲伤的事，而是喜乐的事，且那位为我们而死的主是活着的。因为祂的存在不是从虚无中得来，而是出于父。这真是一件喜乐之事：我们借着主的身体，能看见战胜死亡的标记，就是我们自己的不朽性。因为祂光荣复活了，显然我们众人的复活也必会发生；祂的身体既未腐朽，我们自己的不朽也就毫无疑问了。因为，如保禄所说（这是真理），就因一人的过犯，\BibleRef{罗 5:12}；照样，借着我们的主耶稣基督的复活，我们众人也都要复活。因为他说：\BibleQuote{这可朽坏的必须穿上不可朽坏的，这可死的总要穿上不可死的}\BibleRef{格前 15:53}。这在苦难时期实现了，那时我们的主为我们死了，因为\BibleQuote{我们的逾越节羔羊基督，已被祭杀了。}所以，祂既被祭杀，让我们每个人都以祂为食，并热切勤勉地分享祂的滋养；因为祂毫无吝啬地赐给了众人，且在每个人内成为\BibleQuote{成为涌到永生的水泉}\BibleRef{若 4:14}。\par

15. 我们从\OriginalTerm{帕梅诺特月}{Phamenoth}初九日开始四旬期的斋戒；\EditorialNote{马尔谷福音第5章}在那些日子里，我们以守斋侍奉上主，并先洁净了自己，之后我们从\OriginalTerm{帕尔穆提月}{Pharmuthi}十四日（4月9日）开始圣复活节。随后，将斋戒延伸到第七日，即该月十七日晚，我们深夜安息。主的光首先照耀我们，我们主复活的光荣神圣主日照耀我们时，我们当以源自善工的喜乐，在剩余七周内，直至五旬节，欢欣踊跃，并将光荣归于父，且说：\BibleQuote{这是上主所安排的一天，我们应该为此鼓舞喜欢}，借着我们的主和救主耶稣基督，愿光荣和权能，藉着祂，归于同一位的祂和祂的圣父，至于无穷之世。阿们。你们要以圣吻彼此致候。与我同在的众弟兄都问候你们。亲爱的弟兄们，我祈求你们在主内安康。\par

\Person{圣亚大纳修}的第十一封书信到此结束。\par

\SourceRecord{Translated by R. Payne-Smith.；Nicene and Post-Nicene Fathers, Second Series；Vol. 4.；Edited by Philip Schaff and Henry Wace.；Buffalo, NY: Christian Literature Publishing Co.,；1892.；<http://www.newadvent.org/fathers/2806011.htm>.}

% structure: p0001 source=h1 target=section reason=page-title

\section{第十三封书信}

\Emph{（公元341年）执政官\Person{马尔切利努斯}与\Person{普罗比努斯}，长官\Person{朗吉努斯}；第十四印狄克季；复活节，四月望前三日，法穆提月二十四日；戴克里先纪元五十七年。}\par

再次，我亲爱的弟兄们，我准备好向你们通报那救恩的庆节，这庆节将按每年惯例举行。因为尽管基督的敌对者以磨难和忧愁压迫你们和我们；然而，天主因我们彼此的信德\OriginalTerm{信德}{faith}安慰了我们，看，我甚至从罗马给你们写信。我在此与弟兄们一起守这庆节，也同你们在意愿和精神上一起守它，因为我们共同向天主献上祈祷：\BibleQuote{祂不但赐我们相信祂，也赐我们为祂受苦。} \BibleRef{斐 1:29} 因为我们因远离你们而忧愁，祂推动我们写信，好藉由书信彼此安慰，并互相激励行善。因为确实，众多磨难和针对教会的残酷迫害临到了我们。因为那些心思败坏、在信德上未受试炼的异端者\OriginalTerm{异端者}{heretics}，起来反对真理，猛烈地迫害教会，在弟兄们中，有的受鞭笞，有的被鞭打致伤，最严重的是，他们的侮辱甚至殃及主教。然而，不应因此就忽视这庆节。我们反而应特别记念它，绝不可随时忘记这纪念。那些不信者不认为有庆节的时节，因为他们一生沉溺于宴乐和愚昧；他们所守的节庆是哀伤的场合，而非欢乐。但对我们今生而言，这庆节首先是一条通往天堂的连续通道——这确实是我们的时节。因为这些事为锻炼和考验而服务，好使我们经考验后，显为基督热忱而被选召的仆人，得以与圣人们同为承继者。因为约伯如此说：\BibleQuote{整个世界是地上人类受试探的场所。} 然而，人在这世上通过磨难、劳苦和忧愁被试炼，为使各人能从天主领受相称的赏报，正如祂藉先知所说：\BibleQuote{我是天主，究察人心，考验人肺腑，照各人的行径赐与各人。} \BibleRef{耶 17:10}\par

2. 并非祂在人受考验时才初次知道人的事（因为在事发生之前祂已全然知道），而是因为祂是美善而爱人的，祂按各人的行为分赐相称的赏报，好使人人能高呼：\BibleQuote{天主的判断是正义的！} 如先知又说：\BibleQuote{主试探义人，洞察肺腑。} 再者，为此缘故祂考验我们每一个人，或是为使那不知德性的人，通过受考验者而显示出来，如同论及约伯所说的：\BibleQuote{你想我向你显现，岂有别故，不是为叫你显为义人吗？} 或是为使人们觉察己行时，能知道自己属于何种，从而或悔改其邪恶，或在信德上坚定不移。可赞颂的\Person{保禄}，当受磨难、迫害、饥渴困苦时，\BibleQuote{靠着那爱我们的主耶稣基督，我们在这一切事上，大获全胜。} \BibleRef{罗 8:37} 经过苦难，他身体固然软弱，但他藉着信德和\OriginalTerm{望德}{hope}，在心灵上变得强壮，而他的能力在软弱中得以圆满。\BibleRef{格后 12:9}\par

3. 其他圣人同样对天主怀有信赖，欢欣接受类似的考验，正如约伯说的：\BibleQuote{愿主的名受赞美！} \BibleRef{约 1:21} 圣咏作者却说：\BibleQuote{主，求祢察看我，试验我，熬炼我的肺腑和心肠。} 因为当力量被试验时，它就能说服愚顽人，他们认识到从天主之火中而来的洁净与益处，并不因这样的试炼而气馁，反而在其中喜乐，丝毫不因所发生的事受损，反而显得更加光辉，如同从火中炼净的金子；正如那受过这种纪律训练的人所说：\BibleQuote{祢已试验我的心，祢在夜间鉴察我；祢熬炼我，却未发现我有不义，我口未曾出恶言。} 但那些行为不受法律约束的人，他们只知吃喝与死亡，将考验视作危险。他们很快就绊倒，\BibleQuote{以致他们在信德上未经试炼，就陷入谬误的心智，行那些不合宜的事。} \BibleRef{罗 1:28} 因此，可赞颂的\Person{保禄}在敦促我们从事这类操练时，自己先以此衡量，说：\BibleQuote{所以我乐于忍受磨难、困苦。} 又说：\BibleQuote{要操练自己达到虔敬。} 因为他知道凡选择虔敬生活者必遭受迫害，他愿门徒们事先默想与虔敬相关的困难；好让试炼来临时，磨难兴起时，他们因已在其中受过操练，就能轻松忍受。因为一个人在心思上对于那些事习练了，通常会有一种隐藏的喜乐。可赞颂的殉道者\OriginalTerm{殉道者}{martyrs}们，起初就习练于困难，很快便在基督内达到完全，他们无视肉身的损害，而默观所期待的安息。\par

4. 然而，所有那些\Quote{以自己的名号称呼自己的地业}、并在思想中存有木、草、禾秸\BibleRef{格前3:12}的人；这些人，既然对艰难困苦感到陌生，就成了天国的外人。倘若他们知道，\Quote{磨难生坚忍，坚忍生经验，经验生望德，而望德不使人蒙羞}，他们必会效法\Person{保禄}的榜样操练自己，他说：\BibleQuote{我痛击己身，叫它服我，免得我给别人传报了福音，自己反而落选。}他们也会轻易忍受那不时临到他们身上为考验他们的苦难，如果他们听从了先知的劝诫\BibleRef{哀3:27}：\BibleQuote{人从青年时背负你的轭，原是好的；他当独坐无言，因为这是主加在他身上的；他当把脸转向打击他的人，饱受凌辱。因为主决不会永远抛弃；祂虽使人降卑，但必依照祂丰厚的慈爱施恩。}因为尽管这一切——鞭打、侮辱、凌辱——都来自仇敌，它们在天主丰厚的慈爱面前却毫无用处；我们必会迅速复原，因为这一切只是暂时的，而天主常施恩惠，将祂的慈爱倾注在悦乐祂的人身上。因此，我亲爱的弟兄们，我们不应顾念这暂时的事物，而应专注那永恒的事物。磨难虽会来临，终有止境；凌辱和迫害虽在，但与那摆在我们面前的望德相比，实在算不了什么。因为现今的一切与将来之事相比都微不足道；现时的苦楚，与将要显示给我们的望德，是不足比较的。有什么能与天国相比？有什么能与永生相比？我们在此所能给予的一切，怎能与我们将要在那边继承的相比？因为我们是\BibleQuote{天主的承继者，与基督同为承继者}\BibleRef{罗8:17}。因此，我亲爱的，我们不当顾念磨难和迫害，而当顾念那因迫害而为我们存留的望德。\par

5. 关于这一点，圣祖\Person{依撒加尔}的榜样可作说服，正如\WorkTitle{圣经}\BibleRef{创49:14}所说：\BibleQuote{依撒加尔喜爱美好的事物，安居在产业之间；他见安居之处美好，土地肥沃，便屈肩负重，成了农人。}他被神圣爱火所焚，如同\WorkTitle{雅歌}中的新娘，从\WorkTitle{圣经}中聚集丰盛，因为他的心神不为旧约所独占，而是被两种产业所吸引。于是，他好像展开双翼，远远望见那在天的\Quote{安息}，并且——既然这\Quote{土地}由如此美好的工程组成——那天上的国土更是如此；因为另一个[国度]常新，永不变老。因为这\Quote{土地}会逝去，正如主所说的；但那预备接待圣徒的却是永存不朽的。圣祖\Person{依撒加尔}看见这些事，便欢乐地在磨难和劳苦中夸胜，屈肩负重以便劳作。他不与打击他的人争斗，也不因凌辱而扰乱；反如勇士，藉着这些事越发得胜，越发热切地耕种他的土地，并从其中得利。\OriginalTerm{圣言}{Word}撒下种子，而他警醒地栽培，使之结实，甚至百倍。\par

6. 我亲爱的，这是什么意思呢？岂不是说，我们也当在仇敌列阵攻击我们时，以磨难为荣耀\BibleRef{罗5:3}，受迫害时不灰心，反而更要追求在基督耶稣我们的主内蒙召的冠冕？受凌辱时不扰乱，反而将面颊转给打击者，并屈下肩膀？因为贪爱享乐与热衷仇恨的人，如真福宗徒\Person{雅各伯}所说，\BibleQuote{各人为私欲所勾引、所饵诱}\BibleRef{雅1:14}。但我们既知道自己是为真理受苦，也知道否认主的人正在打击、迫害我们，就应照\Person{雅各伯}的话，\BibleQuote{我的弟兄们，几时你们落在各种试探里，要认为是大喜乐；因为知道你们的信德经过考验，才能生出坚忍}。弟兄们，让我们欢乐地守这庆节，因为我们知道，我们的救恩是在磨难的时候安排的。因为我们救主的救赎并非藉无所作为，而是藉祂为我们受苦，废除了死亡。关于此事，祂曾预先警戒我们说：\BibleQuote{在世界上你们要受苦难}\BibleRef{若16:33}。但祂并非对每一个人说这话，而是对那些殷勤忠信地服侍祂的人，祂预先知道，凡立志在祂内虔敬生活的人，都要受迫害。\par

7. \BibleQuote{作恶与行邪术的人，必日见败坏，既欺骗人，也受欺骗。}\BibleRef{弟后 3:13}因此，倘若这些无知的人，像那些自称能行征兆的解梦者和假先知一样，不是被酒灌醉，而是被自己的邪恶所醉，伪称司祭职，又仗着他们的恐吓而夸耀，你们不要相信他们；我们既受试探，就应谦卑自下，不要被他们引开。因为天主曾藉著梅瑟警告祂的子民说：\BibleQuote{如果你们中间兴起一位先知或一位作梦者，给你们显示一个神迹或奇事，而他所说的神迹和奇事实现了，以致他对你们说：\InnerQuote{让我们去随从事奉别的神，你们素不相识的神！}你们决不可听从这先知或这作梦者的话，因为上主你们的天主在试探你们，要知道你们是否全心全灵爱慕上主你们的天主。}\BibleRef{申 13:1}所以，当我们受这些事的试探时，我们决不离开对天主的爱。但是，我亲爱的，让我们现在遵守这庆节，不是把它当作受苦的日子，而是当作在基督内喜乐的日子，因为我们日日藉着祂获得滋养。我们要记念那位在逾越节期间被祭献的；因为我们庆祝这个，正是因为基督——我们的逾越节——已被祭献。\BibleRef{格前 5:7}那位曾将祂的子民从埃及领出，如今又废除了死亡，并废除了那握有死亡权势的——即是魔鬼——的，\BibleRef{希 2:14}现在也要同样使他蒙羞，并再次援助那些受苦受难、日夜向天主呼求的人。\BibleRef{路 18:7}\par

8. 我们从\OriginalTerm{法梅诺特月}{Phamenoth}十三日（3月9日）开始四十天斋期，并于\OriginalTerm{法尔穆提月}{Pharmuthi}十八日（4月13日）开始复活节圣周；在第七天，即二十三日（4月18日），安息休息，而同月法尔穆提月二十四日（4月19日）\OriginalTerm{大周}{great week}的第一天既已破晓，我们就从此日算至五旬节。我们要时时歌唱赞美，呼求基督，藉我们的主耶稣基督从仇敌手中获得解救；愿光荣与权能，藉着祂，归于圣父，及于无穷世之世。阿们。你们要以圣吻彼此问候。凡在此与我同在的人都问候你们。我祈求，我亲爱的弟兄们，愿你们在主内身体健康。\par

他也是在罗马写了这封信。第十三封信到此结束。\par

\SourceRecord{Translated by R. Payne-Smith.；Nicene and Post-Nicene Fathers, Second Series；Vol. 4.；Edited by Philip Schaff and Henry Wace.；Buffalo, NY: Christian Literature Publishing Co.,；1892.；<http://www.newadvent.org/fathers/2806013.htm>.}

% structure: p0001 source=h1 target=section reason=page-title

\section{第十四封书信}

\Emph{（为342年。）执政官：奥古斯都\Person{君士坦提乌斯三世}、\Person{君士坦斯二世}；长官：同前\Person{朗吉努斯}；第十五小期；复活节为四月十一日，即法尔穆提月十六日；戴克里先纪元五十八年。}\par

弟兄们，我们盛宴的喜乐常在我们身边，凡愿意庆祝这盛宴的人，它从不落空。因为\OriginalTerm{圣言}{Word}就在近处，祂为了我们成为一切，就是我们的主耶稣基督，祂曾许诺永远与我们同住，因此曾高声说：\BibleQuote{看，我同你们天天在一起，直到今世的终结}\BibleRef{玛 28:20}。因为祂既是牧者、大司祭、道路与门，为我们同时是一切，同样，祂又作为节期和圣日显示给我们，正如真福宗徒所说：\BibleQuote{我们的逾越节羔羊基督，已被祭杀作了牺牲}\BibleRef{格前 5:7}。祂就是那所期待的，祂在圣咏作者祈祷时使光闪耀，圣咏作者说：\Quote{我的喜乐，求你救我脱离围绕我的人；}这才是真正的欢跃，真正的节期，即从邪恶中得解救，人唯有彻底度正直的生活，心灵上全心顺服天主，才可抵达这一境界。因为圣人们终其一生，都如同在节期中欢跃的人。有人在向天主的祈祷中得享安息，如真福\Person{达味}，他夜里起来，不只一次而是七次。有人在赞美诗歌献上光荣，如伟大的\Person{梅瑟}，他为战胜\Person{法郎}和那些监工而唱了赞美歌。\EditorialNote{出 15}另有人以不懈的勤奋举行崇拜，如伟大的\Person{撒慕尔}和真福\Person{厄里亚}；他们已结束了尘世的征程，如今在天上守这节期，并从先前借预像所学得的事中欢跃，由这些预像认出了真理。\par

2. 但我们现在庆祝这节期，该使用什么洒水礼仪呢？我们急忙赶赴这节庆时，谁将做我们的向导？我亲爱的诸位，惟有你们将要与我同呼其名的那一位才能做到，就是我们的主耶稣基督，祂曾说：\Quote{我是道路。}因为正如真福\Person{若望}所说，祂\Quote{除免世罪}。祂净化我们的灵魂，正如\Person{耶肋米亚}先知在某处说：\BibleQuote{你们站在路上察看，访问古道，哪是善道，便行在其上；如此，你们心灵必得安息}\BibleRef{耶 6:16}。古时，\BibleQuote{洒在不洁之人身上的公羊和母牛的血，并母牛的灰，尚且能洁淨他们的肉身}\BibleRef{希 9:13}；但如今，藉着\OriginalTerm{天主圣言}{God the Word}的恩宠，人人都得以彻底净化。跟随祂，我们甚至在此世，就能如同站在天上\Place{耶路撒冷}的门槛上，预先默想那永恒的节期；正如真福宗徒们，一同跟随了身为他们领袖的救主，如今成为这同样恩宠的教师，他们说：\BibleQuote{看，我们舍弃了一切，而跟随了你}\BibleRef{谷 10:28}。因为跟随主，以及属于主的节期，并非只凭言语，而要凭行为来完成，每一律法的颁布和每一命令，都涉及明确的行动。就如伟大的\Person{梅瑟}传授神圣法律时，要求百姓作出许诺\BibleRef{出 19:8}，要按照法律行事，这样许诺之后，他们便不致忽视，并被指控为说谎者，同样，这最小的逾越庆典，并不引发疑问，也不要求回答；一旦发出话语，行动随之而来，因为祂说：\Quote{以色列子民应举行逾越节；}祂的命令意在使诫命即刻执行，而命令又辅助其实现。至于这些事，我信赖你们的智慧，以及你们对教导的关心。此类要点，我已屡次并在多封书信中提及。\par

3. 但现在，比一切都更为必要的，我愿提醒你们，并偕同你们提醒我自己，即诫命要求我们来到逾越节期，不可亵慢而无准备，而应行圣事与教义之礼，和规定的仪式，正如我们从历史记载所学到的：\Quote{外方人或买来的人，或是未受割损的人，都不可吃逾越节羔羊。}也不可在\Quote{任何}房屋内吃，祂吩咐要急速地吃；因为从前我们曾在\Person{法郎}的奴役和监工的命令下呻吟忧愁。因为古时以色列子民如此行时，他们被算为配领受预像，那预像原是为这节期而设立的；而如今，节期并非因预像而被引入。正如\OriginalTerm{天主圣言}{Word of God}切愿如此时，也对祂的门徒说：\BibleQuote{我渴望又渴望，在我受难以前，同你们吃这一次逾越节晚餐}\BibleRef{路 22:15}。那真是个奇妙的帐目，因为那时，人可以看见他们束着腰，仿佛为着游行或跳舞，拿着棍杖，穿着凉鞋，拿着无酵饼出去。这些从前发生在预像中的事，都是预表。但如今，真理已临近我们，\BibleQuote{祂是不可见的天主的肖像}\BibleRef{哥 1:15}，我们的主耶稣基督，真光，祂代替棍杖，是我们的权杖；代替无酵饼，是自天降下的食粮；代替凉鞋，给我们装备了福音的准备\BibleRef{弗 6:15}；简言之，祂借这一切引导我们到祂圣父那里。如果仇敌磨难我们、迫害我们，祂将再次代替\Person{梅瑟}，用更佳的话语鼓励我们，说：\Quote{你们放心，我已战胜了恶者。}而如果我们在渡过\Place{红海}后，炎热再次折磨我们，或我们遭遇苦水，主仍将再次显现给我们，将祂的甘甜和赐生命的泉源分施给我们，说：\Quote{谁若渴，到我这里来喝吧。}\par

4. 那么，我们为何还徘徊，为何迟延，不带着全部热忱和勤奋前来赴此庆节，相信是\Person{耶稣}在召叫我们呢？祂为我们成为万有，以无数方式担负我们的救恩；祂为我们又饥又渴，却以自己的救恩恩赐赐我们饮食。因为这是祂的光荣，是祂天主性的奇迹：祂将我们的苦难转变为祂的喜乐。因为，祂本是生命，却为我们而死，使我们活过来；祂本是\OriginalTerm{圣言}{Word}，却成了血肉，为以圣言教导血肉；祂本是生命之泉，却感受我们的干渴，好催促我们赴此庆节，说：\BibleQuote{人若渴了，可以到我这里来喝。}\BibleRef{若 7:37}那时，\Person{梅瑟}宣告这庆节的开端，说：\BibleQuote{本月为你们是月份之首。}\BibleRef{出 12:2}但主在末世降世\BibleRef{希 9:26}，却宣告了另一日，并非要废除法律，绝非如此，而是要坚固法律，并成为法律的终向。\BibleQuote{因为基督就是法律的终向，使凡信祂的人得称为义；}正如圣\Person{保禄}所说：\BibleQuote{我们因信德就废了法律吗？绝对不是！我们反而坚固了法律。}这些事甚至惊动了\Person{犹太人}派来的差役，以致他们惊讶地对\Person{法利塞人}说：\BibleQuote{从来没有像祂这样说话的！}\BibleRef{若 7:46}那么，是什么惊动了那些差役？或是什么如此触动了这些人，使他们称奇？无非是我们救主的胆识与权威。因为古时先知和经师研读圣经时，他们察觉所读的不指自己，乃指他人。例如\Person{梅瑟}：\BibleQuote{主天主将从你们弟兄中给你们兴起一位先知像我，你们要听从祂的一切吩咐。}\Person{依撒意亚}又说：\BibleQuote{看，一位贞女将怀孕生子，人称祂为\Person{厄玛奴耳}。}还有其他先知以不同方式预言了主的事。但主却以自己、而非他人为这些预言的指向；祂将一切归于自己，说：\BibleQuote{人若渴了，可以到我这里来。}\BibleRef{若 7:37}——不是到别人那里，而是到\Quote{我}这里。人固然可以从他人那里听说我的来临，但他从此不可再从他人那里饮水，而要从我这里。\par

5. 因此，我们也应当这样：当我们来赴此庆节时，不再像赴旧约的影子，因为它们已经实现；也不像赴普通的庆节，而要快跑奔向主，祂自己就是这庆节\OriginalTerm{庆节}{feast}，不视之为放纵与口腹之乐，而视之为德行的彰显。因为外邦人的庆节充满贪婪和全然怠惰，因为他们以为闲暇便是过节；他们过节时，却行丧亡之事。但我们的庆节在于修炼德行和实行节制；正如先知的话在某处作证说：\BibleQuote{四月斋戒、五月斋戒、七月斋戒、十月斋戒，必为犹大家成为欢欣喜乐的日子和愉快的庆节。}\BibleRef{匝 8:19}既然这修炼的时机已陈明在我们面前，这样的日子已经来临，先知的声音已发出庆祝庆节的宣告，我们就当全心全力响应这美好的宣告，像赛跑场上竞技的人，彼此争先恪守斋戒的洁净\BibleRef{格前 9:24}，藉着祈祷的警醒、研读圣经、施舍穷人，并与仇敌和好。让我们聚拢四散的人，摒除骄傲，回归心谦，与众人和睦，并勉励弟兄们相亲相爱。圣\Person{保禄}也曾时常斋戒警醒，甚至甘愿为弟兄受诅咒。圣\Person{达味}也因斋戒自谦，而放胆说：\BibleQuote{主，我的天主，我若行了这事，我手若有罪过，我若以恶报善，就让我的仇敌追赶我，将我践踏。}我们若这样行，就必战胜死亡，并领受天国的凭据。\par

6. 我们在法尔穆提月十日（4月5日）开始神圣的复活节庆，在同月的法尔穆提月十五日（4月10日）第七天的晚上，停止神圣的斋戒。我们当在同月的法尔穆提月十六日（4月11日）庆祝神圣的庆节；并逐日持续直到神圣的五旬节，经过这一连串的庆节，让我们为圣神举行庆典，祂就在我们身边，在\Person{耶稣基督}内，藉着祂并同祂，愿光荣和权能归于圣父，至于无穷之世。阿们。\par

\EditorialNote{第十五封和第十六封信件缺失。}\par

\SourceRecord{Translated by R. Payne-Smith.；Nicene and Post-Nicene Fathers, Second Series；Vol. 4.；Edited by Philip Schaff and Henry Wace.；Buffalo, NY: Christian Literature Publishing Co.,；1892.；<http://www.newadvent.org/fathers/2806014.htm>.}

% structure: p0001 source=h1 target=section reason=page-title

\section{第十七封信}

\Emph{（为345年。）执政官\Person{亚曼提乌斯}、\Person{阿尔比努斯}；\Place{加沙}长官\Person{聂斯托利}；税收周期三；复活节日，4月7日，法尔穆提月12日；月龄十九；戴克里先纪元61年。}\par

\Person{亚大纳修}致\Place{亚历山大里亚}的司铎与执事，并亲爱的弟兄们，在基督内致候。\par

依照惯例，亲爱的诸位，我将复活节的通知送达你们，以便你们也可照常将同样的事通知给远方各区的弟兄们。因此，在本节日——我指的是法尔穆提月二十日的这个节日——之后，随后的复活节将在4月7日，或依\Place{亚历山大里亚}人的算法，在法尔穆提月十二日。所以你们要在所有那些地区通知，复活节将在4月7日，或依\Place{亚历山大里亚}的算法，在法尔穆提月十二日。愿你们在基督内平安，我祈求，亲爱的弟兄们。\par

\SourceRecord{Translated by R. Payne-Smith.；Nicene and Post-Nicene Fathers, Second Series；Vol. 4.；Edited by Philip Schaff and Henry Wace.；Buffalo, NY: Christian Literature Publishing Co.,；1892.；<http://www.newadvent.org/fathers/2806017.htm>.}

% structure: p0001 source=h1 target=section reason=page-title

\section{第十八封信}

\Emph{（为346年。）执政官：奥古斯都\Person{君士坦提乌斯}四世、\Person{君士坦斯}三世；总督：同一位\Person{聂斯托利}；小纪4；复活节日：罗马历四月朔前三日，法尔穆提月初四；月龄二十一；戴克里先纪元六十二年。}\par

\Person{亚大纳修}致\Place{亚历山大里亚}的司铎和执事，在主内亲爱的弟兄们，问候。\par

你们做得很好，最亲爱的弟兄们，你们已在那些区域发布了关于复活节的惯常通知；因为我已见到并认可了你们的严谨。通过其它信件，我也曾通知你们，以便当今年结束时，你们得以知晓关于下一年的事。然而现在，我仍认为有必要写同样的事，好使你们在确切知晓后，也能谨慎地致信。因此，在此即将结束的庆节之后——即法尔穆提月十二日，亦即罗马历四月七日——复活节日将是罗马历四月朔前第三日，即\Place{亚历山大里亚}人的法尔穆提月初四。所以，当庆节结束后，依照古老习惯，再次在这些区域如此通知：复活主日在罗马历四月朔前第三日，按\Place{亚历山大里亚}人的历法，即法尔穆提月初四。任何人都不要对此日子犹豫，也不要争辩说，复活节应在法梅诺特月二十七日举行；因为这件事在\OriginalTerm{神圣的公会议}{the holy Synod}上讨论过，并且全体在那里确定它为四月朔前第三日。所以我说，它是在法尔穆提月初四；因为此前的一周太早了。因此，不应有任何争论，而要采取合宜的行动。因为我也同样致信给了罗马人。所以，一如已通知你们的，发布通告：复活节在罗马历四月朔前第三日，按\Place{亚历山大里亚}人的历法，即法尔穆提月初四。\par

我祈求你们在主内得健康，我最亲爱的弟兄们。\par

\SourceRecord{Translated by R. Payne-Smith.；Nicene and Post-Nicene Fathers, Second Series；Vol. 4.；Edited by Philip Schaff and Henry Wace.；Buffalo, NY: Christian Literature Publishing Co.,；1892.；<http://www.newadvent.org/fathers/2806018.htm>.}

% structure: p0001 source=h1 target=section reason=page-title

\section{书信十九}

\Emph{（主后347年）执政官 \Person{鲁菲努斯} 与 \Person{欧瑟比乌斯}；总督仍为 \Person{奈斯托留斯}；课税周期第五年；复活节日在四月十二日，即发尔穆提月十七日；戴克里先纪元六十三年；月龄十五}\par

\BibleQuote{愿我们的主耶稣基督的天主和父受赞美\BibleRef{弗 1:3}}，这样的开头正适合一封信函，尤其在此刻，它以上主的名带来感恩，正如宗徒所说的，因祂从远方把我们带来，并恩许我们可以再次公开地，按惯例，给你们寄出复活节书信。因为这是节日的时期，弟兄们，且已临近了；现在不再像史书所载那样，用号角宣告，而是由为我们受苦并复活的救主，使它为我们所知并临近我们，正如 \Person{保禄} 所宣讲的：\BibleQuote{我们的逾越节羔羊基督，已被祭献了。}因此，从现在起，逾越节是我们的，不是外方人的，也不再是犹太人的。因为预像的时期已废除，旧事已过，现今，初春之月就在眼前，在此月中，人人都应谨守节期，顺从那位曾说过：\Quote{应遵守阿彼布月，为上主你的天主举行逾越节}。甚至外邦人也自以为在过节，犹太人也伪善地假装过节。但祂谴责外邦人的节庆，视之如哀悼者的饼\BibleRef{欧 9:4}，祂也转面不顾犹太人的节庆，因他们是被弃绝的人，说：\BibleQuote{你们的月朔和安息日，我的心灵憎恶\BibleRef{依 1:14}}。\par

2. 因为不按法律和虔敬所行的事，并没有益处，纵然它们看似如此，反而更暴露出那些胆敢行之者的伪善。因此，尽管这些人伪称献祭，但他们却听见父说：\Quote{你们的全燔祭不再悦纳，你们的祭献也不中悦我；即使你们奉献细面粉，也是徒然，乳香于我亦是可憎。}因为天主什么都不需要；既然对祂来说没有不洁之物，祂已厌饫了这些，正如祂藉 \BibleBook{依撒意亚}{先知书} 所证实的，说：\BibleQuote{我厌饫了\BibleRef{依 1:11}}。关于这些事，曾颁布了一项法律，为教导民众，并预表将来的事，因为 \Person{保禄} 在 \BibleBook{迦拉达}{书} 中说：\BibleQuote{在信德来临以前，我们都被看守在法律的监禁下，被禁锢着，直到那将要来临的信德启示出来；因此，法律成了我们的启蒙师，领我们归于基督，好使我们因信德而成义}。但犹太人既不知道，也不明白，因此他们在白日里如在黑暗中行走，摸索却未能触及我们所拥有的、包含在法律中的真理；他们拘泥于字句，却不顺从于精神。当 \Person{梅瑟} 蒙着帕子时，他们看他，但当他揭开帕子时，他们却转面不顾。因为他们不明白所读的，却错误地用一事代替另一事。因此，先知对他们呼喊说：\BibleQuote{虚妄与不忠在你们中间盛行}。主也关于他们这样说：\BibleQuote{异民之子欺骗了我，异民之子已衰老}。但祂多么温和地谴责他们，说：\BibleQuote{如果你们相信 \Person{梅瑟}，必会相信我，因为他指着我而写过\BibleRef{若 5:46}}。但他们既无信心，便继续错待法律，随自己的喜好说话，却不理解圣经；此外，他们伪善地假装遵守圣经的字句，并以此自恃，祂就对他们发怒，藉 \BibleBook{依撒意亚}{先知书} 说：\BibleQuote{谁向你们要求这些东西\BibleRef{依 1:12}}？又藉 \BibleBook{耶肋米亚}{先知书}，因他们非常大胆，祂威吓说：\BibleQuote{你们积聚全燔祭和祭牲，吃肉罢！因为在领你们出离埃及地的时候，我并没有同你们的祖先谈论或吩咐过全燔祭和祭牲的事\BibleRef{耶 7:21}}。因为他们行事不正直，他们的热心也不符合法律，反而在这样的日子里寻求自己的享乐，正如先知所指责的，他们殴打奴仆，聚集起来争闹斗殴，用拳头击打卑微的人，所做的一切都只为满足自己。因此，他们直至终末仍在节庆之外，尽管他们现在不合时宜、不合季节地炫示吃肉。因为，代替法律指定的羔羊，他们学会向巴耳献祭；代替真正的无酵饼，\BibleQuote{他们拾柴，父亲们点火，妇女们揉面，为要向天上的万象做糕饼，并向异神浇奠祭，这是要惹我发怒——上主说}。他们得到这些做法的公正报应，因为尽管他们假装遵守逾越节，欢乐和愉快的声音却被从他们口中夺去，正如 \BibleBook{耶肋米亚}{先知书} 所说：\BibleQuote{在犹大的城邑和耶路撒冷的街道上，已听不到欢乐的声音和愉快的声音，听不到新郎和新娘的声音}。因此，现今，\BibleQuote{他们中宰杀牛犊的，等于击杀一个人；祭杀羔羊的，等于勒死一条狗；奉献细面粉的，等于献猪血；奉上乳香作为记念的，等于亵渎者\BibleRef{依 66:3}}。这些事永远不会悦乐天主，上主的话也从未这样要求他们。但祂说：\BibleQuote{这些人选择了自己的道路，他们灵魂所喜悦的乃是那些可憎之事}。\par

3. 弟兄们，这是什么意思呢？我们有权利探究先知的话，尤其因为有些\OriginalTerm{异端者}{heretics}心思背离了法律。当初天主借着\Person{梅瑟}吩咐了有关祭献的诫命，而整本称为\WorkTitle{肋未纪}的书完全在安排这些事，为叫祂悦纳献祭的人。因此借着众先知，祂斥责那些轻视这些事的人，因为他们违背了那诫命，说：\BibleQuote{我并没有向你们要求过这些事。关于祭献，我也没有同你们的祖先说过，关于\OriginalTerm{全燔祭}{whole burnt-offerings}，我也没有命令过他们。}现在有人以为圣经彼此不合，或者以为那颁布诫命的天主是说谎的。然而没有任何不合，绝对没有，那真理的圣父也不能说谎；\BibleQuote{因为天主决不能说谎}\BibleRef{希 6:18}，正如\Person{保禄}所肯定的。但这一切对正确思考的人，及对所有凭信德接受法律书卷的人，都是显而易见的。依我看来——愿天主因你们的祈祷，恩准我冒昧发表的言论不离真理太远——起初并不是先有关于祭献的诫命和法律，那立法的天主的心意也不在于全燔祭，而是在于那些由它们所指示和预表的事物。\BibleQuote{因为法律只含有未来美物的\OriginalTerm{影子}{shadow}。}又说：\BibleQuote{这一切都规定下来，直到\OriginalTerm{革新时期}{time of reformation}。}\par

4. 因此，整个法律并非专论祭献，虽然法律中有关于祭献的诫命，目的是要借祭献开始教导人，使他们远离\OriginalTerm{偶像}{idols}，亲近天主，为当时而教导。所以起初当天主带领百姓出埃及时，祂并没有吩咐关于祭献或全燔祭的事，甚至他们来到\Place{西乃山}时也没有。因为天主不像人，不会事先操心这些事；祂颁赐诫命，是为叫他们认识那真天主和祂的圣言，并且鄙视那些虚假称为神、实非神而只有外表的东西。祂借着领他们出埃及，使他们渡过\Place{红海}，将自己显示给他们。但等到他们情愿事奉\Person{巴耳}，竟敢向那些虚无之物献祭，忘掉在埃及为他们所行的奇迹，想再回那里去的时候；那时确实在颁赐法律之后，那有关祭献的诫命才被立为法律；为叫他们那曾思念虚无之物的心思，转归那真天主，并学习首先不是献祭，而是要远离偶像，遵从天主的命令。因为当祂说：\BibleQuote{关于祭献我并没有说过话，关于全燔祭我也没有出过命，}立刻又加上：\BibleQuote{然而这是我命令他们的：你们应听从我的声音，这样我必作你们的天主，你们作我的子民；你们应走我吩咐你们的一切道路。}\BibleRef{耶 7:22} 这样，他们事先受了教导，便学会只事奉主。他们明白那影子应当存留到何时，也不忘记那即将来到的时期，在那时不再有牛犊、绵羊、公山羊作为献给天主的祭品\BibleRef{出 12:5}，这一切都应以纯粹属神的方式，借着恒常的祈祷、正直的生活和虔敬的言语而完成；正如\Person{达味}歌咏说：\BibleQuote{愿我的默想蒙祂悦纳。愿我的祈祷如馨香上升于祢面前，我举起的双手如晚祭}。在他内的圣神也命令说：\BibleQuote{你们应向天主献上颂谢祭，又应向至高者还你的誓愿。应献上正义的祭献，并要倚靠上主}。\par

5. 伟大的\Person{撒慕尔}也同样清楚地斥责\Person{撒乌耳}说：\BibleQuote{言语岂不胜过礼物吗？}\BibleRef{德 18:17} 因为借此，人便满全了法律，中悦天主，就如祂说：\BibleQuote{颂谢祭的光荣在于我}。人应\BibleQuote{学习：\InnerQuote{我喜欢仁爱胜过祭献}这句话的意思}，我就不谴责对手。但这事使他们厌倦，因为他们不热心去了解，\BibleQuote{如果他们认识了，决不至于将\OriginalTerm{荣耀的主}{Lord of glory}钉在十字架上。}\BibleRef{格前 2:8} 他们的结局怎样，先知曾预言说：\BibleQuote{祸哉，他们的灵魂！因为他们想出了恶念，说：我们要捆绑义人，因为他使我们厌恶}。这种自暴自弃的结局只能是错误，正如主斥责他们说：\BibleQuote{你们错了，因为不明白圣经}\BibleRef{玛 22:29}。后来当他们受斥责，本应醒悟过来，反而傲慢无礼地说：\BibleQuote{我们是\Person{梅瑟}的门徒；我们知道天主曾同\Person{梅瑟}讲过话}\BibleRef{若 9:28}，他们这样说更显得虚伪，并自证己罪。因为他们如果信了他们所听从的那一位，当主亲身临在，借\Person{梅瑟}说话时，就不会否认祂。\WorkTitle{宗徒大事录}中的太监却不是这样，当他听到\BibleQuote{你明白你所诵读的吗？}\BibleRef{宗 8:30} 时，并不耻于承认自己的无知，而是恳求教导。因此，他既然成了学习者，就蒙受了圣神的恩宠。至于那些执迷不悟的\OriginalTerm{犹太人}{Jews}，如\WorkTitle{箴言}所说：\BibleQuote{死亡临到他们身上。因为愚昧人死于自己的罪恶中}。\par

6. 异端者也是如此；他们既丧失了真正的辨别力，便胆敢为自己捏造 \OriginalTerm{无神论}{atheism}。\BibleQuote{愚人在心中说：没有天主。他们都已败坏，行径可憎。} 凡在心思上这样愚昧的人，其行为也是邪恶的，正如祂所说的：\BibleQuote{你们既是恶人，怎能说出善来？}\BibleRef{玛 12:34} 因为他们心思邪恶，所以行为邪恶。心思本就充满欺诈的人，怎能行公义呢？事先就准备好憎恨的人，又怎能爱呢？一心想贪爱钱财的人，怎能慈悲呢？注视妇女而动淫念的人，又怎能贞洁呢？\BibleQuote{因为由心里发出来的，是恶念、邪淫、奸淫、凶杀。} 愚人就被这些情欲所害，如同被海浪抛掷，受肉身的欢乐引诱牵引；因为经上记着说：\BibleQuote{愚人的肉身，大受颠簸。} 他与愚昧为伍时，就在暴风中颠簸，终致丧亡，正如\Person{撒罗满}在\BibleBook{}{箴言}中所说的：\BibleQuote{愚昧人与无知识的人，必一同丧亡，他们的财富留给外人。} 他们之所以遭遇这些事，是因为他们中间没有一个心灵健全的人来引导他们。因为哪里有明智，哪里就有圣言——灵魂的舵手——与船同在；\BibleQuote{有智识的，必获得指导；} 而没有指导的人，就如叶子般飘落。有谁像\Person{依默纳约}和\Person{菲肋托}那样完全背弃了呢？他们对复活持邪恶的见解，并因此对复活的信德遭遇了船破之患。叛徒\Person{犹达斯}离弃了舵手，便与犹太人一同丧亡。但门徒们因为有智慧，所以留在主身边；虽然海上波涛汹涌，船被波浪掩盖，因为起了风暴，又是逆风，他们却没有离开。因为他们唤醒了与他们同船的圣言，立时海就在主的一声命令下平静了，他们便得了救。他们同时成为宣讲者和教师，传述我们救主的奇迹，也教导我们效法他们的榜样。这些事是为我们而记载的，也是为了我们的益处，好使我们借着这些标记，认识那位行了这些奇事的主。\par

7. 因此，让我们怀着门徒们的信德，时常与我们的师傅交谈。弟兄们，因为世界对我们来说就像海洋，经上记着说：\BibleQuote{这是广阔无边的大海，船舶往来其中；有你所造的勒维雅堂，在其中游戏。} 我们在这海上漂浮，如同藉着风，凭我们自己的自由意志航行；每个人按照自己的意愿掌舵：要么在圣言的引领下进入安息，要么被享乐抓住，遭遇船破之患，在风暴中陷入危险。因为在海洋中有风暴和波浪，同样在世界上也有许多苦难和考验。因此，不信者\BibleQuote{一遭受患难或迫害，立刻就跌倒了，}\BibleRef{谷 4:17} 正如主所说的。因为他们没有在信德上稳固，又只注目于暂时的事物，就无法抵抗由患难而来的困难。但如同愚人在沙土上盖的房子，同样，没有理智的人\BibleRef{路 6:49} 在试探的冲击下，就像被风吹袭，便倒塌了。然而圣人呢，他们的官能因着节制获得了训练\BibleRef{希 5:14}，在信德上坚固，并且领悟了圣言，在考验中并不胆怯；虽然有时会有更严厉的考验临到他们，他们却仍保持忠信，并唤醒与他们同在的主，便得到解救。于是，他们经过水火，获得了安宁，并按规矩守节，向那救赎了他们的天主献上祈祷，怀着感恩之情。因为他们或是像\Person{亚巴郎}一样，因受试探而显为可靠；或是像\Person{约伯}一样，因受苦而蒙悦纳；或是像\Person{若瑟}一样，虽受压迫和欺诈，却耐心忍受；或是虽受迫害，却不被打倒；反而如经上所记，靠着天主，他们\BibleQuote{跳过墙垣} 正是那使弟兄分裂隔绝、使人离开真理的邪恶之墙。就这样，有福的\Person{保禄}，以软弱、凌辱、艰难、迫害、困苦为喜乐，为了基督而欢欣，并愿我们众人也都喜乐，他说：\BibleQuote{应常欢乐，事事感谢。}\BibleRef{得前 5:18}\par

8. 因为还有什么比远离邪恶、保持纯洁的品行、不断怀着感恩向天主献上祈祷，更适合这节日的呢？所以，弟兄们，让我们期待着庆祝天堂里永恒的喜乐，也在此守好这地上的节期：常喜乐，不断祈祷，事事感谢主。我感谢天主，不仅为祂所行的其它奇事，也为如今赐给我们的各种帮助：虽然祂严厉地惩罚了我们，却没有将我们交付于死亡，而是从远方，如同从地极，将我们领回，使我们与你们重新团聚。我在守节之时，也想到要预先通知你们复活节的大节日，好使我们可以一同上去，就像前往\Place{耶路撒冷}，同吃逾越节羔羊，不是分开，而像在同一个家中；我们不可像水煮的一样，用水稀释天主的圣言；也不可像折断骨头一样，破坏福音的命令。却要像用火烤过的，带着苦菜，心神热切，以斋戒和守夜，甚至席地而卧，怀着补赎和感恩之情来守这节。\par

9. 我们在法麦努特月第六日\EditorialNote{谷 2}开始守四旬期；我们若以斋戒和祈祷善度这段时期，就能抵达那圣日。因为凡疏忽不守四旬期的人，就如冒失而不洁净地践踏圣物一般，不能庆祝复活节。此外，让我们互相提醒，彼此勉励，不要疏忽，尤其要在那些日子守斋，好让斋戒接续不断，使我们能正确地结束这节日。\par

10. 如前所述，四十天的斋期始于法梅诺特月第六日（3月2日），苦难圣周始于法尔穆提月第十一日（4月6日）。我们应在此月的第十六日（4月11日）第七天深夜，结束斋期。当同一法尔穆提月的第十七日（4月12日）一周的第一日破晓时，我们应欢庆节日。然后，我们要一个接一个地加上五旬节的七个圣周，欢欣赞美天主，因为祂藉这些事提前向我们昭示了喜乐与永恒的安息，这安息已在天堂为我们及那些真正信仰我们的主\Person{耶稣基督}的人预备好了；愿光荣和权能，藉着祂，偕同祂，归于圣父，及圣神，世世无穷。阿们。\par

你们要以圣吻彼此问候。与我同在的弟兄们问候你们。\par

我想也有必要将任命主教之事告知你们，这些任命是为接替我们蒙福的已故同工，好使你们知道该向谁写信，又该从谁那里收信。因此，在\Place{西耶内}，\Person{尼拉蒙}接替同名的\Person{尼拉蒙}；在\Place{拉托波利斯}，\Person{马西斯}接替\Person{阿摩尼乌斯}；在\Place{科普托斯}，\Person{普塞诺西里斯}接替\Person{狄奥多鲁斯}；在\Place{帕诺波利斯}，因\Person{阿泰米多鲁斯}年老体弱，本人愿意，任命\Person{亚略}为辅理主教；在\Place{希普塞莱}，\Person{阿尔塞尼乌斯}因与教会和好而被任命；在\Place{吕科波利斯}，\Person{尤德蒙}接替\Person{普鲁西亚努斯}；在\Place{安提诺波利斯}，\Person{阿里翁}接替\Person{阿摩尼乌斯}和\Person{提兰努斯}；在\Place{奥克西林库斯}，\Person{狄奥多鲁斯}接替\Person{贝拉基乌斯}；在\Place{尼罗波利斯}，接替\Person{特翁}的是彼此和好的\Person{阿马图斯}和\Person{依撒格}；在\Place{阿尔西诺伊特}，\Person{安德肋}接替\Person{西尔瓦努斯}；在\Place{普罗索皮提斯}，\Person{特里亚德尔夫斯}接替\Person{塞拉潘蒙}；在\Place{迪奥斯法库斯}（河边），\Person{狄奥多鲁斯}接替\Person{塞拉潘蒙}；在\Place{萨伊斯}，\Person{帕弗努提乌斯}接替\Person{内梅西翁}；在\Place{克索伊斯}，\Person{狄奥多鲁斯}接替\Person{阿努比翁}；与他同在的还有\Person{伊西多鲁斯}，他已与教会和好；在\Place{塞斯罗伊特}，\Person{奥里翁}接替\Person{波塔蒙}；在\Place{克利斯马}，\Person{提托纳斯}接替\Person{雅各伯}；与他同在的还有\Person{保禄}，他已与教会和好。\par

\SourceRecord{Translated by R. Payne-Smith.；Nicene and Post-Nicene Fathers, Second Series；Vol. 4.；Edited by Philip Schaff and Henry Wace.；Buffalo, NY: Christian Literature Publishing Co.,；1892.；<http://www.newadvent.org/fathers/2806019.htm>.}

% structure: p0001 source=h1 target=section reason=page-title

\section{书信二十}

\Emph{（写于348年）执政官：\Person{斐理伯}与\Person{萨利亚}；总督仍为\Person{聂斯托利}；第六纪；复活节在4月3日，即法尔穆提月8日；戴克里先纪元64年；月龄18。}\par

弟兄们，让我们现在守这庆节吧，因为正如我们的主那时预先通知了祂的门徒，如今祂也预先告诉我们：\BibleQuote{过几天就是逾越节}\BibleRef{玛 26:2}，在这逾越节中，犹太人固然出卖了主，但我们却庆祝祂的死如同庆节，欢欣喜乐，因为我们那时脱离了我们的苦难，获得了安息。我们勤于聚集在一起，因为从前我们离散迷失，如今已被寻回；我们曾远离，如今被领近；我们曾是外人，如今成为属于祂的人——祂为我们受苦，被钉在十字架上，祂背负了我们的罪，正如先知所说的\BibleRef{依 53:4}，祂为我们受难受苦，为要从我们众人身上除去悲伤、哀愁和叹息。当我们口渴时，祂在节日本身就满足我们；祂站着大声说：\BibleQuote{谁若渴，到我这里来喝吧！}\BibleRef{若 7:37}。因为这就是圣徒们时刻持守的\OriginalTerm{爱}{love}：他们绝不间断，而是向主献上不停息的、恒常的祭献，并且持续地渴慕，向祂祈求畅饮；正如\Person{达味}歌唱说：\BibleQuote{我的天主，我的天主，我清晨寻求你；我的灵魂渴慕你；在这干旱、无路、无水之地，我的心和肉身多次渴想你。在圣所里，我这样瞻仰了你。}\Person{依撒意亚}先知也说：\BibleQuote{夜间我的灵魂寻求你，我内在的心神在清晨寻求你，因为你的诫命是光。}\BibleRef{依 26:9} 又有另一位说：\BibleQuote{我的灵魂因时时刻刻切慕你的判断而昏晕。}他又说：\BibleQuote{我仰望你的判断，我要时时刻刻遵守你的法律。}另一位大声宣告说：\BibleQuote{我的眼目时常仰望主。}同他一道还有人说：\BibleQuote{我心的默念常在你面前。}\Person{保禄}更劝勉说：\BibleQuote{时时感恩，不断祈祷。}那些如此恒心操练的人，全然等候主，并且说：\BibleQuote{让我们认识主，竭力追求认识祂。祂必出现如黎明，祂必临到我们像甘霖，像春云时雨滋润大地。}\BibleRef{欧 6:3} 因为祂不仅在清晨满足他们，也不仅按他们所求的给他们水喝，而是按祂丰盛的慈爱，丰丰富富地赐给他们，常将圣神的\OriginalTerm{恩宠}{grace}赐予他们。并且祂立刻接着说他们所渴慕的是甚么：\BibleQuote{信从我的人。}因为，\BibleQuote{如同冷水使口渴的人欢畅}\BibleRef{若 7:38}，按\BibleBook{箴言}{}所说，照样，对于那些信从主的人而言，圣神的降临胜过一切畅快与喜乐。\par

2. 那么，在这逾越节的日子里，我们应当与圣徒们一同清晨早起，以全灵、以身体的洁净、以宣认和对祂的虔诚信德接近主；这样，当我们在此地先畅饮了、并且满溢了从祂流出的神圣之水时，我们便能与天上的圣人同席坐席，并能分享那地方同一的欢乐之声。罪人们因此被理所当然地赶出，因为他们厌烦此事，并且听见这话：\BibleQuote{朋友，你没有穿婚宴礼服，怎么进到这里来？}\BibleRef{玛 22:12} 罪人固然也渴，但并不渴求圣神的恩宠；倒被邪恶所煽动，完全被宴乐焚烧，正如\BibleBook{箴言}{}所说：\BibleQuote{他终日贪求邪恶的贪欲。}但先知大声反对他们，说：\BibleQuote{祸哉，那些清晨早起追逐烈酒的人！他们流连到晚上，以致酒火焚烧他们。}而且，由于他们放纵淫乱，他们竟敢渴求毁灭他人。他们先喝了虚谎和不可信的水，就有先知所论的那些事临到他们；他说：\BibleQuote{我的创伤为何长久不愈？为何不肯得治？你真的对我像虚假干涸的溪水一般，全不可信！}\BibleRef{耶 15:18} 其次，当他们与同伴共饮时，他们误导并扰乱了正直的心灵，使单纯的人偏离正路。他喊着说甚么呢：\BibleQuote{祸哉，那将浑浊的毁灭之水给邻舍喝，又使他喝醉，好观看他隐秘之事的人。}至于那些假冒为善、偷窃真理的人，他们熄灭了内心。他们喝了这些之后，便继续说那\BibleBook{箴言}{}中淫妇所说的话：\BibleQuote{偷来的水是甜的，暗中吃的饼是好的。}\BibleRef{箴 9:17} 他们暗中布置罗网，因为他们没有德行的自由，也没有\OriginalTerm{智慧}{Sophia}的胆量；智慧在城门口赞扬自己，在通衢大道上放胆发言，在高墙上宣讲。为此，他们被吩咐\Quote{欣然抓住}，因为他们有信德与宴乐之间的抉择，他们偷窃了真理的甘甜，并伪装了自己的苦水，以逃避那原本急速而公开的邪恶之诟。因此，豺狼披上羊皮，粉饰的坟墓以外表白欺人。那就是\Person{撒殚}。\par

\SourceRecord{Translated by R. Payne-Smith.；Nicene and Post-Nicene Fathers, Second Series；Vol. 4.；Edited by Philip Schaff and Henry Wace.；Buffalo, NY: Christian Literature Publishing Co.,；1892.；<http://www.newadvent.org/fathers/2806020.htm>.}

% structure: p0001 source=h1 target=section reason=page-title

\section{选自第二十二封书信}

\Emph{（约三五〇年）}\par

我们的主耶稣基督，为了众人而受死，伸开双手，不是在地上某处，而是在空中，为要显示那由十字架完成的救恩\OriginalTerm{救恩}{Salvation}是为全人类共有的：祂摧毁了那在空中运行的魔鬼\OriginalTerm{魔鬼}{the devil}，祝圣了我们上升天堂的道路，并使之畅通无阻。\par

\SourceRecord{Translated by R. Payne-Smith.；Nicene and Post-Nicene Fathers, Second Series；Vol. 4.；Edited by Philip Schaff and Henry Wace.；Buffalo, NY: Christian Literature Publishing Co.,；1892.；<http://www.newadvent.org/fathers/2806022.htm>.}

% structure: p0001 source=h1 target=section reason=page-title

\section{书信二十四}

\Emph{（写于352年）}\par

当初，当他们出离并越过\Place{埃及}时，他们的敌人成了大海的玩物；但如今，当我们由地进入天堂时，\Person{撒殚}自己从此便如闪电般从天坠落。\par

\SourceRecord{Translated by R. Payne-Smith.；Nicene and Post-Nicene Fathers, Second Series；Vol. 4.；Edited by Philip Schaff and Henry Wace.；Buffalo, NY: Christian Literature Publishing Co.,；1892.；<http://www.newadvent.org/fathers/2806024.htm>.}

% structure: p0001 source=h1 target=section reason=page-title

\section{书信第二十七}

\Emph{（主后355年。）\Place{亚历山大}主教及宣信者\Person{亚大纳削}的第二十七封\WorkTitle{庆节书信}；其开端曰：\Quote{再次，生活的逾越节的日子来临了。}}\par

因为谁是我们的喜乐和夸耀，岂不是我们的主和救主\Person{耶稣基督}么？祂为我们受苦，并藉着自己使我们认识圣父。因为祂并非另一位，就是古时藉众先知说话的那一位；但现在祂对每个人说：\Quote{我就是那说话者，临近了。}\BibleRef{若 4:26} 这话说得真好，因为祂并非一时说话，一时缄默；而是持续不断，从起初就没有停息，祂唤醒每个人，并对每个人的心说话。\par

\SourceRecord{Translated by R. Payne-Smith.；Nicene and Post-Nicene Fathers, Second Series；Vol. 4.；Edited by Philip Schaff and Henry Wace.；Buffalo, NY: Christian Literature Publishing Co.,；1892.；<http://www.newadvent.org/fathers/2806027.htm>.}

% structure: p0001 source=h1 target=section reason=page-title

\section{选自 \WorkTitle{第28封信}}

\Emph{（为356年）}\par

……为使祂得以成为我们众人的祭献，我们，受到真理之言的滋养，并分享祂活泼的教导，能够与圣人一同领受天堂的喜乐。因为在那里，正如祂召唤门徒进入楼上的大厅，同样，\OriginalTerm{圣言}{the Word}也呼召我们与他们同赴那神圣而不朽的筵席；祂在此为我们受苦，却在那里为那些最乐意听从召唤、不懈地，[注视]着目标、追求他们崇高召叫之奖赏的人预备天上的帐幕；那里，对于赴筵并与阻碍他们者争斗的人，存留着冠冕和不朽的喜乐。因为，虽从人性角度说，如此旅程的劳苦是大的，但救主自己已使它变得轻松而仁慈。\par

\SourceRecord{Translated by R. Payne-Smith.；Nicene and Post-Nicene Fathers, Second Series；Vol. 4.；Edited by Philip Schaff and Henry Wace.；Buffalo, NY: Christian Literature Publishing Co.,；1892.；<http://www.newadvent.org/fathers/2806028.htm>.}

% structure: p0001 source=h1 target=section reason=page-title

\section{摘自第29封书信}

\Emph{（357年）摘自第二十九封书信，其开头是\Quote{我们已写的，为当前时刻已足够。}}\par

主考验了门徒们，当祂在枕头上睡着之时，那时行了一个奇迹，这奇迹特别足以令恶人也蒙羞。因为当祂起来，斥责了海，平息了风暴时，祂清楚地表明了两件事：海上的风暴并非来自风，而是出于对行走于海上的主之恐惧；并且斥责它的主并非受造物，而是它的创造者，因为受造物不会服从另一受造物。因为尽管\Place{红海}先前被\Person{梅瑟}分开\BibleRef{出 14:21}，然而那并非\Person{梅瑟}所为，因为那事之所以成就，不是因为他说话，而是因为天主命令。而如果太阳停在\Place{基贝红}\BibleRef{苏 10:12}，月亮停在\Place{阿雅隆}谷，这也不是农的儿子\OriginalTerm{若苏厄}{Jesus son of Nun}的工作，而是主的工作，是主垂听了他的祈祷。正是祂斥责了海，并在十字架上使太阳昏暗\BibleRef{玛 27:45}。\par

% structure: p0004 source=h2 target=subsection reason=relative-heading-rank; base=h2

\subsection{另一片段}

至于人的事物都有终结，神圣的事物却没有。因此之故，当我们死去，当我们的本性疲惫不堪时，祂使我们复活，并引领我们这些虽生于尘土者上升天堂。\par

\SourceRecord{Translated by R. Payne-Smith.；Nicene and Post-Nicene Fathers, Second Series；Vol. 4.；Edited by Philip Schaff and Henry Wace.；Buffalo, NY: Christian Literature Publishing Co.,；1892.；<http://www.newadvent.org/fathers/2806029.htm>.}

% structure: p0001 source=h1 target=section reason=page-title

\section{选自书信三十九}

\Emph{（论367年）关于教会所接受的具体书卷及其数目。摘录自圣\Person{亚大纳修}，\Place{亚历山大}的主教，论复活节庆的第三十九封书信；他在其中以正典方式确定了教会所接受的神圣书卷。}\par

...1. 他们编造了称为表册的书，在其中展示星辰，并以\OriginalTerm{圣人}{Saints}之名命名。确实，他们因此加倍羞辱了自己：编写这些书的人，因他们完善了一种虚谎可鄙的学识；至于愚昧无知的人，他们则用邪恶的思想，使他们对那建立在一切真理之上、在天主面前正直无误的\OriginalTerm{正确信德}{right faith}，引人走入歧途。\par

...2. 但既然我们提到\OriginalTerm{异端者}{heretics}已死，而我们却拥有为得\OriginalTerm{救恩}{salvation}的\OriginalTerm{圣经}{Divine Scriptures}；又因为我担心，正如\Person{保禄}致\Place{格林多}人书中所写，\BibleRef{格后 11:3}，少数\OriginalTerm{纯朴的人}{simple}可能因某些人的诡诈，偏离了他们的纯朴清洁，从此去读别的书——那些称为\OriginalTerm{伪经}{apocryphal}的——因着名字与真经相似而被误导；我恳请你们忍耐，若我也为记念之故，写下你们已经知道的事，乃是出于教会急需与益处的影响。\par

...3. 在开始提及这些事时，为引荐我的工作，我要效法\OriginalTerm{圣史}{Evangelist}\Person{路加}的榜样，以自己的口吻说：\BibleQuote{因为已有许多人着手，} \BibleRef{路 1:1} 将那些称为\OriginalTerm{伪经}{apocryphal}的书卷整理编订，并把它们与受天主默感的圣经相混淆——关于这圣经，我们已完全信服，正如那些从起初亲眼看见并为\OriginalTerm{圣言}{Word}的服役者，所传授给教父们的；既蒙真弟兄们的催促，又从一开始就学习过，我也觉得应当把那些列入\OriginalTerm{正典}{Canon}、传承下来并被认可为神圣的书卷，陈列在你们面前；好使那陷入错误的人，可以谴责那诱导他走入歧途的人；而那在纯洁中保持坚定的人，又可以因这些事被再次提醒而喜乐。\par

...4. 那么，旧约的书卷，共有二十二卷；因为，据我所闻，这是希伯来人字母的数目；它们的次序与名称如下。第一卷是\BibleBook{创世}{纪}，然后是\BibleBook{出谷}{纪}，接着是\BibleBook{肋未}{纪}，之后是\BibleBook{户籍}{纪}，再后是\BibleBook{申命}{纪}。随后有\BibleBook{若苏厄}{书}，接着是\BibleBook{民长}{纪}，然后是\BibleBook{卢德}{传}。然后，再下来是四卷\BibleBook{列王}{纪}，前两卷合为一卷，后两卷也合为一卷。再是\BibleBook{编年}{纪}，前、后二卷合为一卷。\BibleBook{厄斯德拉}{书}前、后二卷也合为一卷。之后是\BibleBook{圣咏}{集}，接着是\BibleBook{箴}{言}，再是\BibleBook{训道}{篇}，以及\BibleBook{雅}{歌}。然后是\BibleBook{约伯}{传}，之后是先知书，其中十二小先知书合为一卷。然后\BibleBook{依撒意亚}{书}一卷，\BibleBook{耶肋米亚}{书}连同\Person{巴路克}、哀歌和那封书信，合为一卷；再后是\BibleBook{厄则克耳}{书}和\BibleBook{达尼尔}{书}，各一卷。以上就是旧约。\par

...5. 再来说新约的书卷，也并不乏味。这些是：四福音，即\BibleBook{玛窦}{福音}、\BibleBook{马尔谷}{福音}、\BibleBook{路加}{福音}、\BibleBook{若望}{福音}。然后是\BibleBook{}{宗徒大事录}，以及七封公函：即\BibleBook{雅各伯}{书}一封，\BibleBook{伯多禄}{前书}、\BibleBook{伯多禄}{后书}两封，\BibleBook{若望}{一书}、\BibleBook{若望}{二书}、\BibleBook{若望}{三书}三封，之后\BibleBook{犹达}{书}一封。此外，还有\Person{保禄}书信十四封，其顺序如下：第一是\BibleBook{罗马}{书}，然后\BibleBook{格林多}{前书}、\BibleBook{格林多}{后书}，以后是\BibleBook{迦拉达}{书}，再是\BibleBook{厄弗所}{书}，再是\BibleBook{斐理伯}{书}，再是\BibleBook{哥罗森}{书}，以后是\BibleBook{得撒洛尼}{前书}、\BibleBook{得撒洛尼}{后书}，\BibleBook{希伯来}{书}，再是\BibleBook{弟茂德}{前书}、\BibleBook{弟茂德}{后书}，\BibleBook{弟铎}{书}，最后是\BibleBook{费肋孟}{书}。此外还有\BibleBook{若望}{默示录}。\par

...6. 这些是救恩的泉源，使渴慕的人能因其中所蕴藏的活语而满足。惟有在这些书中，才宣扬着虔敬的教训。谁也不可添加，谁也不可删减。因为关于这些书，主曾使\OriginalTerm{撒杜塞人}{Sadducees}蒙羞，说：\BibleQuote{你们错了，因为不明了经书。} 祂也责备\OriginalTerm{犹太人}{Jews}说：\BibleQuote{你们查考经典，因为你们认为其中有永生，正是这些为我作证。} \BibleRef{玛 22:29}\par

...7. 但为了更为精确，我出于必需要加上这一点：在这些书以外，还有一些书，虽未列入\OriginalTerm{正典}{Canon}，却是教父们指定给那些新近加入我们、并切望学习虔敬之道的人去阅读的。即\WorkTitle{智慧篇}、\WorkTitle{德训篇}、\WorkTitle{艾斯德尔传}、\WorkTitle{友弟德传}、\WorkTitle{多俾亚传}、所谓的\WorkTitle{宗徒训诲录}和\WorkTitle{牧者}。前者，弟兄们，是正典中的书卷；后者则仅[供]阅读；且在任何地方都未提及\OriginalTerm{伪经}{apocryphal writings}。伪经是\OriginalTerm{异端者}{heretics}的虚构，他们随意编写，予以认可，并指派日期，好能把这些书当作古老的著作，从而寻机误导\OriginalTerm{纯朴的人}{simple}。\par

\SourceRecord{Translated by R. Payne-Smith.；Nicene and Post-Nicene Fathers, Second Series；Vol. 4.；Edited by Philip Schaff and Henry Wace.；Buffalo, NY: Christian Literature Publishing Co.,；1892.；<http://www.newadvent.org/fathers/2806039.htm>.}

% structure: p0001 source=h1 target=section reason=page-title

\section{选自书信四十}

\Emph{（写于368年）}\par

\BibleQuote{你们在我的试探中，始终同我在一起；因此，我将王权给你们预备下，正如我父给我预备了一样，为使你们在我的国里，一同在我的筵席上吃喝。}我们既蒙召赴那宏大而属天的盛宴，进入那已打扫干净的上层厅堂，就让我们如宗徒所劝勉的，\BibleQuote{洁净自己，除去肉身和灵魂的一切玷污，以敬畏天主之情成就圣德}\BibleRef{格后 7:1}；这样，我们内外无玷——外，以节制和正义为衣；内，靠着圣神，正确分解真理的话——便能听到：\BibleQuote{进入你主人的福乐吧！}\BibleRef{玛 25:21}。\par

\SourceRecord{Translated by R. Payne-Smith.；Nicene and Post-Nicene Fathers, Second Series；Vol. 4.；Edited by Philip Schaff and Henry Wace.；Buffalo, NY: Christian Literature Publishing Co.,；1892.；<http://www.newadvent.org/fathers/2806040.htm>.}

% structure: p0001 source=h1 target=section reason=page-title

\section{书信四十二}

\Emph{（作于370年。）}\par

弟兄们，因为我们已蒙召，且现在被\OriginalTerm{智慧}{Wisdom}，并依照福音的比喻，召集到那伟大而天上的盛宴，足为一切受造物享用；我指的是\OriginalTerm{逾越节}{Passover}——指向被祭献的基督；因为\BibleQuote{基督，我们的逾越节羔羊，已被祭献。}\BibleRef{格前 5:7} \Emph{（后来又说：）}因此，那些如此预备好的人将听到，\BibleQuote{进入你主人的福乐吧！}\BibleRef{玛 25:21}\par

\SourceRecord{Translated by R. Payne-Smith.；Nicene and Post-Nicene Fathers, Second Series；Vol. 4.；Edited by Philip Schaff and Henry Wace.；Buffalo, NY: Christian Literature Publishing Co.,；1892.；<http://www.newadvent.org/fathers/2806042.htm>.}

% structure: p0001 source=h1 target=section reason=page-title

\section{选自\WorkTitle{书信四十三}}

\Emph{（写于371年）}\par

那么，逾越节也是我们的，我们的召唤是来自天上的，\BibleQuote{我们的家乡原是在天上}，正如\Person{保禄}所说；\BibleQuote{因为我们在此没有常存的城邑，而是寻求那将来的城邑}，我们仰望它，因而适当地守这节日。\Emph{（又，此后：）} 天确实是高的，距离我们无限遥远；因为经上说，\BibleQuote{天是上主的天}。但我们不能因此疏忽或惧怕，好像通往那里的路是不可能的；相反，我们应当更加热心。然而，我们并不需要像当初那些从东方迁移，在\Place{史纳尔}地方找到一片平原，便开始[建造一座塔]的人那样，用火烤砖，寻找沥青作灰泥；因为他们的语言被混乱了，他们的工程被毁灭了。但是，上主藉着自己的血，为我们祝圣了一条道路，并使之变得容易。\Emph{（又：）} 因为他不但在这距离上给了我们安慰，而且他也已经来到，为我们打开了那从前关闭的门。确实，从他因亚当喜爱乐园而将他赶出，并安置\OriginalTerm{革鲁宾}{Cherubim}和四面旋转\OriginalTerm{火剑}{flaming sword}，把守通往\OriginalTerm{生命树}{tree of life}的道路之时起，这门就关闭了；而今，却大大敞开了。那坐在革鲁宾之上的，带着更大的恩宠和慈爱显现，将那位认他的强盗一同带入乐园，并且作为我们的前驱进入了天堂，为众人开启了大门。\Emph{（又：）} \Person{保禄}也\BibleQuote{朝着目标奔跑，要得到天主在基督耶稣内召我向上的奖品}\BibleRef{斐 3:14}，藉此被提升到第三重天，看见了那些天上的事物，然后又降下来，教导我们，向希伯来人宣讲经上所记载的，说：\BibleQuote{你们原来并不是走近了那可触摸的山，那里有烈火、浓云、黑暗、暴风、号筒的响声以及说话的声音，然而你们却接近了\Place{熙雍}山和永生天主的城，天上的\Place{耶路撒冷}，接近了千万天使的盛会，和那些已被登记在天上的首生者的集会}。谁不愿意享受与这些崇高同伴的交往！谁不希望被登记在这些同伴之中，好能同他们一起听到：\BibleQuote{我父所祝福的，你们来吧，承受自创世以来为你们预备的国度吧}\BibleRef{玛 25:34}。\par

\SourceRecord{Translated by R. Payne-Smith.；Nicene and Post-Nicene Fathers, Second Series；Vol. 4.；Edited by Philip Schaff and Henry Wace.；Buffalo, NY: Christian Literature Publishing Co.,；1892.；<http://www.newadvent.org/fathers/2806043.htm>.}

% structure: p0001 source=h1 target=section reason=page-title

\section{选自书信四十四}

\Emph{（372年）再次选自第四十四封书信，其开篇为：\Quote{我们的主和救主耶稣基督代替我们并为我们所做的一切。}}\par

当司祭长和经师的仆人们看到这些事，又听到耶稣说：\BibleQuote{谁若渴，到我这里来喝吧；}\BibleRef{若 7:37}便意识到这人并非像他们一样的凡人，而是那位赐水给圣人的，并且正是依撒意亚先知所预言的那位。因为他确实是\OriginalTerm{光明的辉耀}{splendour of the light}，是天主的圣言。因此，他像由泉源流出的江河，在古时也曾灌溉乐园；但如今他赐给所有人同样的圣神恩赐，并说：\BibleQuote{谁若渴，到我这里来喝吧。}凡\BibleQuote{信从我的，就如经上所说：从他的腹中要流出活水的江河。}\BibleRef{若 7:37}。这些话不是凡人能说的，而是生活天主说的，他确实赐予生命，并赐下圣神。\par

\SourceRecord{Translated by R. Payne-Smith.；Nicene and Post-Nicene Fathers, Second Series；Vol. 4.；Edited by Philip Schaff and Henry Wace.；Buffalo, NY: Christian Literature Publishing Co.,；1892.；<http://www.newadvent.org/fathers/2806044.htm>.}

% structure: p0001 source=h1 target=section reason=page-title

\section{书信四十五}

\Emph{（公元373年）}\par

让我们都拿起我们的祭献，留意分施给穷人，并进入圣所，如经上所载：\BibleQuote{为了我们作前驱的耶稣已为我们进入了，获得了永远的救赎。} ...（\Emph{同上：}）...并且这是一个伟大的证据，我们本来是外人，如今却被称为朋友；从前是外方人，如今却成了与圣人们同国的公民，并被称为那在上的\Place{耶路撒冷}的儿女，这耶路撒冷，正是\Person{撒罗满}所建的\OriginalTerm{预像}{type}。因为如果\Person{梅瑟}是按照在山上所指示的样式制造一切，那么，在帐幕内所行的敬礼显然就是天上\OriginalTerm{奥秘}{mysteries}的预像，主愿意我们进入那里，就为我们预备了一条新的、永存的道路。并且正如一切旧的事物都是新事物的预像，现今的庆节也是天上喜乐的预像，我们应以圣咏和灵歌，进入这喜乐，让我们开始斋戒吧。\par

\SourceRecord{Translated by R. Payne-Smith.；Nicene and Post-Nicene Fathers, Second Series；Vol. 4.；Edited by Philip Schaff and Henry Wace.；Buffalo, NY: Christian Literature Publishing Co.,；1892.；<http://www.newadvent.org/fathers/2806045.htm>.}

% structure: p0001 source=h1 target=section reason=page-title

\section{第46封信}

\Emph{自\Place{萨底卡}致\Place{马雷奥提斯}书}，公元343--344年。\par

\Person{亚大纳削}致\Place{马雷奥提斯}天主教会的司铎、执事及民众，即我们所爱慕和思念的弟兄们，在主内问候你们。\par

神圣公会议赞许了你们在基督内的虔诚。他们全都认可你们在一切事上的精神和刚毅，因为你们不畏惧恐吓，并且虽然为你们的虔诚不得不忍受凌辱与迫害，你们仍然坚守不渝。你们的书信向他们全体宣读时，令人潸然泪下，激起普遍的同情。他们虽不在你们跟前，却爱着你们，将你们的迫害视同己受。他们致你们的信就是其情谊的证明；尽管将你们与神圣的\Place{亚历山大里亚}教会相提并论已足够，但神圣公会议仍单独致函你们，为使你们因此得鼓励，不因所受苦难而退缩，反要感谢天主；因为你们的坚忍必将结出善果。\par

先前，异端者的品行并不明显。但现在已显露出来，并向所有人公开了。因为神圣公会议注意到了这些人向你们捏造的诽谤，并且憎恶之，经全体同意，在\Place{亚历山大里亚}和\Place{马雷奥提斯}罢免了\Person{德奥多若}、\Person{瓦伦斯}及\Person{乌尔撒奇乌斯}。同样的通知也已送达其他教会。由于他们对各教会所施行的残忍与暴政已不可再容忍，他们已被革除主教职，并被逐出所有人的共融。至于\Person{额我略}，他们甚至不愿提起；因为此人连主教的名分都不具备，他们认为提他的名字是多余的。但是，为了那些受他欺骗的人，他们提及了他的名字；并非因他值得被提及，而是为使那些受骗者由此认清他的恶名，并羞愧于与何等样人相通。从先前的文件你们将得知有关他们的记述；尽管并非所有主教都前来签署，然而信函是由全体起草的，并代全体签了名。\BibleQuote{你们要以圣吻彼此问候。}所有的弟兄问候你们。\par

\Signature{我，普洛托根尼斯，主教，愿你们在主内得蒙保守，我所爱与思念的。}\par

\Signature{我，雅典诺多鲁斯*，主教，愿你们在主内得蒙保守，至爱的弟兄们。}［其他署名］\Signature{犹利安}、\Signature{阿摩尼乌斯}、\Signature{阿普里安}、\Signature{玛尔凯路斯}、\Signature{格隆提乌斯*}、\Signature{波尔菲里乌斯*}、\Signature{佐西默斯}、\Signature{阿斯克勒庇乌斯}、\Signature{阿皮安}、\Signature{欧洛吉乌斯}、\Signature{欧根尼乌斯}、\Signature{李奥多鲁斯（26）}、\Signature{玛尔提里乌斯}、\Signature{欧卡尔普斯}、\Signature{路基乌斯*}、\Signature{卡洛厄斯}。\Signature{马克西穆斯：藉着来自\Place{高卢}的书信，我愿你们在主内得蒙保守，亲爱的。}\Signature{我们，阿尔基达穆斯和菲洛克塞努，司铎，以及执事莱奥，来自\Place{罗马}，愿你们得蒙保守。}\Signature{我，高登提乌斯，\Place{奈苏斯}的主教，愿你们在主内得蒙保守。}［亦］\Signature{\Place{潘诺尼亚}的\Place{梅里亚}的弗洛伦提乌斯}、\Signature{\Place{潘诺尼亚}的\Place{卡斯泰卢姆}的阿米阿努斯（9）}、\Signature{\Place{贝内文托}的雅努阿里乌斯}、\Signature{\Place{潘诺尼亚}的\Place{纳尔基多农}的普雷泰克斯塔图斯}、\Signature{\Place{色萨利}的\Place{希帕塔}的希佩尔内里斯（48）}、\Signature{\Place{凯撒奥古斯塔}的卡斯图斯}、\Signature{\Place{色萨利}的\Place{卡尔基苏斯}的塞维鲁}、\Signature{\Place{七城的泰拉}的犹利安}、\Signature{\Place{维罗纳}的路基乌斯}、\Signature{\Place{赫克莱亚尔·基克比奈}的欧根尼乌斯（35）}、\Signature{\Place{阿普利亚}的\Place{利克尼·苏诺西翁}的佐西默斯（92）}、\Signature{\Place{锡凯翁}的赫尔摩格尼斯}、\Signature{\Place{马加拉}的特里佛斯}、\Signature{\Place{卡斯皮}的帕雷哥里乌斯*}、\Signature{\Place{卡斯特罗马蒂斯}的卡洛厄斯（21）}、\Signature{\Place{锡科尼斯}的伊雷内乌斯}、\Signature{\Place{吕皮亚农}的马凯多尼乌斯}、\Signature{\Place{纳乌帕克蒂}的玛尔提里乌斯}、\Signature{\Place{迪乌斯}的帕拉迪乌斯}、\Signature{\Place{高卢}的\Place{卢格杜努姆}的布罗修斯（87）}、\Signature{\Place{布雷西亚}的乌尔撒奇乌斯}、\Signature{\Place{维米纳基乌姆}的阿曼提乌斯，由司铎马克西穆斯代签}、\Signature{\Place{阿哈伊亚}的\Place{吉帕拉}的亚历山大}、\Signature{\Place{莫托纳}的欧提基乌斯}、\Signature{\Place{潘诺尼亚}的\Place{佩塔维奥}的阿普里安}、\Signature{\Place{马其顿}的\Place{帕勒内}的安提哥努斯}、\Signature{\Place{阿卡里亚·康斯坦蒂亚斯}的多美提乌斯*}、\Signature{\Place{埃诺罗多佩}的奥林匹乌斯}、\Signature{\Place{奥雷奥马尔加}的佐西默斯}、\Signature{\Place{米兰}的普罗塔西乌斯}、\Signature{\Place{萨瓦河畔的西西亚}的马尔谷}、\Signature{\Place{阿哈伊亚}的\Place{奥普斯}的欧卡尔普斯}、\Signature{\Place{阿非利加}的\Place{维尔塔拉}的维塔利斯*}、\Signature{\Place{提尔塔纳}的赫利阿努斯}、\Signature{\Place{克里特}的\Place{赫拉皮塔}的辛佛鲁斯}、\Signature{\Place{赫拉克拉}的莫西尼乌斯（64）}、\Signature{\Place{基萨穆斯}的欧基苏斯}、\Signature{\Place{基多尼亚}的基多尼乌斯}。\par

\SourceRecord{Translated by Archibald Robertson.；Nicene and Post-Nicene Fathers, Second Series；Vol. 4.；Edited by Philip Schaff and Henry Wace.；Buffalo, NY: Christian Literature Publishing Co.,；1892.；<http://www.newadvent.org/fathers/2806046.htm>.}

% structure: p0001 source=h1 target=section reason=page-title

\section{书信四十七}

\Emph{致\Place{亚历山大}教会，论同一事由。}\par

\Person{亚大纳削}致\Place{亚历山大}及\Place{帕伦波拉}圣天主教会所有司铎与执事，最亲爱的弟兄们，问候。\par

我最亲爱的弟兄们，在写这封信时，我必须首先感谢基督。但现在这尤其合宜，因为主所成就的许多大事理应获得我们的感谢，而信祂的人不应为祂的众多恩惠而忘恩。因此我们感谢主，祂常使我们在信仰上显明于众人，祂也在此时为教会行了诸多奇事。因为\Person{优西比乌}的异端党派及\Person{亚略}的继承者们所主张和散播的一切，聚集的主教们已宣布为虚假和捏造。那些被许多人视为可畏的人，像那些被称为巨人的，竟被算为虚无，这也理所当然；因为正如光明来到，黑暗就被照亮；同样，正义者来到，邪恶就被揭露；当善者临在时，无价值者便暴露无遗。\par

因为，亲爱的，你们自己并非不知道那恶名昭彰的\Person{优西比乌}异端的继承者，即\Person{德奥多若}、\Person{那尔基索}、\Person{瓦伦斯}、\Person{乌尔撒奇}，以及他们中最坏的\Person{乔治}、\Person{斯德望}、\Person{亚加修}、\Person{梅诺凡托}及其同党所行的事，因为他们的疯狂已显明于众人；他们对诸教会所犯的罪行也没有逃过你们的眼目。因为你们是首先被他们伤害的，你们的教会是首先被他们企图败坏的。但行了许多大事、且如前所述在众人心目中极为可怖的他们，竟惊恐得超乎想象。因为他们不但畏惧\Place{罗马}会议，不但在受邀时推辞，就是现在到了\Place{撒尔迪加}，他们也是良心备受谴责，以至见到审判者时便惊惶失措。于是他们在心中昏厥。诚然，有人会对他们说：\BibleQuote{死亡，你的刺在哪里？死亡，你的胜利在哪里？}因为事情并未如他们所愿，让他们能随心所欲地判决；这一次他们无法再欺骗他们想欺骗的人了。但他们看见了忠实的人们，关心正义的人；不，更确切地说，他们看见吾主亲自在他们中间，就像古时从坟墓里出来的魔鬼一样；因为身为虚谎之子，他们受不了见真理。于是\Person{德奥多若}、\Person{那尔基索}、\Person{乌尔撒奇}和他们的同伙这样说：\Quote{且住，我们与你何干，基督的人？我们知道你是真实的，并怕被定罪；我们不愿当面承认我们的诽谤。我们与你无干；因为你们是基督徒，而我们是基督的仇敌；在你们那里真理强大，我们却学会了欺骗。我们以为所行的事是隐藏的；我们没想到现在要受审判；你们为何在我们的时期未到之前揭露我们的行为？为何揭露我们，在日子之前就折磨我们？}虽然他们品质极坏，行走在黑暗中，但他们终于学到了光明与黑暗不能相通，基督与\Person{贝里雅耳}不能和谐。因此，亲爱的弟兄们，他们既知道自己的所作所为，又看见了受害者们已准备作原告，且证人就在眼前，便效法\Person{加音}，像他一样逃跑了；他们这样大大地误入歧途，因为他们模仿了他的逃跑，也就接受了他的处罚。因为神圣的会议知道他们的作为；它已听见我们的血高声呼喊，亲自听见了受伤者的声音。所有主教都知道他们如何犯罪，以及他们对我们的众教会和其他教会行了多少事；因此，他们已把这些像\Person{加音}一样的人逐出教会。因为当你们的信被宣读时，谁不流泪？谁不叹息，看见那些人放逐的是谁？谁不把你们的磨难看作自己的？最亲爱的弟兄们，从前他们加害于你们时，你们受了苦，或许战争停息才不久。但如今，所有聚集并听见你们所受苦难的主教们，都如同你们当初受害时一样悲伤哀叹，并分享了那时的忧伤...\par

因着这些行为，以及他们对各教会所犯的其他一切恶行，圣大公会议已将他们全部罢免，不仅判定他们与教会隔绝，更认为他们不配被称为基督徒。因为否认基督的人怎能被称为基督徒？做恶事反对各教会的人怎能被接纳进入教会？因此，圣会议已致信各地教会，使这些事在众人中显明，好让那些受他们欺骗的人，如今能回归完备的确信与真理。因此，亲爱的弟兄们，不要丧胆；要如同天主的仆人，宣认基督信仰的人，在主内经受考验；不要因患难而灰心，也不要因那些图谋害你们的异端者所带来的烦恼而忧愁。因为全世界都同情你们的忧伤，更甚者，整个世界记念着你们众人。如今我想，那些受他们欺骗的人，一旦看到会议的严厉判决，就会离开他们并弃绝他们的不虔敬。然而，即使在此之后他们仍举手反抗，你们也不要惊讶，如果他们发怒也不要害怕；却要祈祷，举起你们的手向天主，并且确信主必不迟延，必会按照你们的意愿成就一切。我确实希望能给你们写一封更长的信，详细叙述所发生的一切；但既然长老和执事们能当面将他们的所见告诉你们，我就节制不多写了。我只嘱咐你们一件事，视之为必要：你们要将敬畏主放在眼前，把主放在首位，并以你们惯常的同心合意办理一切，如同有智慧明达的人。请为我们祈祷，记念寡妇们的急需，尤其因为真理的敌人夺取了属于她们的。但让你们的爱德胜过异端者的恶意。因为我们相信，因着你们的祈祷，主必施恩，让我快快见到你们。同时，你们将从所有主教写给你们的信中得知会议的进程，并从附上的信函中了解对\Person{狄奥多尔}、\Person{纳尔基索斯}、\Person{斯德望}、\Person{阿卡西乌斯}、\Person{乔治}、\Person{梅诺凡图斯}、\Person{乌尔撒基乌斯}和\Person{瓦伦斯}的罢免。至于\Person{格里高利}，他们不愿提及，因为认为提名一个根本不配主教之名的人是多余的。但为了那些受他欺骗的人，他们提到了他的名字，并非他的名配被提及，而是为使受他欺骗的人知道他的丑行，并为他们与之共融的那种人感到羞愧。...我祈求你们得蒙保守在主内，至爱和深切思念的弟兄们。\par

\SourceRecord{Translated by Archibald Robertson.；Nicene and Post-Nicene Fathers, Second Series；Vol. 4.；Edited by Philip Schaff and Henry Wace.；Buffalo, NY: Christian Literature Publishing Co.,；1892.；<http://www.newadvent.org/fathers/2806047.htm>.}

% structure: p0001 source=h1 target=section reason=page-title

\section{第48封信}

\Emph{致\Person{阿孟}的信。写于公元354年之前。}\par

凡天主所造的都是美好纯洁的，因为圣言没有造任何无用或不洁之物。因为 \BibleQuote{我们就是献与天主的基督的馨香} \BibleRef{格后 2:15}，正如宗徒所说的。然而魔鬼的毒箭既多样又诡诈，他设法搅扰那些心思单纯的人，并试图阻挠弟兄们通常的操练，在他们中间暗中散播不洁与污秽的思想；来，让我们靠着救主的恩宠，简短地驱散恶者的错谬，并坚固纯朴人的心。因为 \BibleQuote{为洁净的人一切都是洁净的} \BibleRef{铎 1:15}，但为不洁净的人，连他们的良心和所有的一切都污秽了。我也惊讶魔鬼的狡猾：他本是腐败与灾祸，却在纯洁的外表下暗示种种思想；但这不过是一个陷阱，而非考验。正如我之前说过的，他旨在使隐修士们分心，离开他们惯常而有益的默想，并看似战胜他们，便煽动起这类如嗡嗡之声的思想，于生活毫无裨益，不过是虚妄的问题和琐事，本应弃之一旁。请告诉我，亲爱的，最虔诚的朋友，任何一种自然的排泄物中，有什么罪恶或不洁呢？难道有人竟想把擤鼻涕或口中吐出的痰当作可责的事吗？我们还可以加上腹中的排泄物，那是动物生命生理上所必需的。再者，若我们按照圣经所教导的，相信人是天主的工程，那么任何污秽的工程怎能出自纯洁的全能者呢？并且，按照神圣的\BibleBook{宗徒}{大事录} \BibleRef{宗 17:28}，\BibleQuote{我们也是天主的子孙}，那么我们本身并没有任何不洁。因为只有在我们犯罪，那最污秽的事时，我们才招致污秽。但任何身体的排泄，若出于不由自主，那么我们就如其他事物一样，因着自然的必然性而经历此事。然而，有些人专以反驳正确的话，更准确地说，反驳天主的造化为乐，他们甚至曲解\WorkTitle{福音}中的一句话，声称 \BibleQuote{不是入于口的使人污秽，而是出于口的才使人污秽} \BibleRef{玛 15:11}，我们不得不点明他们这种无理——我甚至不能称之为问题。首先，他们像不稳重的人，强解圣经，自取无知 \BibleRef{伯后 3:16}。这神圣神谕的真义如下：当时有些人，像今天的这些人一样，对食物心存疑虑。主亲自，为消除他们的无知，或许也是揭露他们的诡诈，确言不是入于口的使人污秽，而是出于口的。然后祂确切地加上它们从何处出，即是从心里出来。因为祂知道，那里有邪恶的宝藏，就是亵渎的思想和其他罪恶。但宗徒更简洁地教导同一件事，说：\BibleQuote{食物不能使我们更亲近天主} \BibleRef{格前 8:8}。此外，我们可以合理地认为，任何自然的排泄也不会把我们带到祂面前受罚。但也许医生们（为叫这些人即使在教外人面前也自觉羞愧）会在这点上支持我们，告诉我们动物身体被赋予某些必要的通道，用以排出我们各部位分泌的过剩之物；例如，头部的过剩，有头发和头部的水状排出物，腹部的排泄，以及生殖管道的过剩排泄。那么，因天主的名，至爱天主的尊长，如果造了身体的主愿意并造了这些部分具有这样的通道，何罪之有呢？然而，我们必须对付恶人的异议，因为他们可能会说：\Quote{若器官是由造物主分别塑造的，那么正常使用它们便无罪。}让我们提出这个问题来堵住他们：你们所说的使用是什么意思？是指天主允许的合法使用，当祂说 \BibleQuote{你们要生育繁殖，充满大地} \BibleRef{创 1:28}，而宗徒也肯定说 \BibleQuote{婚姻应受尊重，床笫应是无玷的} \BibleRef{希 13:4}；还是那种公开却偷偷摸摸、以奸淫方式进行的滥用呢？因为生活中的其他事情也一样，我们也会发现依情况而有分别。例如，杀人是不对的，但在战争中，消灭敌人却是合法且值得称赞的；因此，不仅那些在战场上表现卓越的人被认为配得极大的荣耀，而且还树立碑碣宣扬他们的功绩。所以同样的行为，有时在某些情形下是不合法的，而在另一些情形下，在适当的时机，却是合法且容许的。同样的道理也适用于男女关系。那年轻时自由婚配，自然生子的人，是有福的。但他若放纵本能，那宗徒所写的、等待淫乱和奸淫者的惩罚 \BibleRef{希 13:4} 必将临到。\par

因为在生命中，关于这些事，有两条道路。一条是较为温和且普通的，我指的是\OriginalTerm{婚姻}{marriage}；另一条是天使般的和无可比拟的，即\OriginalTerm{童贞}{virginity}。现在，若有人选择了世俗的道路，即婚姻，他固然无可指责，但不会领受像另一条道路那样大的恩赐。因为他也将领受，因为他同样结果实，即\OriginalTerm{三十倍}{thirtyfold}。但若有人拥抱那条圣善而非现世的道路，尽管与前者相比，它崎岖难行，却拥有更奇妙的恩赐：因为它结出完美的果实，即\OriginalTerm{一百倍}{hundredfold}。因此，他们那些不洁与邪恶的反对，早就在\WorkTitle{圣经}中有了恰当的解答。所以，父亲啊，你要坚固你属下的羊群，以\OriginalTerm{宗徒著作}{Apostolic writings}劝勉他们，以福音的教训引导他们，以\BibleBook{圣咏}{集}劝戒他们，并且说：\BibleQuote{依照祢的话使我恢复生命}；但所说的\InnerQuote{祢的话}，乃是指我们应当以纯洁的心侍奉祂。因为知道这一点，先知说，好像在解释自己，\BibleQuote{天主，求祢给我再造一颗纯洁的心}，以免污秽的思念扰乱我。\Person{达味}又说：\BibleQuote{以祢\OriginalTerm{自由的精神}{free spirit}坚固我}，这样即使思念扰乱我，来自祢的某种强大力量会坚固我，成为我的支撑。然后，给予这些及类似的劝告后，对于那些迟于服从真理的人，你要说：\BibleQuote{我要教导恶人祢的道路}，并且深信在主内，你会说服他们停止这样的邪恶，就歌唱：\BibleQuote{罪人将归向祢}。但愿那些提出恶意问题的人停止这种徒劳无功的劳碌，而那些单纯地怀疑的人能以\InnerQuote{自由的精神}获得坚固。至于你们中确实认识真理的人，要在我们的主\Person{耶稣基督}内，保持这真理完整而不动摇；愿光荣和权能，藉着他，归于圣父，及圣神，至于无穷之世。阿们。\par

\SourceRecord{Translated by Archibald Robertson.；Nicene and Post-Nicene Fathers, Second Series；Vol. 4.；Edited by Philip Schaff and Henry Wace.；Buffalo, NY: Christian Literature Publishing Co.,；1892.；<http://www.newadvent.org/fathers/2806048.htm>.}

% structure: p0001 source=h1 target=section reason=page-title

\section{第49封信}

\Emph{致\Person{德拉孔提}的信。写于公元354或355年。}\par

我真不知道该如何写信。我是要责备你的拒绝呢？还是责备你顾及试探，因怕\OriginalTerm{犹太人}{Jews}而躲藏起来？无论如何，亲爱的\Person{德拉孔提}，你所做的都应受责备。因为领受\OriginalTerm{恩宠}{grace}后不该躲藏，身为明智之人，也不该给别人提供逃跑的借口。因为许多人听见这事就跌倒；不仅因为你做了这事，更因为你做这事是顾及时代和压在\OriginalTerm{教会}{Church}身上的苦难。我担心你为了自己逃跑，反而因别人在\OriginalTerm{主}{Lord}面前处于危险之中。因为如果\BibleQuote{凡使这些信我的小子中的一个跌倒的，倒不如拿一块驴拉的磨石，系在他的颈上，沉在海的深处里}\BibleRef{玛 18:6}，那么你若使这么多人跌倒，将有什么等着你呢？因为在\Place{亚历山大里亚}地区，你当选时那惊人的一致必将因你的退隐而瓦解：该地区的\OriginalTerm{主教职务}{episcopate}将被许多人——正如你所深知，许多不合格的人——觊觎。许多\OriginalTerm{外教人}{heathen}曾承诺在你当选后就成为基督徒，如果你的虔诚轻视了所赐的恩宠，他们将继续留在外教中。\par

2. 对于这种行为，你将提出什么辩护？你能用什么论证来洗刷和消除这样的指控？你将如何医治那些因你而跌倒受伤的人？或者你将如何修复破碎的和平？亲爱的\Person{德拉孔提}，你给我们带来了忧伤而非喜乐，哀叹而非安慰。因为我们原期望有你作我们的安慰；如今却见你逃跑，并且知道你将在审判中被定罪，受审时必会后悔。正如\OriginalTerm{先知}{Prophet}所说：\BibleQuote{谁将怜悯你}\BibleRef{耶 15:5}，当人看见\OriginalTerm{基督}{Christ}为之而死的\OriginalTerm{弟兄们}{brethren}因你的逃跑而受伤害时，谁会为你的平安着想呢？因为你必须明白，且不可疑惑：你在当选之前是为自己活，当选之后，你就是为你的羊群活。在你领受主教职务的恩宠之前，没有人认识你；但当你成为主教后，\OriginalTerm{平信徒}{laity}期望你带给他们食粮，即\OriginalTerm{圣经}{Scriptures}的教导。那么当他们期待，忍受饥饿，你却\BibleQuote{只喂养自己}\BibleRef{则 34:2}，而当我们的主耶稣基督来临，我们站在祂面前时，当祂看到自己的羊忍饥挨饿时，你将提出什么辩护？因为如果你没有拿那笔钱，祂不会责备你。但如果你拿了，却把它埋藏起来——愿天主不让你的虔诚听到这样的话：\BibleQuote{你该把我的钱交给钱庄里的人，待我来时，便连本带利收回。}\BibleRef{玛 25:27}\par

3. 我恳求你，顾惜自己和我们。顾惜自己，免得你陷入危险；顾惜我们，免得我们因你而忧伤。你要为教会着想，免得许多小子因你受害，也给别人提供退缩的机会。不，如果你是因为畏惧时势而胆怯行事，你的心思就不像男子汉了；因为在这种情况下，你本当为基督显出热忱，勇敢面对环境，用真福\Person{保禄}的话说：\BibleQuote{然而，靠着那爱我们的主，我们在这一切事上，大获全胜}\BibleRef{罗 8:37}；尤其因为我们应当服事主，而非时机。但如果组织教会令你厌烦，你认为主教职务的职分没有回报，那么，你就是在轻看安排这些事的\OriginalTerm{救主}{Saviour}。我恳求你，抛弃这些想法，也不要容忍那些这样建议你的人，因为这配不上\Person{德拉孔提}。因为主藉\OriginalTerm{宗徒}{Apostles}们所建立的秩序，是美好而稳固的；但弟兄们的怯懦终将停止。\par

4. 因为如果人人都像你现在的建议者那样想法，你又怎能成为基督徒呢，因为那时将没有\OriginalTerm{主教}{bishops}？或者如果我们的继任者将继承这种心态，教会又如何能维系呢？还是你的建议者认为你没有领受什么，因而轻视它？果真如此，他们肯定错了。因为若有轻视者，他们该想想\OriginalTerm{圣洗之泉}{Font}的恩宠是否无关紧要了。但你已领受了，亲爱的\Person{德拉孔提}；不要容忍你的建议者，也不要自欺。因为赐予恩宠的天主将会向你索要。你没有听过宗徒说：\BibleQuote{不要疏忽你心内的神恩}\BibleRef{弟前 4:14}吗？或者你没有读过他如何接纳那赚了双倍钱的人，而定了那埋藏钱的人的罪吗？但愿你快快回来，好使你也能成为受称赞的人之一。或者告诉我，你的建议者希望你效法谁呢？因为我们应当以\OriginalTerm{圣人}{saints}和\OriginalTerm{教父}{fathers}为标准行事，效法他们，并确信若离开他们，我们也就不属于他们的团体了。那么他们希望你效法谁呢？是那个迟延、想要跟从却因家庭缘故耽延并商量的人\BibleRef{路 2:61}，还是真福\Person{保禄}，他一受托\OriginalTerm{管家职分}{stewardship}，就\BibleQuote{没有与属血肉的人商量}\BibleRef{迦 1:16}？因为虽然他说：\BibleQuote{我连当宗徒的资格也没有}\BibleRef{格前 15:9}，然而，他既知道自己所领受的，也不忽视赐予者，便写道：\BibleQuote{我若不传福音，我就有祸了。}\BibleRef{格前 9:16} 但是，既然不传福音是\InnerQuote{有祸了}，那么，当他教导和宣讲\OriginalTerm{福音}{gospel}时，他所引领归主的人就成了他的\BibleQuote{喜乐和冠冕}\BibleRef{得前 2:19}。这就解释了为何这圣人是如此热切地传福音，甚至直到\Place{伊利里亚}，也不退缩前往\Place{罗马}\BibleRef{罗 1:15}，甚至远至\Place{西班牙}，为的是越劳苦，就越能获得更大的赏报。他夸口\BibleQuote{已打完美好的仗}\BibleRef{弟后 4:7}，并确信要获得那巨大的冠冕。因此，亲爱的\Person{德拉孔提}，你现在的行为是在效法谁呢？保禄，还是与他不同的人？至于我，我祈求你和我也能证明自己是效法所有圣人的人。\par

5. 或许有些人劝你躲藏起来，因为你已发誓，若当选就不接受这项职务。因为我听说他们正这样在你耳边聒噪，且自认为这是凭着良心行事。然而，如果他们真有良心，就应先敬畏那位将这职务托付给你的天主。或者，如果他们读过圣经，就不会劝你违背圣经了。因为如今他们连\Person{耶肋米亚}也该责备，并要控告伟大的\Person{梅瑟}，指责他们没有听从他们的劝告，而是敬畏天主，履行了自己的职务，且在履行先知职务上臻于圆满。原来，当他们接受使命和预言的恩宠时，也曾推辞。但后来他们却起了敬畏，没有轻看那位派遣他们的。所以，就算你口舌迟钝，说话结巴，仍要敬畏那位创造你的天主；或者，若你自称太年轻，不宜传道，仍要敬畏那位在你受造之前就已认识你的天主。再者，即使你已许下承诺（须知，他们是将对圣人所说的话视同誓言的），仍要看看\Person{耶肋米亚}：他也曾说过：\BibleQuote{我不再提念上主的名}\BibleRef{耶 20:9}，但后来他畏惧内心燃起的火，便没有照自己所说的去做，也没有好像受誓言束缚那样躲藏起来，反而敬畏那位将职务托付给他的，并完成了先知的召叫。亲爱的弟兄，难道你不知道\Person{约纳}也曾逃避，但遭遇了那事之后，就回来履行了先知的职务吗？\par

6. 因此，不要听从相反于此的劝告。因为主比我们自己更了解我们的情况，他知道将他的众教会托付给谁。因为一个人即使不配，也不应只顾以往的生活，而应履行他的职务，免得除了他本来的生活，还招致疏忽职守的诅咒。亲爱的\Person{德拉康提乌斯}，我问你：你既知道这些事，又是个明智之人，灵魂岂不感到刺痛吗？你岂不挂虑，恐怕那些托付给你的人中会有人丧亡吗？你的良心岂不犹如烈火焚烧吗？你岂不畏惧审判之日，那时现今给你出主意的人没有一个能在那里帮助你吗？因为每人都要为交在他手中的人交账。那埋藏银钱的人，他的推辞对他有什么益处呢？\Person{亚当}说：\BibleQuote{女人哄骗了我}\BibleRef{创 3:12}，这话对他又有什么益处呢？亲爱的\Person{德拉康提乌斯}，即使你确实软弱，你仍应当承担起职责，免得教会无人照管，敌人趁机，利用你的逃避来伤害她。你应当束起腰来，免得让我们独自作战；你应当与我们一同劳苦，以便将来也与众人一同领受赏报。\par

7. 所以，亲爱的，赶快行动吧，不要再迟延，也不要容忍那些阻拦你的人；只要记念那位赐下恩宠的，并到我们这些爱你的人这里来，我们给你的劝告是出自圣经的，好使你能由我们亲自祝圣，并当你管理教会时记念我们。因为你并不是唯一从隐修士中被选出来的人，也不是唯一管理过隐修院，或被隐修士们所爱的人。你知道，\Person{塞拉彼雍}不仅曾是隐修士，还管理着一大群隐修士；你也并非不知道，\Person{阿颇罗}曾是许多隐修士的父亲；你认识\Person{阿加通}，也不缺少对\Person{阿里斯通}的了解。你记得曾与\Person{塞拉彼雍}一同出外的\Person{阿摩尼乌斯}。也许你还听说过上\Place{特拜德}的\Person{穆伊图斯}，并能打听到\Place{拉托波利斯}的\Person{保禄}和其他许多人。然而，这些人在被选立时并没有推辞；他们以\Person{厄里叟}为榜样，知道\Person{厄里亚}的事迹，并学习了关于门徒和宗徒的一切，就承担起职责，没有轻视这职务，也没有辜负自己原来的身份，反而期待着劳苦的赏报，使自己精进，并引领他人向前。多少人因他们而远离偶像？多少人因他们的规劝而不再是邪魔的密友？他们领了多少仆人归向上主，使那些见到如此奇事的人惊叹不已？使一位少女度童贞生活，使一位少年人度节德生活，使一位拜偶像者认识基督，这岂不都是伟大的奇迹吗？\par

8. 因此，不要让隐修士们劝阻你，好像只有你一个是出自隐修士而被选立似的；也不要找借口说自己会退步。其实，你若效法\Person{保禄}，随从圣人们的行事，你反而可能更进步。因为你知道，那些像这样的人，当他们被委任为奥迹的管家时，更是奋力前进，向着他们\BibleQuote{崇高召叫的标竿}\BibleRef{斐 3:14}。\Person{保禄}是在被派遣去教导之后，才接受殉道，期待获得荣冠吗？\Person{伯多禄}是在宣讲福音，成为\BibleQuote{渔人的渔夫}\BibleRef{玛 4:19}之后，才宣认信仰吗？\Person{厄里亚}不是在完成他先知的职务之后，才被提升天吗？\Person{厄里叟}不是在舍弃一切，追随\Person{厄里亚}之后，才获得双倍的神恩吗？救主拣选门徒，不就是为了派遣他们作宗徒吗？\par

9. 所以，亲爱的\Person{德拉康提乌斯}，你应以这些人为榜样，不要听信人说主教的职务是犯罪的机会，也不要相信它会引发犯罪的诱惑。因为即使你身为主教，也仍可以如\Person{保禄}一样忍饥受渴。\BibleRef{斐 4:12}你也可以像\Person{弟茂德}那样不饮酒，\BibleRef{弟前 5:23}并像\Person{保禄}那样时常\BibleQuote{守斋}\BibleRef{格后 11:27}，好使你如此效法他守斋之后，能以你的言语饱饫他人，并在缺乏饮料而干渴时，以你的教导滋润他人。因此，不要让那些劝你的人提出这些借口。因为我们知道有守斋的主教，也有用饭的隐修士。我们知道有不饮酒的主教，也有饮酒的隐修士。我们知道有行奇迹的主教，也有不显奇迹的隐修士。许多主教甚至没有结婚，而隐修士却作了父亲；反过来说，我们也知道有些主教是子女的父亲，而有些隐修士却是\Quote{最成全的}。我们也知道有忍饥的圣职人员，也有守斋的隐修士。因为在后者的身份中守斋是可能的，在前者的身份中也不是不允许的。所以，一个人不论身在何处，都应热心努力；因为冠冕的赏赐不是凭地位，而是凭行为。\par

10. 因此，不要容忍那些提出相反意见的人。反而要赶紧，不要迟延；尤其因为圣节临近了；免得平信徒在没有你的情况下守节，而你给自己带来巨大的危险。因为，如果你不在，谁会向他们宣讲复活节讲道呢？如果你躲藏起来，谁会向他们宣告复活的大日子呢？如果你逃亡，谁会劝勉他们恰当地守节呢？啊，如果你出现，多少人会受益，如果你逃跑，多少人会受损！谁又会因此认为你好呢？他们为什么劝你不要接受主教职位，而他们自己却希望有司铎呢？因为如果你不好，就让他们不要与你交往。但如果他们知道你好，就不要嫉妒别人。因为，如果如他们所说，教导和管理是犯罪的机会，就让他们自己不要受教导，也不要拥有司铎，免得他们堕落，他们自己和教导他们的人都堕落。但不要理会这些世人的言辞，也不要容忍那些提出此类建议的人，正如我已经多次说过的。反而要赶紧转向主，好使你既能照顾他的羊群，也能记念我们。为此，我已吩咐我们亲爱的司铎\Person{希拉克斯}和读经员\Person{马克西穆斯}前去，也亲口嘱咐你，使你这便可以明白我写信时的心情，以及违抗教会命令的危险。\par

\SourceRecord{Translated by Archibald Robertson.；Nicene and Post-Nicene Fathers, Second Series；Vol. 4.；Edited by Philip Schaff and Henry Wace.；Buffalo, NY: Christian Literature Publishing Co.,；1892.；<http://www.newadvent.org/fathers/2806049.htm>.}

% structure: p0001 source=h1 target=section reason=page-title

\section{第五十封信}

\Emph{致\Person{路济弗尔}的第一封信。}\par

致我们的主、最亲爱的兄弟、主教兼\OriginalTerm{宣信者}{Confessor}\Person{路济弗尔}。\Person{亚大纳削}在主内问候。\par

托天主的福，我们身体安好，现在差派我们最亲爱的执事\Person{欧提克斯}，以便您至为虔诚的圣德，正如我们深切渴想的，能将您本人及与您同在者平安的消息告知我们。因为我们相信，正是藉着你们这些\OriginalTerm{宣信者}{Confessors}、天主之仆的生命，\OriginalTerm{天主教会}{Catholic Church}的状况得以更新；而异端者所企图撕裂的，我们的主\Person{耶稣基督}通过你们得以完全恢复。\par

因为，虽然\OriginalTerm{敌基督}{Antichrist}的前驱藉着这世俗的权力，已尽一切努力想要熄灭真理之灯，然而\OriginalTerm{天主性}{Deity}因着你们的\OriginalTerm{宣认}{confession}，却更清晰地显示出它的光明，使人不能不看见他们的诡诈。从前他们或许还能掩饰：如今他们被称为\OriginalTerm{敌基督者}{Antichrists}。因为，谁不诅咒他们，如同躲避污秽或蛇毒一般远离他们的共融呢？整个教会处处哀恸，每座城市都在呻吟，年迈的主教们正在流放中受苦，而异端者却伪装着，他们否认基督，使自己成为税吏，坐在教会中征收贡银。啊，魔鬼所发明的新种类的人与迫害，竟使用如此的残忍，甚至以圣职人员作邪恶的帮凶。但是，尽管他们如此行事，且极其骄傲与亵渎，你们的宣认、你们的虔诚和智慧，仍将是兄弟团体最大的安慰与慰藉。因为有人告知我们，您的圣德已经致书于\Person{君士坦提乌斯·奥古斯都}；我们愈加惊奇，您虽如同居于蝎子群中，却仍保持着精神的自由，为能以劝告、教导或纠正，引领那迷途者归向真理之光。因此，我请求，且所有宣信者与我一同请求，您惠允寄给我们一份抄本；好使众人不仅凭传闻，更凭书信，认识您精神的勇气，以及您信德的信心与坚定。与我同在的众位问候您的圣德。我问候所有与您同在的人。愿天主常保佑您平安无恙，并记念我们，最亲爱的主、天主真正的仆人。\par

\EditorialNote{收到这封信后，真福\Person{路济弗尔}寄出了他致\Person{君士坦提乌斯}的书；\Person{亚大纳削}读后，寄出了第五十一封信。}\par

\SourceRecord{Translated by Archibald Robertson.；Nicene and Post-Nicene Fathers, Second Series；Vol. 4.；Edited by Philip Schaff and Henry Wace.；Buffalo, NY: Christian Literature Publishing Co.,；1892.；<http://www.newadvent.org/fathers/2806050.htm>.}

% structure: p0001 source=h1 target=section reason=page-title

\section{第51封信}

\Emph{致\Person{路西弗}的第二封信}\par

致最荣耀的\OriginalTerm{主}{lord}、理应极其渴望的同主教\Person{路西弗}，\Person{亚大纳削}在主内问候。\par

虽然我相信你的圣善也已得到消息，关于基督的仇敌近日试图煽起的迫害，他们图谋我们的血，但我等最亲爱的使者能将此事详告你的虔敬。他们竟如此疯狂，借助士兵，不仅驱逐了城中的圣职界，还去到隐修士那里，向独修者们下了毒手。因此我也远避他处，免得款待我的人因我遭受他们的伤害。\OriginalTerm{亚略派}{Arians}连自己的灵魂也不顾惜，又岂会宽恕谁呢？或者，当他们坚持否认我们的主基督、天主的唯一圣子时，又怎能放弃他们可耻的行为呢？这是他们邪恶的根源；他们在这沙土的基础上建筑其行径的乖僻，正如我们在\BibleBook{}{圣咏}第十三篇所读到：\BibleQuote{愚妄的人心中说：没有天主。}紧接着又说：\BibleQuote{他们都丧尽天良，恣意作恶。}因此，那些否认天主子的犹太人，理应被称为\BibleQuote{邪恶的民族，罪孽沉重的人民，作恶者的后裔，没有法律的子女}\BibleRef{依 1:4}。为何是\Emph{没有法律}呢？——因为你们离弃了上主。如此，至福的保禄，当他不但开始相信天主之子，而且宣讲祂的神性时，写道：\BibleQuote{我自己毫无所愧}\BibleRef{格前 4:4}。因此，我们也愿按照你的信德宣认，持守宗徒的传统，并按神圣法律的诫命生活，好能在你左右，如今圣祖、先知、宗徒及殉道者正在欢欣踊跃的那一群体中寻得我们。所以，尽管\OriginalTerm{亚略的疯狂}{Arian madness}借助外来权势如此猖獗，以致我们的弟兄因他们的暴怒甚至无法自由见天日，但赖天主的恩佑，并因你的祈祷，我虽经历困苦危险，仍得以见到那位常为我带来必需品和你的圣善信函以及其他弟兄信函的信使。因此我们收到了你那最智慧敬虔的灵魂的著作，其中我们看见了宗徒的形象，先知的胆识，真理的教导，真信仰的教义，天国之路，殉道的荣耀，对抗\OriginalTerm{亚略异端}{Arian heresy}的凯旋，先辈们毫无受损的传统，以及教会秩序的正当规章。啊，真正的\Person{路西斐尔}，真如你的名字，带来真理之光，并把它放在灯台上，照亮众人。因为除了亚略派，有谁不能从你的教导中清楚看见真信仰与亚略派的污点呢。你强有力且令人钦佩地，如同光明自黑暗中，将真理从异端者的诡诈和虚伪中分离出来，捍卫了\OriginalTerm{天主教会}{Catholic Church}，证明了亚略派的论调不过是一种幻觉，并教导人那魔鬼的切齿是可鄙的。你对殉道的劝勉何其美好而可贵；你将为了基督、永生天主之子的死显示得何等值得渴求。你对来世和天上生活的爱是何等深厚。你似乎是救主的真实殿宇，祂居住在你内，借你口说出这些确切的话，并赋予了你的讲论如此恩宠。你从前已为众人所爱，如今众人心中对你充满如此热烈的感情，以致他们称你为我们时代的\Person{厄里亚}；这也不足为奇。因为若那些似乎中悦天主的人，被称为天主的子女，那么将这名字加之于先知们的伙伴——即\OriginalTerm{精修者}{Confessors}，而且特别加之于你，就更合宜了。相信我，\Person{路西斐尔}，说出这些话的不只是你，而是\OriginalTerm{圣神}{Holy Spirit}与你同在。从何处你获得了如此卓越的圣经记忆？从何处有了对圣经毫无损伤的领悟与理解？从何处构建起如此论述的条理？从何处你得到了对天国之路的这般劝勉，从何处得了对抗魔鬼的这般胆识，以及反驳异端的这般明证，若非圣神寓居在你内？因此，欢乐吧，因见到你已经在那里，你的先行者殉道者们所在之处，即在那天使的群体中。我们也欢喜，因有你作我们英勇、忍耐与自由的表率。至于你在著述中提到我名字的话，我羞于说什么，以免显得我在谄媚。然而我知道并相信，主自己已将一切知识启示给你圣善而虔敬的心灵，也必为这辛劳在天国中赐你赏报。你既是这样的人，我们恳求主，好使你能为我们祈求，求主大发慈悲，俯视\OriginalTerm{天主教会}{Catholic Church}，拯救祂所有的仆人脱离迫害者之手；好使那些因现世的恐惧而跌倒的人，也能最终得以重新站起来，回到义路上；他们离开了义路，正四处漂流，可怜的人们，不知道自己身处何种深坑。我特别请求，如果我言谈有何不妥，请你宽仁包容，因为从如此丰沛的泉源，我的愚拙未能汲取它本可汲取的份。至于我们的弟兄，我再次请求你包容我未能见到他们。因为真理本身为我作证，我切愿并渴望做到此事，深以不能为苦。因我的眼目不止流泪，我的心灵不止叹息，因为我们连弟兄们也不得一见。但天主为我作证，由于他们的迫害，我甚至连自己的父母也不能得见。因为亚略派还有什么事做不出来呢？他们监视道路，窥探出入城市的人，搜查船只，巡行旷野，搜查房屋，骚扰弟兄，扰乱众人。但感谢天主，这般行径使他们日益招惹众人的憎恶，也正应了你圣善对他们之称：\OriginalTerm{假基督}{Antichrist}的奴仆。而这些可怜虫，虽为人憎恶，却仍坚持他们的恶行，直到他们被判定其祖先\Person{法郎}那样的死亡。与我同在者问候你的虔敬。请代问候与你同在者。愿天主的恩宠护佑你，记念我们，永远蒙福；你堪当被称为天主的人，基督的仆人，宗徒的伙伴，弟兄们的安慰，真理之师，且在一切事上最被人渴慕。\par

\SourceRecord{Translated by Archibald Robertson.；Nicene and Post-Nicene Fathers, Second Series；Vol. 4.；Edited by Philip Schaff and Henry Wace.；Buffalo, NY: Christian Literature Publishing Co.,；1892.；<http://www.newadvent.org/fathers/2806051.htm>.}

% structure: p0001 source=h1 target=section reason=page-title

\section{第52封信}

\Emph{致隐修士的第一封信。（写于358-360年。）}\par

1. 致生活在各处的度\OriginalTerm{隐修生活}{monastic life}者，他们在天主的信仰上站稳，在基督内被圣化，他们曾说：\BibleQuote{看，我们舍弃了一切，跟随了祢}\BibleRef{玛 19:27}。我至爱并渴念的弟兄们，致以主内的最热切问候。\par

1. 为回应你们多次向我恳切提出的深情请求，我撰写了关于我们自己以及教会所经受之苦痛的简短记述，竭尽所能驳斥亚略疯子的可咒\OriginalTerm{异端}{heresy}，并证明它完全背离了真理。我想有必要向你们的虔诚说明，撰写这些内容花费了我多少辛劳，好让你们能够由此理解真福宗徒所言是何等真切：\BibleQuote{啊，天主的智慧和知识，何其丰盛深邃}\BibleRef{罗 11:33}；并且能够善意地宽容像我这样天生软弱的人。因为我越想写作，并努力强迫自己理解\OriginalTerm{圣言的神性}{Divinity of the Word}，那知识就越发远离我；我本以为领会了多少，就越发觉自己未能领会。此外，我甚至无法用文字表达自己以为理解了的内容；我所写的，根本无法与我心中存有的那不完全的真理之影相称。\par

2. 因此，我想起\BibleBook{训道}{篇}所载：\BibleQuote{我曾说：我要有智慧，但智慧离我很远；那遥远且深奥已极的，谁能找着呢？}\BibleRef{训 7:23}以及\BibleBook{圣咏}{集}中所言：\BibleQuote{祢的知识奇妙莫测，是我不能企及的，崇高，是我不能达到的}；还有\Person{撒罗满}说：\BibleQuote{将事情隐藏，是天主的光荣}\BibleRef{箴 25:2}；我屡次决意停笔不再写作；请相信我，我确实如此。但为免令你们失望，或因我的沉默而使那些向你们求问、又好争辩的人导入\OriginalTerm{不虔}{impiety}，我勉强自己简要写出，即现在送呈你们虔诚之处的这封信。因为，尽管由于肉身的软弱，目前我们尚无法完全领悟真理，但如训道者本人所言，仍可能看透不虔者的狂妄，并发现后说，它\BibleQuote{比死亡更苦}\BibleRef{训 7:26}。为此缘由，我既然能看透并能发现它，便写了这信，知道对于信者而言，揭露不虔便足以说明虔诚在于何处。因为尽管无法理解天主是什么，但仍可能说出祂不是什么。我们知道，祂不像人；我们也不可想象有任何\OriginalTerm{受造的本性}{originated nature}存于祂内。同样，关于天主子，尽管我们本性上远远无法理解祂，但仍然可能且容易驳斥异端者关于祂的言论，并说，天主子并非如此；我们甚至不可在心里设想他们关于祂\OriginalTerm{天主性}{Godhead}所讲的那些事，更不可用口唇说出来。\par

3. 因此，我已尽己所能写了；而你们，至爱者，接受这些通信，不要视其为对\OriginalTerm{圣言天主性}{Godhead of the Word}的完备阐述，而仅是反驳基督之敌的不虔，并为渴求者提供导向在基督内虔诚且健全信仰的提示。如果这些文字有任何缺欠（我认为它们在各方面都有缺欠），请以纯洁的\OriginalTerm{良心}{conscience}予以宽恕，只欣然接纳我支持\OriginalTerm{敬虔}{godliness}之善意的胆量。对于彻底谴责亚略异端，你们知道主在\Person{亚略}之死中所施的审判便已足够，对此你们已由他人告知。\BibleQuote{因为圣洁天主所定旨的，谁能破坏呢？}\BibleRef{依 14:27}主所谴责的，谁能称为义呢？既然给出了这样的标记，现在谁不承认这异端为天主所憎恶，不论它有怎样的人作其庇护者？你们读过此记述后，请为我祈祷，并彼此劝勉这样做。并立即将它送回给我，不要让任何人抄录，也不要为自己誊写。但要像精明的兑钱商那样，满足于阅读；若渴望，可反复阅读。因为我们这些饶舌者和个人的作品，落到后来的人手中是不安全的。在爱中彼此问候，也问候所有因虔诚和信仰来到你们这里的人。因为如宗徒所说：\BibleQuote{若有人不爱主，当受诅咒。愿主耶稣基督的恩宠与你们同在。阿们。}\BibleRef{格前 16:22}\par

\SourceRecord{Translated by Archibald Robertson.；Nicene and Post-Nicene Fathers, Second Series；Vol. 4.；Edited by Philip Schaff and Henry Wace.；Buffalo, NY: Christian Literature Publishing Co.,；1892.；<http://www.newadvent.org/fathers/2806052.htm>.}

% structure: p0001 source=h1 target=section reason=page-title

\section{\WorkTitle{第53封信}}

\Emph{\WorkTitle{致隐修士的第二封信}。}\par

\Person{亚大纳削}，\Place{亚历山大里亚}总主教，致\OriginalTerm{独修士}{Solitaries}们。\par

\Person{亚大纳削}向那些度独修生活、在天主内信德坚固的至爱弟兄们，致以主内的问候。\par

我感谢主，因为他赐予你们在他内信从的恩宠，好使你们也能同诸圣获得永生。但是，因为有些人附从\Person{亚略}，并在各隐修院之间游走，没有别的目的，只是假借拜访你们之名，并从我们这里回去之后，欺骗那些纯朴的人；而另有一些人，虽然声称自己并不附从\Person{亚略}，却自甘妥协，与他的党羽一同崇拜；因着一些极其真诚的弟兄们的请求，我不得不立即写信，好使你们能以忠诚无伪的心，保守天主的恩宠在你们内所成就的虔诚信德，而不致给弟兄们立下恶表。因为当有人看见你们这些在基督内的忠信者，竟与这样的人交往共融\EditorialNote{或一同参与崇拜}，他们必会认为这是无关紧要的事，并陷入不虔敬的泥淖。因此，亲爱的，为避免这样的事发生，请你们远避那些持守不虔敬\EditorialNote{亚略}的人，更要躲避那些人，他们虽然装作不附从\Person{亚略}，却与不虔敬者一同崇拜。我们特别有责任逃避那些我们深恶痛绝其主张的人的共融。\EditorialNote{因此，如果有人来到你们那里，正如\Person{真福若望}所说：\BibleRef{若二 1:10}，带着纯正的教义，你们应对他说：愿你平安，并接纳这人如同弟兄。}但若有人声称宣认纯正的信仰，却显然与那些人保持共融，就要劝他戒绝这种共融；如果他答允这样做，就待他如同弟兄；如果他固执好争，就要避开他。\EditorialNote{我本可以大幅增长这封信的篇幅，从神圣的圣经中援引这项教导的纲要。然而，你们既是智慧人，能预先领会写信者的意思，又更专心于舍己克修，足以教导他人，所以，我因着信赖，就口述了这封短信，如同一位爱友写给另一些爱友，深信你们照现在的生活，必会保持纯洁真诚的信德，并且那些人会因看到你们不与他们一同崇拜而获益，他们害怕自己被视为不虔敬者，并与这样的人同流合污。}\par

\SourceRecord{Translated by Archibald Robertson.；Nicene and Post-Nicene Fathers, Second Series；Vol. 4.；Edited by Philip Schaff and Henry Wace.；Buffalo, NY: Christian Literature Publishing Co.,；1892.；<http://www.newadvent.org/fathers/2806053.htm>.}

% structure: p0001 source=h1 target=section reason=page-title

\section{第五十四封信}

\Emph{致\Person{塞拉皮翁}，论\Person{亚略}之死}\par

\Person{亚大纳削}致\Person{塞拉皮翁}，兄弟及\OriginalTerm{同工}{fellow-minister}，在主内祝安。\par

我阅读了您虔敬的信函，信中您请求我向您讲述与我时代相关的有关我本人的事件，并陈述那\OriginalTerm{亚略派}{Arians}最邪恶的异端——我因此异端而遭受这些苦难——以及\Person{亚略}死亡的情形。对您的三项请求，我欣然应允了其中两项，并已将我所写与隐修士们的书信寄予您的\OriginalTerm{虔敬}{Godliness}；从中您可以了解到我本人的经历以及这异端的原委。至于另一件事，我指的是他的死，我心中犹豫了许久，唯恐有人以为我因那人之死而幸灾乐祸。然而，既然你们当中关于这异端的争论，最终归结为这个问题：\Person{亚略}究竟是否先与教会\OriginalTerm{共融}{communion}而后死去？我因此不得不愿对此死况作一交代，认为如此便可平息这疑问，又考虑到借此披露，亦可同时堵住那些好争辩之人的口。因为我想，当与他死亡相关的奇异情形被人知晓后，就连那些先前质疑的人，也不敢再怀疑亚略异端在天主眼中是可憎的。\par

2. 他死时我并不在\Place{君士坦丁堡}，但\Person{马卡里}\OriginalTerm{长老}{Presbyter}在场，我从他那里听到了叙述。\Person{亚略}因\Person{欧瑟比}及其同党的运作，受\Person{君士坦丁}皇帝邀请；当他进见时，皇帝问他，是否持守天主教会的信德？他起誓宣称自己持守正信，并书面递交了一份自述信仰的陈词，但隐瞒了那些致使他被\Person{亚历山大}主教逐出教会的要点，只巧言引用圣经中的词句。因此，当他誓称自己不持守\Person{亚历山大}绝罚他所依据的那些观点后，皇帝便打发他走，说：\Quote{如果你的信德正确，你发誓是好的；但如果你信德邪恶，却发了誓，愿天主按你的誓言审判你。}当他如此从皇帝面前出来后，\Person{欧瑟比}及其同党，以其惯常的强横，想要领他进入教会。但那位可敬记念的\Place{君士坦丁堡}主教\Person{亚历山大}，抵制他们，说异端的创立者不应获准\OriginalTerm{共融}{communion}；于是\Person{欧瑟比}及其同党威胁说，宣称：\Quote{我们既已违你心愿，让皇帝邀请了他，同样明天，即便违你所愿，\Person{亚略}也必将在这教堂内与我们共融。}他们说这话时，正是安息日。\par

3. 主教\Person{亚历山大}听闻此言，极其忧苦，进入圣堂，向天主伸出双手，哀哭自己；并在\OriginalTerm{司祭所}{chancel}内俯伏于地，躺卧在石板地上祈祷。\Person{马卡里}当时也在场，与他一同祈祷，并听见了他的话语。他恳求这两件事，说：\Quote{若明天\Person{亚略}被领入共融，就让祢的仆人离去吧，不要将虔敬者与不虔敬者一同毁灭；但若祢怜悯祢的教会（我知道祢必怜悯），就垂顾\Person{欧瑟比}及其同党的话语，不要将祢的产业交予毁灭和羞辱\BibleRef{岳 2:17}，并除去\Person{亚略}，免得他若进入教会，那异端似乎也随他而入，从此以后，\OriginalTerm{不虔敬}{impiety}反被视作\OriginalTerm{虔敬}{piety}。}主教这样祈祷之后，便在深重的焦虑中退去；这时发生了一件奇妙非凡的事。正当\Person{欧瑟比}及其同党威胁，主教祈祷之时，\Person{亚略}因对\Person{欧瑟比}及其同党信心满满，且语多狂妄，在这时因自然需要而退到一边，突然，按圣经的话说，\BibleQuote{倒头直下，在中间崩裂}\BibleRef{宗 1:18}，当场仆倒在地，即刻气绝身亡，同时失去了共融与生命。\par

4. \Person{亚略}的下场便是如此：\Person{欧瑟比}及其同党，羞愧难当，埋葬了他们的同谋；而真福\Person{亚历山大}，在教会的欢腾之中，虔诚而正统地举行了\OriginalTerm{圣体圣事}{Communion}，与所有弟兄一同祈祷，极力光荣天主，并非因他的死而幸灾乐祸（天主不许！），因为\BibleQuote{按照命定，人人都有一死}\BibleRef{希 9:27}，而是因为这事以一种超越人意料的奇妙方式被显明出来。因为主亲自在\Person{欧瑟比}及其同党的威胁，与\Person{亚历山大}的祈祷之间施行审判，谴责了亚略异端，表明它不配与教会共融，并向众人显明：尽管它获有皇帝及全人类的支持，却为教会自身所谴责。因此，\OriginalTerm{亚略派}{Arian}疯子的敌基督团伙，已显示出不为天主所悦、且是邪恶的；许多先前受其欺骗的人，都改变了看法。因为正是那受他们亵渎的主自己，谴责了这起来反抗祂的异端，并再次表明：无论\Person{君士坦提乌斯}皇帝如今如何为这异端对主教们施以暴虐，它仍被排斥于教会共融之外，与\OriginalTerm{天国}{kingdom of heaven}无分。因此，你们当中所兴起的疑问，从今以后也当平息（因为你们之间曾有如此约定）；谁也不要附和这异端，连那些受骗者也当悔改。因为主所谴责的，谁还敢接受呢？凡支持祂所\OriginalTerm{绝罚}{excommunicate}之物的，岂不犯下大不敬之罪，且公然成为基督的仇敌吗？\par

这已足以驳倒好争辩之人；因此，请将它读给先前提出这问题的人，也读给那封致隐修士们的、反对\OriginalTerm{异端}{heresy}的简短书信，好让他们由此受到引导，更严厉地谴责\OriginalTerm{亚略派狂人}{Arian madmen}的不虔和邪恶。然而，不要同意将这些抄本给予任何人，你自己也不要抄写（我同样已向隐修士们表示过）；但是作为真诚的朋友，如果我写的内容有什么欠缺，就请补充，并立即把它们送回给我。你会从我写给弟兄们的信中得知，我写这封信付出了多大的辛劳，同时也能认识到，一个人的私人著述被公开是不安全的（尤其当它们涉及最高和最首要的信理时），原因在于——因软弱或语言晦涩而表达不完美的内容，恐会伤害读者。因为大多数人并不考虑信德或作者的意图，而是出于嫉妒或争辩之心，按照自己预先形成的意见，随意接受所写的内容，并随己意曲解它。但愿主使有关我们的主\Person{耶稣基督}的真理和健全信德在所有人中间盛行，尤其是在你将要读给他们听的那些人中间。阿们。\par

\SourceRecord{Translated by Archibald Robertson.；Nicene and Post-Nicene Fathers, Second Series；Vol. 4.；Edited by Philip Schaff and Henry Wace.；Buffalo, NY: Christian Literature Publishing Co.,；1892.；<http://www.newadvent.org/fathers/2806054.htm>.}

% structure: p0001 source=h1 target=section reason=page-title

\section{第55封信}

\Emph{致鲁菲尼亚努斯书}\par

致我们的主、儿子、最亲爱的同工\Person{鲁菲尼亚努斯}。\Person{亚大纳削}在主内问候。\par

你写信给我，正如一位可爱的儿子对父亲写信时所当做的：因此，当你以书信亲近我时，我最亲爱的\Person{鲁菲尼亚努斯}，我便拥抱了你。虽则我也许可在信首、信中或信尾以儿子待你，我却加以抑制，以免我的赞许和见证借着文字显示出来。因为你就是我的书信，如经上所载\BibleRef{格后 3:2}，铭刻在心，为人所共知共读。你既如此，当深信，且要深信。我写信给你，并邀请你也写信。因为你这样做，便给了我最大的喜乐。然而，你既以一种高尚且合乎教会的精神，一如你的虔敬所当然，向我询问那些受势所逼而被掳去、却未受错谬所败坏的人，并希望我写信告知对此事，无论在会议上或其他处，已作何决议；你当知道，最亲爱的主，首先，当暴力停止后，曾召开了一次会议，有来自外地的众主教参加；同时，在\Place{希腊}的我们的同工，以及在\Place{西班牙}和\Place{高卢}的同工，也举行了其他会议；并且这里和各地都得出了同样的决定，就是：对于那些曾跌倒并作了不虔敬首领的人，只要他们悔改，可得宽赦，但不授予圣职的品位；至于那些并非蓄意行不虔敬之事，而是势不得已、遭到强暴而被掳去的人，则应不仅获得宽赦，而且可保留圣职的品位：尤因他们提出了情有可原的辩护，而所发生的事似乎出于某种特殊的理由。因为他们向我们保证，他们并未投向不虔敬；而是唯恐某些极不虔敬的人当选并败坏诸教会，他们情愿屈从强暴、忍受重担，而不愿失去民众。但他们这样说，我们觉得似有道理；因为他们引述\Person{梅瑟}的哥哥\Person{亚郎}作为理由，\Person{亚郎}在旷野容忍了百姓的过犯；他如此行，乃因怕百姓返回\Place{埃及}并陷于偶像崇拜中。因为这种看法确有理由：如果他们留在旷野，他们可能得以戒除不虔敬；但若去往\Place{埃及}，他们必遭败坏，且使不虔敬在他们中间加剧。为此之故，那些曾受蒙骗并遭受强暴的人，既蒙宽赦，便被准许保留圣职品位。我将此告知你的虔敬，深信你必接纳所决议的，且不指控那些与会者懈怠，一如我所说的。但务请你读给你手下的圣职人员和平信徒听，好使他们得知此事，而不责怪你对这些人抱持如此态度。因为由我写信并不合宜，既然你的虔敬能自行通知，并向他们宣布我们的心意，且执行其余该办的事。感谢主，祂已如同\BibleRef{格前 1:5}所载，以一切口才和一切知识充满你。因此，那些公开悔改的人，应当指名弃绝\Person{欧多克西乌斯}和\Person{欧佐伊乌斯}的错谬。因为他们仍在亵渎，且写道祂是受造物，他们是\OriginalTerm{亚略异端}{Arian heresy}的魁首。但他们当宣认在\Place{尼西亚}由教父们所宣认的信德，并且不得将其他任何会议置于该会议之上。请代问候你们那里的弟兄团体。我们这里的弟兄团体在主内问候你。\par

\SourceRecord{Translated by Archibald Robertson.；Nicene and Post-Nicene Fathers, Second Series；Vol. 4.；Edited by Philip Schaff and Henry Wace.；Buffalo, NY: Christian Literature Publishing Co.,；1892.；<http://www.newadvent.org/fathers/2806055.htm>.}

% structure: p0001 source=h1 target=section reason=page-title

\section{书信第五十六}

致约维安皇帝论信仰\par

\Emph{约维安皇帝致最神圣的亚历山大大主教亚太纳西的书信副本：}\par

致最虔诚而亲密天主的人，亚太纳西，约维安致候：我们极其钦佩您最可敬的生活之成就、您与万有的天主之相似、以及您对我们救主基督的爱戴，因此我们接纳您，最可敬的主教。您既未因一切劳苦而退缩，也未因迫害者的恐惧而退缩，视危险与刀剑的威胁如粪土，紧握您所珍视的正统信仰之舵，至今仍为真理而奋战，并继续向全体信众展现您为楷模与德行之榜样——我皇陛下召回您，并盼望您重返救恩教导之职。请回到神圣的教会，牧养天主的子民，并热诚为我们的仁政向天主献上祈祷。因为我们知道，因您的恳求，我们及所有与我们一同持守（基督徒信仰）的人，将从至高天主获得极大的助佑。\par

1. 渴望学习并切慕属天之事，正合于一位虔诚的皇帝；因为如此您必将真正拥有\Quote{您的心}也\Quote{在天主手中}。既然您的虔敬希望从我们学习天主教会信仰，我们为此事感谢主，首先劝告提醒您的虔敬父老们在尼西亚所宣认的信仰。因为某些人藐视这信仰，多方图谋反对我们，因为我们不愿附和亚略异端，他们便成了天主教会内异端与分裂的制造者。因为对主真实而虔诚的信仰已向众人显明，从圣经中\Quote{被知晓并被诵读}。圣徒们在此信仰中得以成全并殉道，如今安息于主内；这信仰原可永远保存无损，若非某些异端者的邪恶胆敢篡改它。因为某个亚略及其同伙企图败坏它，并引入不敬虔来代替它，声称天主之子来自无有，是受造物、被造之物且可改变。但这些话欺骗了许多人，甚至\Quote{那些看似有分量的}也被他们的亵渎所迷惑。然而，如前所述，我们神圣的父老们及时聚集于尼西亚大公会议，谴责了他们，并以文字宣认了天主教会信仰，以致这信仰到处被宣讲，异端所燃起的邪说得以熄灭。这信仰于是处处在各教会中被真诚地知晓并宣讲。但如今，某些人想要复兴亚略异端，竟胆敢藐视父老们在尼西亚所宣认的这信仰，并表面承认实则否认，曲解\Quote{\OriginalTerm{同体}{Coessential}}，又擅自亵渎圣神，宣称祂是受造物、是由圣子所造之物，我们鉴于这种亵渎对百姓造成的损害，有义务将尼西亚所宣认的信仰呈递您的虔敬，以便您的虔诚能知晓那准确记载的内容，以及那些悖逆教导的人错得何等之远。\par

2. 因为要知道，最虔诚的奥古斯都，这些事从古时就已宣讲，这信仰是聚集于尼西亚的父老们所宣认的；各地所有教会都赞同它，包括西班牙、不列颠、高卢、全意大利、达尔马提亚、达契亚、默西亚、马其顿、全希腊、全非洲、撒丁、塞浦路斯、克里特，以及旁非利亚、吕西亚、以扫利亚、埃及、利比亚、本都、卡帕多细亚、和我们邻近的地区，以及东方的教会，除了少数附和亚略的人。因为我们既从经验得知上述所有地方的意见，也拥有书信。最虔诚的奥古斯都，您知道，即使有少数人反对这信仰，他们也不能提出抗辩，因为全世界都持守宗徒的信仰。因为他们长期受亚略异端感染，如今更顽固地反对真理。为使您的虔敬得知，虽然您已知道，我们仍急于附上主教们在尼西亚所宣认的信仰。那么，父老们在尼西亚所宣认的信仰如下：——\par

3. 我们信……等等，等等。\par

4. 奥古斯都，所有人都必须持守这信仰，视之为神圣和宗徒的，绝不可用巧辩和争辩之词动摇它，这正是亚略狂人所做的，他们说天主之子来自无有，且曾有一时祂不存在，祂是受造的、被造的和可改变的。因为为此，如上所述，尼西亚大公会议谴责了这样的异端，却宣认了真理的信仰。因为他们不仅说圣子相似圣父，免得祂被相信只是相似天主，而不是从天主而来的真天主；他们写下了\Quote{\OriginalTerm{同体}{Coessential}}，这是真正而真实的圣子所特有的，真正而自然地来自圣父。他们也没有使圣神与圣父和圣子分离，反而将祂与圣父和圣子同受光荣，在至圣天主圣三的同一信仰中，因为在至圣天主圣三中也有一个天主性。\par

% structure: p0009 source=h2 target=subsection reason=relative-heading-rank; base=h2

\subsection{附录}

路求、贝尔尼基安和某些其他亚略派人士在安提约基雅向约维安皇帝提交的反对亚历山大大主教亚太纳西的请愿书。\par

皇帝出发前往军营时，他们在罗马门提交的第一份请愿书。\par

愿陛下、大王与陛下虔敬之心俯听我等。\Emph{皇帝：}\Quote{你们是谁，来自何方？}\Emph{亚略派信徒：}\Quote{基督徒，我主。}\Emph{皇帝：}\Quote{来自何方？来自哪座城？}\Emph{亚略派信徒：}\Quote{\Place{亚历山大}。}——\Emph{皇帝：}\Quote{你们想要什么？}\Emph{亚略派信徒：}\Quote{愿陛下、大王赐予我等一位主教。}\Emph{皇帝：}\Quote{我已下令先前你们所拥有的那位\Person{阿塔纳修}就任该教座。}\Emph{亚略派信徒：}\Quote{愿陛下明鉴：他多年流放，且身负指控。}突然一名士兵愤然答道：\Quote{愿陛下垂询他们是何人、来自何方，因为这些人是\Place{卡帕多细亚}的残渣余孽，是那亵圣的\Person{乔治}的残余，他使此城及天下荒芜。}皇帝听闻，便策马离去，前往军营。\par

亚略派信徒的第二项请愿。\par

\Quote{我等对\Person{阿塔纳修}握有指控与确凿证据，因他在十年、二十年前曾被可记念的\Person{君士坦丁}与\Person{君士坦提乌斯}剥夺职位，并在最虔诚、最富哲学与福乐之心的\Person{尤利安}治下遭受流放。}\Emph{皇帝：}\Quote{十年、二十、三十年前的指控如今已过时效。不要再向我提及\Person{阿塔纳修}，因我知道他为何被指控、如何被流放。}\par

亚略派信徒的第三项请愿。\par

\Quote{如今我们又有某些针对\Person{阿塔纳修}的新指控。}\Emph{皇帝：}\Quote{案件的是非曲直不会因人数众多与喧嚷而显明；你们从你们中选出两人，再从大多数一方选出另外两人，因我不能逐个回答你们。}\Emph{多数方：}\Quote{这些人是亵圣的\Person{乔治}的残余，他毁坏了我们的省份，不许任何谋士居住于城中。}\Emph{亚略派信徒：}\Quote{愿陛下垂允，任何您所愿意之人，除了\Person{阿塔纳修}。}\Emph{皇帝：}\Quote{我告诉你们，\Person{阿塔纳修}的案件已经了结。}（然后愤怒地）\Quote{\OriginalTerm{打}{feri}！\OriginalTerm{打}{feri}！}\Emph{亚略派信徒：}\Quote{愿陛下垂允，若你遣走\Person{阿塔纳修}，我们的城便毁了，无人与他集会。}\Emph{皇帝：}\Quote{然而我费心查明了，他持守正确见解，是正统派，教导有方。}\Emph{亚略派信徒：}\Quote{他口中所言正确，但心怀诡诈。}\Emph{皇帝：}\Quote{够了，你们已为他作证：他言说正确，教导有方；但若他舌讲正道而心怀恶念，那在天主面前自有关涉。因为我们是人，听见所说的；但心中所存，天主知晓。}\Emph{亚略派信徒：}\Quote{请授权我们彼此共融。}\Emph{皇帝：}\Quote{怎么，谁阻止你们？}\Emph{亚略派信徒：}\Quote{愿陛下垂允，他声称我们是分党派者和教条主义者。}\Emph{皇帝：}\Quote{那是他的职责，也是所有正确教导者的职责。}\Emph{亚略派信徒：}\Quote{愿陛下垂允，我等不能容忍此人，他已夺去了教会的土地。}\Emph{皇帝：}\Quote{哦，那么你们来此是为财产，而非为信仰}——接着他补充道——\Quote{走吧，保持和平。}他再次对亚略派信徒说：\Quote{你们到教会去，明天你们有共融礼；礼成之后，这里有主教们，还有\Person{内梅西努斯}在此，你们各人可按自己所信的签字；\Person{阿塔纳修}也在此；凡不明白信仰之言者，可从\Person{阿塔纳修}学习。你们有明天和后日，因我要去军营。}一位属于犬儒学派的律师向皇帝请愿：\Quote{愿陛下垂允，因主教\Person{阿塔纳修}之故，总收税官查封了我的房屋。}\Emph{皇帝：}\Quote{若总收税官查封了你的房屋，这与\Person{阿塔纳修}何干？}另一位律师\Person{帕塔拉斯}说：\Quote{我对\Person{阿塔纳修}有控诉。}\Emph{皇帝：}\Quote{你是异教徒，与基督徒有何相干？}但多数方中有些\Place{安提约基雅}人把\Person{路济乌}带到皇帝面前说：\Quote{愿陛下与大王垂允，看看他们想立为主教的人！}\par

\Emph{另一项在宫殿门廊处代表\Person{路济乌}提出的请愿：}——\Quote{愿陛下垂听我言。}皇帝停住说：\Quote{我问你，\Person{路济乌}，你是如何来此的，走海路还是陆路？}\Emph{\Person{路济乌}：}\Quote{愿陛下垂允，走海路。}\Emph{皇帝：}\Quote{好吧，\Person{路济乌}，愿世界之天主、光辉的太阳和月亮，向那些与你同航的人发怒，因为他们没有把你抛入海中；愿那艘船再不得顺风，风暴中载着乘客时也找不到港湾。}透过\Person{尤西乌}，不信的亚略派信徒请求\Person{普罗巴图斯}及其同伴——即\Person{尤西比乌}与\Person{巴迪奥}的继任者，身为宦官——使他们获准觐见。皇帝得知此事，便对宦官施以酷刑，并说：\Quote{若有人想对基督徒提出请愿，这便是他的下场。}于是皇帝遣散了他们。\par

\SourceRecord{Translated by Archibald Robertson.；Nicene and Post-Nicene Fathers, Second Series；Vol. 4.；Edited by Philip Schaff and Henry Wace.；Buffalo, NY: Christian Literature Publishing Co.,；1892.；<http://www.newadvent.org/fathers/2806056.htm>.}

% structure: p0001 source=h1 target=section reason=page-title

\section{第57封信}

致\Person{奥西斯乌斯}的第一封信\par

\Quote{在那里住了几天之后，他向院长\Person{德奥多鲁斯}说：\InnerQuote{逾越节近了，你照你的惯例去探访弟兄们吧；至于我，主怎样安排我，我便怎样做。} 于是，他拥抱了他，打发他离去，并借他带给院长\Person{奥西斯乌斯}和弟兄们一封信，内容如下：}——\par

我见到了你们的同工、弟兄们的父亲\Person{德奥多鲁斯}，在他身上见到了我们父亲\Person{帕科米乌斯}的导师。我因见到教会的子女而欢欣，他们的临在使我喜悦。但主是他们的酬报者。当\Person{德奥多鲁斯}要离开我回到你们那里时，他对我说：\Quote{请记念我。} 我对他说：\BibleQuote{耶路撒冷！我若不怀念你，愿我的右手枯萎！我若不怀念你，愿我的舌头紧贴上颚！}\BibleRef{咏 137:5-6}\par

\SourceRecord{Translated by Archibald Robertson.；Nicene and Post-Nicene Fathers, Second Series；Vol. 4.；Edited by Philip Schaff and Henry Wace.；Buffalo, NY: Christian Literature Publishing Co.,；1892.；<http://www.newadvent.org/fathers/2806057.htm>.}

% structure: p0001 source=h1 target=section reason=page-title

\section{第58封书信}

\Emph{致\Person{奥西修斯}的第二封信}。\par

\Quote{但是，至圣的总主教\Person{亚大纳削}听说我们的父亲\Person{狄奥多尔}的消息后，深感忧伤，便致信院长\Person{奥西修斯}和弟兄们，为他的逝世安慰他们，内容如下：}——\par

\Person{亚大纳削}致院长、隐修士之父\Person{奥西修斯}，及所有与他一同度独居生活、在天主内坚定信德、在主内最为想念的亲爱的弟兄：问候。\par

我已听闻有福的\Person{狄奥多尔}逝世的消息，这消息令我深感焦虑，因为我知道他对你们的价值。若非狄奥多尔，我定会泪流满面地对你们说许多话，思考死后之事。但既是你们和我都认识的狄奥多尔，我在信中还需说什么呢？唯有\BibleQuote{有福的}狄奥多尔，\BibleQuote{他不从恶人的计谋}。但若\BibleQuote{敬畏主的人是有福的}，我们现在大可放心地称他是有福的，确信他已抵达如避风港一般、无忧无虑的生活。但愿我们每人也都如此；但愿我们每人在奔跑中都能抵达；但愿我们每人航行途中都能将自己的小船停泊在那无风浪的避风港，好能与先辈们一同安息，并能说：\BibleQuote{我要住在这里，因为是我所喜悦的。}因此，亲爱的、最为想念的弟兄们，不要为\Person{狄奥多尔}哭泣，因为他\BibleQuote{不是死了，而是睡着了}\BibleRef{玛 9:24}。愿无人因想起他而哭泣，却要效法他的生活。因为不该为前往无忧之地的人忧伤。我写这些给你们大家；但尤其对你，亲爱的、最为想念的\Person{奥西修斯}，好叫他现在安眠之后，你能挑起全部担子，在弟兄中取代他的位置。当他活着时，你们二人如同一人；一人不在时，两人的工作仍继续进行；两人都在时，你们如同一人，向亲爱的弟兄们宣讲有益之事。所以，就这样行事吧，并写信告诉我你和弟兄团体的平安。我劝你们众人一同祈祷，求主赐予教会更多的和平。因为我们现在欢庆过节，既过了复活节又过了五旬节，并喜乐于主的恩惠。我写信给你们众人。请问候每一位敬畏主的人。与我同在的也问候你们。我祈祷你们在主内安好，亲爱的、极为想念的弟兄们。\par

\SourceRecord{Translated by Archibald Robertson.；Nicene and Post-Nicene Fathers, Second Series；Vol. 4.；Edited by Philip Schaff and Henry Wace.；Buffalo, NY: Christian Literature Publishing Co.,；1892.；<http://www.newadvent.org/fathers/2806058.htm>.}

% structure: p0001 source=h1 target=section reason=page-title

\section{第59封信}

\Emph{致\Person{埃比克泰德}}\par

致我主、亲爱的兄弟及最渴念的同工\Person{埃比克泰德}，\Person{亚大纳削}在主内致候。我曾以为，一切异端者的种种妄谈，无论他们何其众多，已为在\Place{尼西亚}所举行的会议所遏止。因为教父们依据\WorkTitle{圣经}所宣认的信德，本身即足以一举推翻一切不虔敬，并稳固在基督内的宗教信仰。为此之故，现今在各地召开的不同会议——在\Place{高卢}、\Place{西班牙}及伟大的\Place{罗马}——所有与会者，如同被同一\OriginalTerm{圣神}{Holy Spirit}所推动，一致地绝罚了那些仍暗中附和\Person{阿黎乌}的人，即\Place{米兰}的\Person{奥森提乌}、\Person{乌尔撒奇}、\Person{瓦伦斯}及\Place{潘诺尼亚}的\Person{加约}。他们且致书各处，申明鉴于上述诸人正设法假借会议之名来支持己方，在天主教会内，除\Place{尼西亚}所举行的会议外，不应援引任何其他会议，因为那是战胜一切异端的凯旋纪念，尤其战胜阿黎乌派的，而该异端正是当初促使该会议召开的主要原因。那么，这一切之后，为何仍有人试图激起怀疑或问题呢？若他们属于阿黎乌派，则对那为反对他们而制定的信辞吹毛求疵，也就不足为奇，正如外邦人听到\BibleQuote{外邦人的偶像不过是金银，是人手所造的}时，认为神圣十字架的教理是愚妄。但若那些试图通过提问来重新质疑一切的人，属于那些自以为信仰正确、并热爱教父们所宣告的道理的人，那么他们简直是在行先知所描述的事：将浑浊的混乱给邻人喝，并为无益的言词争斗，除了颠覆愚直者外，毫无益处。\par

2. 我写这封信，是在读了你的虔敬提交来的备忘录之后，我深愿它从未被写下，好让这些事甚至连丝毫记录也不会流传给后代。因为谁曾听闻过这等事？谁曾教导或学习过？\BibleQuote{法律将出自熙雍，上主的话将出自耶路撒冷}\BibleRef{依 2:3}；但这说法却出自何处？是哪一个低下境界吐出了这样的陈述：由\Person{玛利亚}所生的身体与\OriginalTerm{圣言}{Word}的\OriginalTerm{天主性}{Godhead}\OriginalTerm{同性同体}{coessential}？或是\OriginalTerm{圣言}{Word}被变化成了血肉、骨骼、毛发和整个身体，并从其自身的本性有了改变？或者谁曾在教会中，甚至从基督徒处听闻过，主穿戴了一个只是假象而非本质的身体？或者又有谁在不信不虔中竟走得如此之远，以至于说并坚持这\OriginalTerm{同性同体}{coessential}于圣父的\OriginalTerm{天主性}{Godhead}遭受了割损，并成为了不完美而不是完美；以及那悬挂于木架上的并非身体，而是那创世的\OriginalTerm{本体}{Essence}与\OriginalTerm{智慧}{Wisdom}本身？或者谁若听到\OriginalTerm{圣言}{Word}为自己转变了一个能受苦的身体，并非出自\Person{玛利亚}，而是出自他自己的\OriginalTerm{本体}{Essence}，能称说这话的人为基督徒？或者又是谁想出这令人憎恶的不虔，竟至让它进入他的想象，并说宣称主的身体出自\Person{玛利亚}，就是在\OriginalTerm{天主性}{Godhead}中主张\OriginalTerm{四元一体}{Tetrad}而非\OriginalTerm{天主圣三}{Triad}？那些如此思想的人，说救主从\Person{玛利亚}所摄取的身体，是属于\OriginalTerm{天主圣三}{Triad}的\OriginalTerm{本体}{Essence}。或者又有某些人吐出了与前述同样严重的不虔，即说身体并不比\OriginalTerm{圣言}{Word}的\OriginalTerm{天主性}{Godhead}更新，而是\OriginalTerm{同永恒的}{coeternal}，因为它是由\OriginalTerm{智慧}{Wisdom}的\OriginalTerm{本体}{Essence}所构成的。或者那些自称为基督徒的人又怎敢甚至怀疑，那位由\Person{玛利亚}所出、依\OriginalTerm{本体}{Essence}与本性为天主之子的主，是否按肉身是\BibleQuote{出自\Person{达味}的后裔}\BibleRef{罗 1:3}，并出自圣\Person{玛利亚}的血肉？或者又有谁如此胆大，竟说那在肉身内受苦并被钉死的基督，不是主、救主、天主和圣父之子？或者他们怎能愿意被称为基督徒，因为他们说\OriginalTerm{圣言}{Word}降在一个圣人身上，如同降在一位先知身上一样，而并未亲自成为人，由\Person{玛利亚}取得身体；而且基督是一位，而在\Person{玛利亚}之前及万世之前即为圣父之子的\OriginalTerm{天主圣言}{Word of God}是另一位？或者他们怎能是基督徒，因为他们说子是一位，而\OriginalTerm{天主圣言}{Word of God}是另一位？\par

3. 这就是备忘录的内容；种种不同的说法，但在意向与含义上却是一致的；皆趋向于不虔敬。正是为了这些事，那些以在\Place{尼西亚}所制定的教父们的信经自夸的人，彼此纷争吵闹。但我诧异的是，你的虔敬竟容忍了这些，你竟没有阻止那些说这些话的人，并向他们提出正确的信德，好叫他们听见后能缄默不语，或者若他们反对，就可被算为异端者。因为这些说法，不适合基督徒去说或听，相反，它们在各方面都与宗徒的教导背道而驰。为此，如我在上文所说，我已将他们所说的赤裸裸地写在这封信中，好让仅仅听闻的人也能看出其中所包含的羞耻与不虔。尽管本应完全谴责并揭露那些持有此类思想者的愚妄，但在此结束我的信、不再多写，也不失为一件好事。因为对于如此明显呈现为邪恶的事，无需再费时进一步揭露，免得好争辩的人认为此事尚有疑虑。仅需如此答复便已足够：我们满足于这一事实，即这不是天主教会的教导，教父们也未曾持守此见。但为了不使\BibleQuote{那些发明邪恶的人}\BibleRef{罗 1:30}以我们的全然沉默为借口而恬不知耻，从\WorkTitle{圣经}中提几点将是有益的，但愿他们能因此而羞愧，并停止这些卑污的伎俩。\par

4. 诸位先生，你们怎么会想到说，身体与\OriginalTerm{圣言}{Logos}的\OriginalTerm{天主性}{Godhead}是同一\OriginalTerm{本质}{ousia}呢？因为从这里开始是好的，为的是通过指出这个观点是错误的，也能证明所有其他观点同样错误。如今我们从神圣的\WorkTitle{圣经}中，并未发现任何类似的说法。因为\WorkTitle{圣经}说，天主以人的身体来到世上。然而，那些曾在\Place{尼西亚}聚会的教父们说，不是身体，而是圣子自己与圣父\OriginalTerm{同质}{homoousios}，且圣子虽然出自圣父的\OriginalTerm{本质}{ousia}，但身体，正如他们按照\WorkTitle{圣经}所承认的，来自\Person{玛利亚}。那么，要么你们否认\Place{尼西亚}大公会议，并如\OriginalTerm{异端者}{heretics}一般从旁引入你们的学说；要么，如果你们愿意作教父的儿女，就不要坚持与他们所写的相反的东西。因为在此你们又可看到那是多么荒谬：如果\OriginalTerm{圣言}{Logos}与那属地性质的身体\OriginalTerm{同质}{homoousios}，而按照你们自己的承认，\OriginalTerm{圣言}{Logos}又与圣父\OriginalTerm{同质}{homoousios}，就会得出连圣父自己也与那由土而生的身体\OriginalTerm{同质}{homoousios}。当你们转向另一种不虔敬的形式，说\OriginalTerm{圣言}{Logos}被变化成了\OriginalTerm{血肉}{flesh}、骨头、毛发、肌肉和整个身体，并从其自己的本性受改变时，为什么还要指责\Person{亚略派}称圣子为受造物呢？因为现在你们该公开说，他是由土而生的了；因为骨头和整个身体的性质原是从土来的。那么，你们这是何等的愚昧，竟至彼此自相矛盾？因为你们既说\OriginalTerm{圣言}{Logos}与身体\OriginalTerm{同质}{homoousios}，是将二者加以区分；又说\OriginalTerm{圣言}{Logos}被变化成了\OriginalTerm{血肉}{flesh}，则是设想\OriginalTerm{圣言}{Logos}自身的改变。如果你们竟至说出这些观点，谁还能容忍你们呢？因为你们在不虔敬上已超越任何异端。因为如果\OriginalTerm{圣言}{Logos}与身体\OriginalTerm{同质}{homoousios}，那么\Person{玛利亚}的记念和工作便是多余的，因为身体本可以在\Person{玛利亚}之前存在，正如\OriginalTerm{圣言}{Logos}也是永恒的：就是说，如果按你们所说，\OriginalTerm{圣言}{Logos}与身体\OriginalTerm{同质}{homoousios}的话。或者，\OriginalTerm{圣言}{Logos}来到我们中间，穿上那与自己\OriginalTerm{同质}{homoousios}的东西，或改变自己的本性而成为身体，又有什么必要呢？因为\OriginalTerm{天主性}{Deity}并不援助自己\BibleRef{希 2:16}，以便穿上自己\OriginalTerm{本质}{ousia}之物，正如\OriginalTerm{圣言}{Logos}也没有犯罪，好使他赎回他人的罪，以便他变成身体后，能为自己献上\OriginalTerm{祭献}{sacrifice}，并救赎自己。\par

5. 但这并非如此，绝无此意。因为正如\OriginalTerm{宗徒}{apostle}所说，他\BibleQuote{援助了\Person{亚巴郎}的后裔}\BibleRef{希 2:16}，于是他应该在各方面相似弟兄们，并取得一个像我们一样的身体。这就是为什么\Person{玛利亚}确实被预定了的原因，好使他能从她取得身体，并作为自己的为我们献上。而\Person{依撒意亚}在他的预言中指向了这一点，说：\BibleQuote{看，\OriginalTerm{贞女}{Virgin}}\BibleRef{依 7:14}；同时\Person{加俾额尔}被派到她那里——不是随便到一个贞女那里，而是\BibleQuote{到一位已许配于男子的贞女那里}\BibleRef{路 1:27}，为藉着这已许配的男子表明\Person{玛利亚}确实是一个人。为此，\WorkTitle{圣经}也提到她生产，并述说她用襁褓裹起他；因此，他所吮吸的乳房也被称为有福的。他被献为\OriginalTerm{祭献}{sacrifice}，因为那位诞生者开了母胎。凡此种种，都是\OriginalTerm{贞女}{Virgin}确实生产的证据。而\Person{加俾额尔}毫不含糊地向她宣讲了福音，不仅说\BibleQuote{那要诞生的圣者}，以免以为身体是从外界附加于她的，而是说\BibleQuote{要由你而生的}，好使人们相信那诞生者确实由她自然而来，因为自然清楚地表明，贞女若未生产，就不可能产生乳汁，身体若未先经自然生产，就不可能用乳汁滋养并用襁褓包裹。这就是他于第八天受割损、\Person{西默盎}将他抱在怀中、他成为孩童、十二岁时成长，以及他到三十岁的意义所在。因为，并非如某些人所设想，是\OriginalTerm{圣言}{Logos}的\OriginalTerm{本质}{ousia}发生了改变并受了割损，因为\OriginalTerm{本质}{ousia}是不能改变或变化的。因为\OriginalTerm{救主}{Saviour}自己说：\BibleQuote{注意，注意，我是上主，我决不改变}\BibleRef{拉 3:6}；而\Person{保禄}则写道：\BibleQuote{耶稣基督昨天、今天、直到永远，常是一样}\BibleRef{希 13:8}。但是，在那受割损、被怀抱、吃喝、疲倦、被钉在木架上并受苦的身体内，存在着\OriginalTerm{不可受难}{impassible}且\OriginalTerm{无形体}{incorporeal}的天主\OriginalTerm{圣言}{Logos}。正是这身体被安放在坟墓中，那时\OriginalTerm{圣言}{Logos}已离开它，却并未与它分离，为如\Person{伯多禄}所说，去向那些在监狱中的灵魂宣讲\BibleRef{伯前 3:19}。\par

6. 这尤其显示出那些声称圣言被转变成了骨肉之人的愚蠢。因为如果真是如此，坟墓就没有必要了。因为身体本可以自行前往阴府，向那里的灵魂宣讲。但事实上，祂亲自去宣讲，而身体则由若瑟用细麻布裹好，安放在哥耳哥达。\BibleRef{谷 15:46} 这样就向众人显明，身体并不是圣言，而是圣言的身体。这正是多默在复活后所触摸的，他看见了钉痕，这些都是圣言亲自承受的，祂看着钉子钉入自己的身体，尽管能够阻止，却没有阻止。相反，无形体的圣言将身体的特性归于自己，视之为自己的身体。当差役击打那身体时，祂如同自己受苦一样问道：\BibleQuote{你为什么打我？}\BibleRef{若 18:23} 而圣言本性上是不可触摸的，却说：\BibleQuote{我将我的背转给打击我的人，把我的腮转给扯我胡须的人；对于侮辱和唾污，我没有遮掩我的面。}\BibleRef{依 50:6} 因为圣言的人性身体所遭受的，居住在身体内的圣言都归于自己，好使我们能够分享圣言的天主性。这实在奇妙：祂既受苦，又未受苦。说祂受苦，是因为祂自己的身体受苦，而祂在这受苦的身体内；说祂未受苦，是因为圣言本性是天主，不能受苦。当无形体的祂在能受苦的身体内时，这身体内有着不能受苦的圣言，圣言正在消除身体固有的软弱。但祂这样做，并且事就这样成了，为的是祂亲自取得属于我们的，并将之作为祭献，好能除免它，反过来又将属于祂的赐给我们，并让宗徒说：\BibleQuote{这可朽坏的，必须穿上不可朽坏的；这可死的，必须穿上不可死的。}\BibleRef{格前 15:53}\par

7. 这事并非如某些人所设想的是虚构的：万万不可这样想！相反，救世主真实地成了人，从而实现了全人的救恩。因为，如果圣言只是虚构地存在于身体内，如他们所说，而虚构是指想象中的，那么，人的救恩和复活也就仅仅是表面的，正如最不虔敬的摩尼所主张的。但事实上，我们的救恩并非仅仅是表面的，也不仅局限于身体，而是整个的人，包括身体和灵魂，都真实地在圣言自己内获得了救恩。所以，由玛利亚所生的，依照圣经而言，其人性是真实的，主的身体也是真实的；而它之所以真实，是因为它与我们的身体相同，因为玛利亚是我们的姐妹，就我们全都出自亚当而论。任何人若记得路加所写的事，就不会对此有所怀疑。因为主从死者中复活后，当有些人以为他们所见的是鬼神，而非那由玛利亚所出的身体中的主时，祂说：\BibleQuote{你们看看我的手，我的脚，是我自己。摸摸我看，因为鬼是没有肉，没有骨头的，如同你们看我却是有的。他说了这话，就把手和脚伸给他们看。}\BibleRef{路 24:39} 由此，那些胆敢声称主转变成了肉和骨的人，便可被驳倒。因为祂并没有说：\Quote{如同你们看我是肉和骨}，而是说：\Quote{如同你们看我却是有的}，为的是不让人以为圣言本身转变成了这些东西，而是让人相信祂在复活后也如同在死亡前一样拥有它们。\par

8. 这些既已证明，再涉及其他论点或进行相关讨论便是多余了，因为圣言在其内的身体并非与天主性同性同体，而是真实地由玛利亚所生，而圣言本身并未转变成骨肉，而是降来在肉身内。因为若望所说：\BibleQuote{圣言成了血肉}\BibleRef{若 1:14}，其意义正如我们从一处类似的经文所见的；因为保禄写道：\BibleQuote{基督为我们成了可咒骂的}\BibleRef{迦 3:13}。正如祂并非自己真的成了可咒骂的，而是因为祂承担了我们的咒骂才这样说，同样，祂成了血肉，也并非被转变成了血肉，而是因为祂为我们承担了生活的血肉，成了人。因为说\Quote{圣言成了血肉}，就等于说\Quote{圣言成了人}；正如岳厄尔所说：\BibleQuote{我要将我的神倾注在一切血肉上}\BibleRef{岳 2:28}；因为这许诺并不涉及无理性的动物，而是针对人，主正是为了人而成了人。既然这是上述经文的意义，那么所有那些以为由玛利亚所出的肉身先于玛利亚而存在，以及以为圣言在她之先就具有一个人性灵魂，并且即便在祂来临之前也常存在于其中的人，都将合理地自讼己罪。那些声称血肉不可能死亡，而是属于不朽本性的人，也将停止其说。因为如果它没有死，保禄怎能对格林多人说\BibleQuote{基督照经上记载的，为我们的罪死了}\BibleRef{格前 15:3}，或者，如果祂不曾死，又怎能复活呢？再者，那些甚至考虑过如果声称肉身出自玛利亚，就会由三位变为四位的人，也将深自羞愧。因为（他们论证说）若我们说肉身与圣言同一本体，三位仍为三位；因为此时圣言并没有引入外来的元素；但若我们承认由玛利亚所出的肉身是人性的，那么由于肉身体本体上是外来的，而圣言在其内，结果肉身的附加就使三位变成了四位。\par

9. 当他们这样论证时，并未察觉自己陷入了何等自相矛盾。因为即使他们说身体并非来自\Person{玛利亚}，而是与圣言\OriginalTerm{同性同体}{coessential}，但（这正是他们故意掩饰，以免其真实见解被相信之处）从他们的前提本身即可证明，这必然涉及\Quote{\OriginalTerm{四一}{Tetrad}}。正如圣子，按教父们的教导，与圣父同性同体，却非圣父自身，而是作为子与父称为同性同体；同样，那被他们称为与圣言同性同体的身体，并非圣言自身，而是另一个实体。果真如此，依他们所示的，他们的\OriginalTerm{三一}{Triad}便将变成四一。因为那真实、绝对完美且不可分割的三一，并不像这些人所设想的三一那样能够有所增加。况且，在真实的天主之外再设想的另一位神，这些人又如何算得基督徒呢？再者，在他们的另一谬误中，也可看出其愚昧之深。因为，如果他们以为，圣经既包含并陈述救主的身体是人的，且出自玛利亚，于是四一便取代了三一，仿佛那身体造成了增加，那么他们就大错特错了，甚至谬至将受造物等同于造物主，以为\OriginalTerm{天主性}{Godhead}可有所增加。他们也未能领悟，圣言所以成了血肉，并非为给天主性增加什么，而是为使肉身复活。圣言由玛利亚而生，也非为使自己更加优越，而是为救赎人类。那么，他们怎能认为，那被圣言所救赎并赋予生命的身体，竟给赋予它生命的圣言在神性方面造成了增加呢？恰恰相反，倒是那人的身体自身，因与圣言的契合与结合，获得了极大的增加。因为它由有死的变为不死的；虽是属血气的身体，却成为属灵的；虽出自尘土，却进入了天门。因此，这天主圣三，虽然圣言从玛利亚取了身体，仍是三一，既无增亦无减；它永远是圆满的，在这三一中，我们认识的是一个天主性，所以在教会中所宣讲的也是唯一的天主，即圣言的父。\par

10. 为此，他们这些曾说过那出自\Person{玛利亚}的并非真正的基督、主或天主的人，今后也将住口无言了。因为如果他在身体内不是天主，那他刚一出自玛利亚，怎么立即被称为\BibleQuote{厄玛奴耳，意思是\InnerQuote{天主与我们同在}}\BibleRef{玛 1:23}呢？再者，如果圣言不在肉身内，\Person{保禄}为什么写给罗马人说：\BibleQuote{基督按血统而言，也是从他们出来的，他是在万有之上的天主，永受赞颂，阿们。}\BibleRef{罗 9:5}因此，甚至那些先前否认被钉十字架者是天主的人，也该承认自己错了；因为神圣的圣经，尤其是\Person{多默}，在看见了他身上的钉痕以后，大声呼喊说：\BibleQuote{我主，我天主！}\BibleRef{若 20:28}，都指证他们。因为圣子既是天主，又是光荣的主\BibleRef{格前 2:8}，就居住在那被不光彩地钉死并受辱的身体内；但那身体受难之时，虽被悬于木架，胁膀被刺透，流出血与水，因为它原是圣言的宫殿，故仍充满\OriginalTerm{天主性}{Godhead}。为此，太阳看见自己的创造者在自己受辱的身体内受苦，便收敛了光芒，使大地昏暗。但那身体自身虽属可死的本性，却因在其内的圣言，超越了它的本性而复活了；它已脱离了本性的腐朽，穿上了那超越人的圣言，成了不朽的。\par

11. 至于某些人的幻想，他们说圣言降临在某个特定的人，即\Person{玛利亚}之子，之上，如同它降临在每位先知上一样，讨论此事实属多余，因为他们的疯狂本身便明显地带着自我谴责。因为如果他是这样降临的，为什么那人由童贞女所生，而不像其他人由男女结合而生呢？因为每位圣徒也都是这样生的。或者，如果圣言是这样来临的，为什么没有说每位圣徒的死都是为了我们，而只说这一位的死呢？或者，如果圣言在每位先知身上都曾这样在我们中间居住，为什么唯独论及那由玛利亚所生者时，才说\BibleQuote{在这世代的末了，只一次地临在了}\BibleRef{希 9:26}呢？或者，如果他的来临像他从前在圣徒们身上临在一样，为什么当众人死后尚未复活时，唯独玛利亚之子在第三日复活了呢？或者，如果圣言在其他事例中也是如此临在的，为什么唯独玛利亚之子被称为厄玛奴耳，仿佛由她所生的身体充满了\OriginalTerm{天主性}{Godhead}呢？因为厄玛奴耳解释为\Quote{天主与我们同在}。或者，如果他是这样来临的，为什么当每位圣徒吃、喝、劳苦、死亡时，没有说那一位（圣言）也曾吃、喝、劳苦、死亡，而只论及玛利亚之子呢？因为那身体所受的，被说成是圣言所受的。而对于其他人，我们仅仅得知他们出生、受生；唯独论及玛利亚之子时说：\BibleQuote{圣言成了血肉}。\par

\OriginalTerm{圣言}{Word}虽为使人能说预言而临于所有其他的人，却从\Person{玛利亚}取了肉躯，以人的形态出现；祂就本性及本质而言是天主的圣言，但按肉躯而言，是\Person{达味}的后裔，出自\Person{玛利亚}的血肉，正如\Person{保禄}所说。圣父曾在\Place{约旦河}和山上指出他，说：\BibleQuote{这是我的爱子，我所喜悦的。}\OriginalTerm{阿里乌斯派}{Arians}否认他，但我们认识并朝拜他，不把圣子与圣言分开，而是知道圣子即是圣言本身，万物都是藉祂而造成的，我们也是藉祂而获得救赎。为此，我们惊讶于你们中间竟因如此清楚的事起了纷争。但感谢主，我们读到你们的记录时虽甚忧伤，却也同样因它们的结论而欢喜。因为他们怀着和谐离去，并在虔诚和正统信德的宣认上平安地达成了一致。这一事实促使我，经过先前许多考虑，写下这几行文字；因为我担心，我的缄默可能会使此事给那些其和谐带给我们喜乐的人带来痛苦而非喜乐。因此，我首先请求尊驾，其次请求读者，善意地看待我的信；若信中有任何欠缺虔诚之处，请予以纠正，并告知我。但若这信因出于不善言辞之人而写得低于主题且不完善，请大家宽容我在言辞上的软弱。问候你们那里的所有弟兄。我们这里的所有人也问候你们；愿你们在主内生活安好，亲爱且深为渴念的诸位。\par

\SourceRecord{Translated by Archibald Robertson.；Nicene and Post-Nicene Fathers, Second Series；Vol. 4.；Edited by Philip Schaff and Henry Wace.；Buffalo, NY: Christian Literature Publishing Co.,；1892.；<http://www.newadvent.org/fathers/2806059.htm>.}

% structure: p0001 source=h1 target=section reason=page-title

\section{\WorkTitle{第六十封信}}

\Emph{致\Person{阿德尔菲乌斯}主教兼精修者：驳斥亚略派}\par

我们已拜读了您的虔敬写给我们的信，并真诚地赞许您对基督的虔敬。首先，我们光荣天主，祂赐给了您这样的恩宠，使您不但持有正确的见解，而且，在尽可能的范围内，不致对魔鬼的诡计\BibleRef{格后 2:11}一无所知。然而，我们对异端者的乖张感到惊异，因为他们竟堕入了如此不虔敬的深渊，以至于连理智都不复存在，悟性也全然败坏了。但这种企图是魔鬼的阴谋，是对悖逆的犹太人的仿效。因为正如后者，处处受驳斥时，仍不断为自己招致伤害的借口，只要能否认主，将先知所预言的灾祸招到自己身上，同样，这些人见自己处处被排斥，察觉自己的异端已为众人憎恶，便显出自已正是\BibleQuote{捏造恶事的，}\BibleRef{罗 1:30}为要持续与真理争斗，好成为始终如一、名副其实的基督的敌人。因为，他们这种新的祸害从何而生？他们怎敢对救主发出这种新的亵渎？然而，不虔敬的人，看来真是一无是处，实在\BibleQuote{是在信德上可废弃的}\BibleRef{弟后 3:8}。因为从前，他们否认天主独生子的天主性时，还至少假装承认祂降生成人。可是现在，逐步变本加厉，连这种观点也丢弃了，在各方面都成了无神者，既不承认祂是天主，也不信祂曾成为人。如果他们相信这事，就不会对您的虔敬说报道的那种话了。\par

然而，亲爱的、我们最切望的您，在驳斥、劝诫并指责这类人时，做了合乎教会传统和您对主的虔敬的事。但因他们受其父魔鬼的煽动，如经上所载：\Quote{他们既不知道，也不明白，仍在黑暗中行走}，就让他们从您的虔敬得知，他们的这种错谬属于\Person{瓦伦廷}、\Person{马西昂}和\Person{摩尼}之流；他们中有些人以幻象代替真实，另有些人分裂那不可分裂的，否认\BibleQuote{圣言成了血肉，寄居在我们中间}\BibleRef{若 1:14}的真理。那么，既然他们赞同那些人，为何不也继承他们的名号呢？按理，他们既持有他们的错谬，也该拥有他们的名号，从此被称为\Person{瓦伦廷派}、\Person{马西昂派}和\Person{摩尼派}。或许这样，因这些名号声名狼藉而感到羞耻后，他们就能发觉自己堕入了何等不虔敬的深渊。本来，按照宗徒的劝告，我们完全有权不理睬他们：\BibleQuote{对异端者，在谴责过一次两次之后，就该远离他；因为像这样的人已经败坏了，犯了罪，自绝自责。}\BibleRef{铎 3:10}，更何况，先知论到这类人曾说：\Quote{愚昧的人会口出愚言，心里图谋虚无的事}。但是，因他们同自己的首领一样，也像狮子般巡游，寻找可吞噬的愚仆之人\BibleRef{伯前 5:8}，我们不得不回信给您的虔敬，好使弟兄们因您的劝诫再受教导，从而更厌弃这些人的虚妄言论。\par

我们不朝拜受造物。绝不是！因为这样的错误属于外教人和亚略派。我们所朝拜的，是降生成人的造物主、天主圣言。因为，肉身本身虽是受造界的一部分，却成了天主的身体。我们既不将这样的身体与圣言分开，单独朝拜它；也不会在朝拜圣言时，将祂与肉身远远隔开。而是如上所述，知道\Quote{圣言成了血肉}，我们在祂带着肉身之后，依然承认祂是天主。那么，有谁会如此糊涂，竟对主说：\Quote{离开你的身体，好让我朝拜你}；或有谁如此不虔敬，竟会因着那身体，与糊涂的犹太人一同说：\BibleQuote{你只是人，却把自己当作天主}\BibleRef{若 10:33}？但那个癞病人却不是这类人，因为他朝拜了居于肉身内的天主，并认出祂是天主，说：\BibleQuote{主，你若愿意，就能洁净我}\BibleRef{玛 8:2}。他既不因肉身而认为天主圣言是受造物，也不因圣言是万物的创造者而轻视祂所穿戴的肉身。他朝拜了那居住在受造殿宇中的宇宙创造者，因而得了洁净。同样，那患血漏的妇人，因着信德，只摸了祂的衣裳繸头，就得到了痊愈；\BibleRef{玛 9:20}波涛汹涌的大海一听从降生成人的圣言，风浪便平息了；而那胎生的瞎子，因圣言藉肉身所吐的唾沫而得到痊愈\BibleRef{若 9:6}。还有更伟大、更惊人的事（或许正是这事触怒了那些极不虔敬的人），就是当主被悬在真实的十字架上时（因为那是祂的身体，而圣言就在其中），太阳昏暗了，大地震动，岩石崩裂，圣殿的帐幔裂开，许多安眠的圣人们的身体复活起来。\par

4. 这些事就这样发生了，没有人怀疑，如同亚略派如今所冒昧怀疑的，是否该相信降生成人的圣言；然而，即使看见那人，他们就认出他是他们的造物主；当他们听到人的声音时，并不因为那是人的声音就说圣言是受造物。相反，他们战栗，且认出那确是从\OriginalTerm{圣殿}{holy Temple}发出的声音。那么，这些不敬虔者怎能不惧怕，恐怕\BibleQuote{他们既不肯认真地认识天主，天主也就任凭他们陷于邪恶的心思，去行那些不合理的事}\BibleRef{罗 1:28}呢？因为受造界不朝拜受造物。她也不因他的肉体而拒绝朝拜她的主。然而，她看见造物主在肉身内，并且\BibleQuote{在耶稣的名字前，无不屈膝}已经跪下，的确，\BibleQuote{将要屈膝——无论天上的、地上的和地下的——一切唇舌无不明认}，不论亚略派认可与否，\BibleQuote{耶稣是主，以光荣天主圣父}\BibleRef{斐 2:10}。因为肉体并未减损圣言的光荣；绝非如此：相反，肉体由他而得光荣。圣子既具有天主的形体，却取了奴仆的形体，这并未剥夺他的天主性。相反，他因而成了全人类和整个受造界的救主。并且，如果天主差遣那由女人所生的圣子，这事并不使我们感到羞耻，反而成为荣耀和大恩宠。因为他成了人，是为使我们在他内神化；他由女人所生，且由童贞女所孕育，为将我们错误的族类转移到他自己身上，使我们从此成为圣洁的族类，并成为\BibleQuote{有分于天主性体的人}\BibleRef{伯后 1:4}，正如圣\Person{伯多禄}所写。并且，\BibleQuote{法律既因肉性的软弱而有所不能行的，天主遂派遣了自己的儿子，带着罪恶肉身的形状，为了罪恶，在肉身上定了罪恶的罪案}\BibleRef{罗 8:3}。\par

5. 因此，既然圣言取了肉体，为拯救众人，使所有从死者中复活，并为罪成就救赎，那些轻视这肉体的人，以及那些因此控告天主圣子为受造物或被造者的人，岂不显得忘恩负义，并该受一切憎恨吗？因为他们无异于向天主喊叫说：\Quote{不要派遣你的独生子在肉身内，不要让他从童贞女取得肉体，以免他从死亡和罪恶中救赎我们。我们不愿他在肉身内来临，以免他为我们受死：我们不愿圣言成为血肉，以免他在其中成为我们的中保，使我们得以进到你面前，并居住天上的居所。让天门关闭吧，免得你的圣言透过\BibleQuote{帐幔，即他的肉身}\BibleRef{希 10:20}，为我们祝圣那道路。}这就是他们所捏造的错误，以魔鬼般的胆大妄为所发泄的言语。凡不愿朝拜成为肉身的圣言的人，就是对他降生成人忘恩负义。而那些将圣言与肉身分开的人，就不认为已发生了一次为罪获得的救赎，或死亡的一次性毁灭。这些不敬虔的人究竟会在哪里找到救主所取的肉身，而能离开他呢？他们竟敢说：\Quote{我们不朝拜有肉身的圣言，而是将那身体分开，只朝拜他。}请看，圣\Person{斯德望}看见在天上，主站在\BibleQuote{天主的右边}\BibleRef{宗 7:55}，而天使们向门徒们说：\Quote{他怎样升了天，也要怎样降来。}主亲自向父说：\BibleQuote{我愿我在哪里，他们也同我在一起。}\BibleRef{若 17:24}的确，如果肉身与圣言不可分离，岂不必然推出：这些人要么放弃他们的错误，今后因我们的主耶稣基督的名朝拜父；要么，如果他们不朝拜或不侍奉那在肉身内而来的圣言，便被从各处驱逐，不再被视为基督徒，而被视为外邦人，或列于犹太人之中。\par

6. 如上所述，那些人的疯狂和胆大妄为就是如此。但我们的信仰是正确的，它源于宗徒们的教导和教父们的传授，并由新约和旧约所证实。因为先知们说：\Quote{发出你的圣言和你的真理}；又说：\Quote{看，一位童贞女将怀孕生子，人将称他的名字为厄玛奴耳，意思是：天主与我们同在。}这若不表示天主已降临于肉身，那又是什么意思呢？而宗徒传统则藉着圣\Person{伯多禄}的话教导说：\Quote{基督既然在肉身上为我们受了苦难，...}；以及圣\Person{保禄}所写的：\BibleQuote{期待所希望的幸福，和我们伟大的天主及救主耶稣基督荣耀的显现。他为我们舍弃了自己，是为救赎我们脱离一切罪恶，并洁净我们，使我们成为他的选民，热心行善}\BibleRef{铎 2:13}。他若没有穿戴肉身，怎会把自己献出呢？因为他献出了肉躯，为我们舍弃了自己，为能藉着在肉身上经历死亡，\BibleQuote{摧毁那握有死亡权势者——魔鬼}\BibleRef{希 2:14}。因此，我们常因耶稣基督的名感谢，并且绝不轻视透过他而临于我们的恩宠。因为救主降生成人的来临，已成为整个受造界的赎价与救恩。所以，蒙爱的、极所渴望的诸位，愿我所说的，能提醒那些爱主的人；至于那些效法\Person{犹达斯}的行为，离弃主而投靠\Person{盖法}的人，愿他们透过这些事得以受教改进，如果他们愿意，如果他们知耻的话。并让他们知道，我们朝拜那在肉身内的主，并不是朝拜受造物，而是如先前所说，朝拜那穿戴了受造身体的造物主。\par

但我们希望你的虔诚能问他们这一点：当以色列人受命上耶路撒冷，到主的圣殿朝拜时，那里有约柜，\BibleQuote{柜上有荣耀的\OriginalTerm{革鲁宾}{Cherubim}遮着\OriginalTerm{赎罪盖}{Mercy-seat}\BibleRef{希 9:5}}，他们做得对还是错？如果他们做了恶事，那么轻慢这法律的人怎会当受惩罚呢？因为经上记载：若有人轻视而不上去，这人必从民中灭绝。但如果他们做得对，并因此蒙天主悦纳，那么，\OriginalTerm{亚略派}{Arians}——可憎且是一切异端中最可耻的——岂不更该受毁灭吗？他们赞同先前的民族对圣殿所表示的尊敬，自己却不肯朝拜那在肉身内如同在圣殿中的主。然而，先前的圣殿是由石头和金子筑成的影儿。但当实体来到时，预像便从此停止了，按主的话说：\BibleQuote{没有一块石头留在另一块石头上，而不被拆毁的}\BibleRef{玛 24:2}。当他们看到石头的圣殿时，并没有认为在圣殿中说话的主是受造物；也没有轻视圣殿而退避远处去朝拜。相反，他们按着法律来到那里，朝拜那从圣殿中发出神谕的天主。既然事实如此，那么朝拜主的身体，岂能不是理所当然的呢？这身体是最神圣、最可敬的，由总领天使\Person{加俾额尔}所预报，由圣神所形成，并成为圣言的\OriginalTerm{衣袍}{Vesture}。至少，那圣言伸出有形的手，拉起患热症的女子\BibleRef{谷 1:31}；以人的声音，命\Person{拉匝禄}从死者中起来\BibleRef{若 11:43}；再者，当祂在十字架上伸开双手时，倾覆了那现今在悖逆之子身上活动的空中权势的首领，并为我们开辟了通往天上的道路。\par

因此，谁若轻慢圣殿，便是轻慢在圣殿中的主；谁若将圣言与身体分离，便是轻视在祂内赐予我们的恩宠。最不虔敬的\OriginalTerm{亚略派}{Arians}疯狂之徒，既不要因为身体是受造的，就以为圣言也是受造物；也不要因为圣言并非受造物，就贬低祂的身体。因为他们的错误实在令人惊讶，他们同时又混淆和扰乱了万事，并且只图谋借口，好将造物主列于受造物之中。\par

但让他们听着：如果圣言是受造物，祂就不会取受造的身体以赋予它生命。因为受造物能从自身也需要拯救的受造物得到什么助益呢？但正因为圣言是造物主，自己造了万有，所以祂也在\OriginalTerm{时代的终结}{consummation of the ages}\BibleRef{希 9:26}时穿上了受造物，好使祂作为造物主能再次祝圣它，并能使它复原。然而，受造物绝不能被受造物拯救，正如万有若非主创造，便不能由受造物所造一样。因此，他们不要用谎言攻击圣经，也不要开罪纯朴的弟兄；但若他们愿意，就让他们回头改念，不再朝拜受造物，而不朝拜创造万有的天主。如果他们执意守着他们的不敬，就任凭他们自食其果，让他们如同他们的父魔鬼一般咬牙切齿吧。因为天主教会的信德清楚知道，天主的圣言是万有的创造者和建树者；我们明白，\BibleQuote{在起初已有圣言，圣言与天主同在}\BibleRef{若 1:1}，如今祂为了我们的救恩也成了人，我们便朝拜祂，不是以为祂带着身体来，使自己与身体同等，而是作为主，取了奴仆的形体，作为造物主与创造者，来到受造物内，为能在其中解救万有，引领世界亲近圣父，并使天地万物都与祂和好。因为如此我们也承认祂的天主性，就是圣父的天主性，并朝拜祂\OriginalTerm{降生成人的临在}{Incarnate Presence}，即使那\OriginalTerm{亚略派}{Arians}的疯狂之徒自身破裂、粉身碎骨。\par

请代我问候所有爱主耶稣基督的人。我们祈求你们安好，并在主前记念我们，我们亲爱的、至为想念的诸位。如有需要，请将这封信读给\Person{依厄拉克斯}司铎听。\par

\SourceRecord{Translated by Archibald Robertson.；Nicene and Post-Nicene Fathers, Second Series；Vol. 4.；Edited by Philip Schaff and Henry Wace.；Buffalo, NY: Christian Literature Publishing Co.,；1892.；<http://www.newadvent.org/fathers/2806060.htm>.}

% structure: p0001 source=h1 target=section reason=page-title

\section{信函 61}

\Emph{致\Person{马克西穆斯}的信。（约写于公元371年）}\par

\Person{亚大纳削}致我们亲爱且最热切渴望的儿子、哲学家\Person{马克西穆斯}，在主内致候。\par

阅读了你今次的来信，我赞成你的虔诚；但对那些\BibleQuote{他们不明白自己所说和所坚持的事}\BibleRef{弟前 1:7}的人轻率鲁莽感到惊奇，我原本决定缄默不语。因为就如此显明、比光更清晰的事作出回答，只会给那些无法无天之人提供厚颜无耻的借口。这是我们从救主那里学到的。因为当\Person{彼拉多}洗了手，默许当时犹太人的诬告时，主就不再回答他，而是在梦中警告他的妻子，好使人相信那位受审判者是在能力上，而非言语上的天主。祂也未曾答复\Person{盖法}的愚妄，却亲自以应许引众人趋向认识。因此我耽搁了一段时间，因着那些无耻之人好争辩，才勉强顺从了你们对真理的热心。我所口授的，无非是你信中所含内容，好使对手从今以后在他所反对的各点上得以信服，并且能\BibleQuote{谨守口舌，不说欺诈的话}。但愿他们不再加入昔日的犹太人，路过讥诮挂在\OriginalTerm{树}{Tree}上的那一位，说：\BibleQuote{如果你是天主子，从十字架上下来吧！}\BibleRef{玛 27:40}。但倘若在这以后他们仍不屈服，你要记着宗徒的训令：\BibleQuote{对异端者，在第一次和第二次劝诫后，就该回避他，因为知道这样的人背弃了正道，犯罪作恶，自我定罪}\BibleRef{铎 3:10}。因为如果他们这样放肆的是外邦人，或是\OriginalTerm{犹太化派}{Judaisers}，那么，让他们要么像犹太人那样以基督的十字架为绊脚石，要么像外邦人那样以它为愚妄吧。\BibleRef{格前 1:23}但他们若自称是基督徒，就让他们学习：那被钉的基督同时是光荣的主、天主的大能和天主的智慧。\par

但如果他们怀疑祂到底是不是天主，就让他们恭敬\Person{多默}吧，他曾触摸被钉者，并宣认祂是主、天主。\BibleRef{若 20:28}或者让他们敬畏主自己，祂在给门徒洗脚后说：\BibleQuote{你们称我\InnerQuote{主}、\InnerQuote{师傅}，说得正对，我原是。}然而，祂以给门徒洗脚时的同一身体，也亲自将我们的罪背负到\OriginalTerm{树}{Tree}上\BibleRef{伯前 2:24}。祂被见证为受造界的主宰，因为太阳收去光芒，大地震动，岩石崩裂，行刑者认出被钉者真是天主子。因为他们所见的身体，并非某个凡人的，而是天主的身体；祂在这身体内，甚至在被钉时，还复活了死人。所以，他们贸然声称天主的圣言进入某个圣洁的人，是毫无益处的；因为对每位先知和其他圣人而言都是如此，若是那样，祂就显然在每个人身上出生和死亡了。但绝非如此，愿这种想法远离。而是\BibleQuote{到了今世的末期}\BibleRef{希 9:26}，为除去罪恶，\BibleQuote{圣言成了血肉}\BibleRef{若 1:14}，由童贞\Person{玛利亚}出生，成为与我们相似的人，正如祂对犹太人所说：\BibleQuote{你们为什么想杀我，一个对你们讲论真理的人呢？}我们被神化，不是因分沾某个凡人的身体，而是因领受圣言自身的身体。\par

对此我也大为惊奇，他们竟敢认为圣言成为人是由祂的本性所致。若果如此，纪念\Person{玛利亚}就成了多余。因为本性不知道童贞女不藉男人而生育。因此，按照圣父的美意，祂本是真天主，就本性而言是天主的圣言和智慧，却为我们的救恩，在肉身内成为人，好能有可献上的，拯救我们众人，\BibleQuote{即那些因死亡的恐怖，一生当奴隶的人}。因为将自己交付于我们的，并非某个凡人；因为按对\Person{亚当}内众人所说的话，人人都受死亡判决：\BibleQuote{你既是土，仍要归于土}。也并非任何受造物，因为一切受造物都易变化。而是圣言自己为我们将自己的身体献上，为使我们的信仰和希望不在于人，而是信仰在天主圣言自己内。看，即使祂成了人，我们仍看到祂的光荣，\BibleQuote{正如父独生者的光荣，满溢恩宠和真理}。因为祂藉肉体所承受的，祂以天主性显扬了。祂虽然肉身饥饿，却以天主性饱饫饥饿的人。若有人因肉身的状况而跌到，就让他因天主所行的事业而相信吧。因为祂按人性询问\Person{拉匝禄}葬在哪里，却以天主性将他复活。所以，谁也不要讥笑，称祂为孩子，引述祂的年龄、成长、饮食和受苦，以免否认肉身所宜有之事时，全然否认祂在我们中间的居留。正如祂成为人并非由祂的本性所致，同样，当祂取了身体，显示属肉身的事乃是应当的，免得\OriginalTerm{摩尼}{Manichæus}的幻影理论得势。再者，当祂带着肉身行事时，不隐藏属天主性的事，也是应当的，免得\Place{撒摩撒塔}人找到借口称祂为人，位格上与天主圣言有别。\par

因此，让不信者明白，并学习：当祂作为婴儿躺在马槽里时，祂使贤士服从并受他们朝拜；当祂作为孩子下到\Place{埃及}去时，祂摧毁了那地偶像的手制品；而在肉身内被钉死后，祂复活了早已腐朽的死者。这也已向众人显明：祂忍受一切不是为了自己，而是为了我们，好使我们由祂的苦难穿上免于苦难和不朽\BibleRef{格前 15:53}，并常住于永生。\par

5. 因此，我依照上述你来信的内容，简要地口授了这些，没有进一步详述任何一点，只提到与\OriginalTerm{圣十字}{Holy Cross}相关的事，好让那些轻蔑之人在他们反感之处受到更好的教导，并能朝拜\OriginalTerm{被钉十字架者}{Crucified}。但你要竭力劝说那些不信者；或许他们能由无知达到认识，并正确相信。尽管你来信所写的已足够，但我仍为那些好争辩的人着想，加添我所写的，以作提醒；目的并非要驳倒他们冒失的言论，使之蒙羞，而是要藉着提醒，使之不致忘记真理。愿在\Place{尼西亚}的教父们所宣认的信理得胜。因为这信理正确无误，且足以推翻一切异端，无论多么不虔敬，尤其是\OriginalTerm{亚略异端}{Arians}——它反对\OriginalTerm{天主的圣言}{Word of God}，并且必然地亵渎祂的圣神。请向所有持守正信的人致意。所有与我们同在的人都问候你们。\par

\SourceRecord{Translated by Archibald Robertson.；Nicene and Post-Nicene Fathers, Second Series；Vol. 4.；Edited by Philip Schaff and Henry Wace.；Buffalo, NY: Christian Literature Publishing Co.,；1892.；<http://www.newadvent.org/fathers/2806061.htm>.}

% structure: p0001 source=h1 target=section reason=page-title

\section{书信六十二}

\Emph{致\Person{若望}与\Person{安提约古}}.\par

\Person{亚大纳削}致\Person{若望}与\Person{安提约古}，我们亲爱的儿子和在主内的同僚司铎，问候。\par

我刚才很高兴收到你们的信，尤其因为你们是从\Place{耶路撒冷}写来的。我感谢你们将聚集在那里的弟兄们，以及那些因争议点想要扰乱纯朴者的人，告诉了我。但对于这些事，愿宗徒吩咐他们，不要听从那些争辩言辞、只求讲述和听闻新奇之事的人。但是你们，既有稳固的根基，就是我们的主耶稣基督，以及教父们关于信仰的宣认，就该避开那些想要在这宣认上增加或删减的人，而要以弟兄们的益处为目标，好使他们敬畏天主并遵守祂的诫命，从而既凭教父们的教导，也凭遵守诫命，他们能在审判之日，在主面前显为蒙悦纳的。但我对那些胆敢攻击我们敬爱的\Person{巴西略}主教——天主的真仆人——的人的大胆妄为，极为惊骇。因为从这类虚妄的言论，就可以证明他们甚至连教父们的宣认也不爱。\par

请问候弟兄们。与我在一起的人也问候你们。我祈求你们在主内安好，亲爱的和极渴慕的儿子。\par

\SourceRecord{Translated by Archibald Robertson.；Nicene and Post-Nicene Fathers, Second Series；Vol. 4.；Edited by Philip Schaff and Henry Wace.；Buffalo, NY: Christian Literature Publishing Co.,；1892.；<http://www.newadvent.org/fathers/2806062.htm>.}

% structure: p0001 source=h1 target=section reason=page-title

\section{书信六十三}

\Emph{致帕拉迪司铎的信}。\par

致我们亲爱的儿子\Person{帕拉迪}司铎，主教\Person{亚大纳削}在主内致候。\par

我高兴地也收到了你单独写的信，尤其是信里洋溢着正统信仰，一如你平日的习惯。我早已知道，不是第一次，而是很久以前就知道，你暂住在我们亲爱的\Person{英诺森}那里的原因，我对你的虔诚感到高兴。既然你如此行事，请写信告诉我那里的弟兄们如何，以及真理的敌人对我们怎样看待。你既告诉我\Place{凯撒利亚}的隐修士们的事，而我已从我们亲爱的\Person{迪亚尼}得知他们感到烦恼，并且反对我们亲爱的主教\Person{巴西略}，我感谢你告知我，并已向他们指出适宜的事，就是他们应当如子女服从父亲，不要反对他所赞同的事。因为如果他在真理方面受到怀疑，他们与他抗争就做得对。但如果他们确信，正如我们众人所确信的，他是教会的光荣，更是为真理而奋斗，并教导那些需要教导的人，那么与这样的人抗争就不对，反而应当怀着感恩接纳他的\OriginalTerm{良心}{conscience}。因为据亲爱的\Person{迪亚尼}所述，他们的烦恼似乎是毫无理由的。因为正如我所确信的，他对软弱的人，就成为软弱的，为赢得那软弱的人。\BibleRef{格前 9:22} 但愿我们亲爱的朋友们看到他的真理目标和他的特别目的，并光荣主，因为主把这样一位主教赐给了\Place{卡帕多细亚}，这是任何地区都渴望得到的。亲爱的，请你按我所写的向他们指出服从的责任。因为这样做，不仅会使他们对父亲有好感，而且会为各教会保全和平。我祈求你在主内安好，亲爱的儿子。\par

\SourceRecord{Translated by Archibald Robertson.；Nicene and Post-Nicene Fathers, Second Series；Vol. 4.；Edited by Philip Schaff and Henry Wace.；Buffalo, NY: Christian Literature Publishing Co.,；1892.；<http://www.newadvent.org/fathers/2806063.htm>.}

% structure: p0001 source=h1 target=section reason=page-title

\section{第64封书信}

\Emph{致狄奥多罗斯（残篇）。}\par

致我主、我子、并最亲爱的同为主仆的\Person{狄奥多罗斯}[\Emph{提洛}的主教]，\Person{亚大纳削}在主内致候。\par

我感谢我主，祂在各地建立祂的教义，而主要是借着祂自己的众子，正如事实所显示的您就是这样的人。因为在您写信之前，我们就已经知道，因着您的坚忍，在\Place{提洛}成就了多么大的恩宠。我们与您一同喜乐，因为通过您，\Place{提洛}也学会了正确的虔诚之言。我确实找到了机会给您写信，您是我长久思念和亲爱的；但我惊讶于您尚未回复我的信。因此不要迟延，请立即写信，因为您知道您如同儿子对父亲一样，给我带来慰藉，并且如同真理的宣报者一样，使我极其喜乐。不要与异端者辩论，而是以沉默克胜他们的好辩，以礼貌克胜他们的恶意。因为这样，您的言语当是\BibleQuote{以恩宠，用盐调和}\BibleRef{哥 4:6}，而他们[将被所有人的良心所判断]。\par

\SourceRecord{Translated by Archibald Robertson.；Nicene and Post-Nicene Fathers, Second Series；Vol. 4.；Edited by Philip Schaff and Henry Wace.；Buffalo, NY: Christian Literature Publishing Co.,；1892.；<http://www.newadvent.org/fathers/2806064.htm>.}
